\documentclass[11pt,a4paper,oneside,openany]{book}

\usepackage[T1]{fontenc}
\usepackage[utf8]{inputenc}
\usepackage{libertinus}
\usepackage{microtype}
\usepackage[a4paper,margin=27mm,headheight=15pt]{geometry}
\usepackage{amsmath,amssymb,amsthm,mathtools}
\usepackage{bm}
\usepackage{mathrsfs}
\usepackage{booktabs,longtable,array,tabularx}
\usepackage{enumitem}
\usepackage{xcolor}
\usepackage{tikz}
\usetikzlibrary{arrows.meta,positioning,calc}
\usepackage[most]{tcolorbox}
\usepackage{fancyhdr}
\usepackage{xurl}
\usepackage{hyperref}
\usepackage[nameinlink,noabbrev,capitalise]{cleveref}
\usepackage{makeidx}
\makeindex

\definecolor{deepolive}{RGB}{67,79,45}
\definecolor{softolive}{RGB}{236,239,226}
\definecolor{darkteal}{RGB}{31,83,84}
\definecolor{warmgray}{RGB}{92,87,78}
\definecolor{softcoral}{RGB}{246,225,219}
\definecolor{linkblue}{RGB}{31,76,122}

\hypersetup{
  colorlinks=true,
  linkcolor=deepolive,
  citecolor=darkteal,
  urlcolor=linkblue,
  pdftitle={Combinatorial Coefficient Calculus: A Unified Proof-Oriented Monograph},
  pdfauthor={A unified mathematical synthesis},
  pdfsubject={Stirling, Bell, Eulerian, Bernoulli, Sheffer, Faa di Bruno, and Lagrange-Good coefficient calculus}
}

\setlist{itemsep=2pt,topsep=4pt,parsep=1pt}
\allowdisplaybreaks[3]
\raggedbottom
\setcounter{tocdepth}{2}
\setcounter{secnumdepth}{3}

\pagestyle{fancy}
\fancyhf{}
\fancyhead[L]{\small\nouppercase{\leftmark}}
\fancyhead[R]{\small\thepage}
\renewcommand{\headrulewidth}{0.4pt}
\fancypagestyle{plain}{%
  \fancyhf{}%
  \fancyfoot[C]{\thepage}%
  \renewcommand{\headrulewidth}{0pt}%
}

\newtheorem{theorem}{Theorem}[chapter]
\newtheorem{lemma}[theorem]{Lemma}
\newtheorem{proposition}[theorem]{Proposition}
\newtheorem{corollary}[theorem]{Corollary}
\newtheorem{identity}[theorem]{Identity}
\theoremstyle{definition}
\newtheorem{definition}[theorem]{Definition}
\newtheorem{example}[theorem]{Example}
\newtheorem{algorithm}[theorem]{Algorithm}
\theoremstyle{remark}
\newtheorem{remark}[theorem]{Remark}
\newtheorem{historical}[theorem]{Historical note}

\newtcolorbox{masterbox}[1][]{
  enhanced,breakable,colback=softolive,colframe=deepolive,
  boxrule=0.7pt,arc=1mm,left=2mm,right=2mm,top=1.5mm,bottom=1.5mm,
  title={Master principle},fonttitle=\bfseries,#1
}
\newtcolorbox{auditbox}[1][]{
  enhanced,breakable,colback=softcoral,colframe=warmgray,
  boxrule=0.6pt,arc=1mm,left=2mm,right=2mm,top=1.5mm,bottom=1.5mm,
  title={Editorial normalization},fonttitle=\bfseries,#1
}

\newcommand{\N}{\mathbb N}
\newcommand{\Z}{\mathbb Z}
\newcommand{\Q}{\mathbb Q}
\newcommand{\R}{\mathbb R}
\newcommand{\C}{\mathbb C}
\newcommand{\F}{\mathbb F}
\newcommand{\Bell}[1]{\mathsf B_{#1}}
\newcommand{\BellMinus}{\mathsf U}
\newcommand{\BellP}[2]{\mathcal B_{#1,#2}}
\newcommand{\BellC}[1]{\mathcal B_{#1}}
\newcommand{\OBellP}[2]{\widehat{\mathcal B}_{#1,#2}}
\newcommand{\stirone}[2]{\genfrac{[}{]}{0pt}{}{#1}{#2}}
\newcommand{\stirtwo}[2]{\genfrac{\{} {\}}{0pt}{}{#1}{#2}}
\newcommand{\eulerian}[2]{\genfrac{\langle}{\rangle}{0pt}{}{#1}{#2}}
\newcommand{\eulertwo}[2]{\left\langle\!\left\langle\genfrac{}{}{0pt}{}{#1}{#2}\right\rangle\!\right\rangle}
\newcommand{\falling}[2]{#1^{\underline{#2}}}
\newcommand{\rising}[2]{#1^{\overline{#2}}}
\newcommand{\coeff}[2]{[#1^{#2}]}
\newcommand{\res}{\operatorname*{Res}}
\newcommand{\Li}{\operatorname{Li}}
\newcommand{\Var}{\operatorname{Var}}
\newcommand{\E}{\mathbb E}
\newcommand{\1}{\mathbf 1}
\newcommand{\wt}{\operatorname{wt}}
\newcommand{\cyc}{\operatorname{cyc}}
\newcommand{\setpart}{\operatorname{Part}}
\newcommand{\dd}{\,\mathrm d}
\newcommand{\bitand}{\mathbin{\&}}
\DeclareMathOperator{\tr}{tr}
\DeclareMathOperator{\lcm}{lcm}
\DeclareMathOperator{\sgn}{sgn}
\DeclareMathOperator{\id}{id}
\DeclareMathOperator{\diag}{diag}

\title{\Huge\bfseries Combinatorial Coefficient Calculus\\[4mm]
\Large Basis Changes, Labelled Structures, Special Polynomials,\\
Higher Composition, and Multivariate Inversion}
\author{A unified proof-oriented monograph}
\date{September 2026}

\begin{document}
\frontmatter
\hypersetup{pageanchor=false}
\maketitle
\hypersetup{pageanchor=true}

\chapter*{Abstract}
\addcontentsline{toc}{chapter}{Abstract}
This monograph is a unified, proof-oriented treatment of the coefficient arrays,
polynomial families, and inversion principles that recur throughout enumerative
combinatorics, formal power-series calculus, finite-difference theory, probability,
and asymptotic analysis.  It consolidates five overlapping draft manuscripts into
one dependency-driven article, removes repeated exposition, reconciles conflicting
normalizations, and incorporates additional material from the supplied Wikipedia
and MathWorld dossiers together with primary mathematical sources.

The central objects are signed and unsigned Stirling numbers of the first kind,
Stirling numbers of the second kind, Lah and generalized factorial coefficients,
Bell numbers and Touchard polynomials, exponential and ordinary Bell polynomials,
Eulerian numbers of ordinary, type-$B$, and second order, Bernoulli, Euler,
Genocchi, N\"orlund, and Cauchy--Bernoulli polynomials, Fa\`a di Bruno composition
coefficients, Sheffer sequences and exponential Riordan arrays, and the
one-variable and multivariate Lagrange--Good inversion formulas.  Their
connections with finite differences, M\"obius inversion, moments and cumulants,
normal ordering, symmetric functions, gamma and polygamma functions, Barnes'
$G$-function, associahedra, rooted trees, Lambert $W$, Kepler's equation, and
Laplace- and Euler--Maclaurin-type asymptotics are developed systematically.

The emphasis is on derivation rather than cataloguing.  Triangular recurrences are
proved from insertions or basis changes; generating functions are obtained from
labelled decompositions; inverse transforms are derived by M\"obius inversion or
matrix algebra; higher chain rules are proved from set partitions; implicit-series
coefficients are proved by formal residues; and analytic versions are accompanied
by explicit hypotheses, branch information, singularity diagnostics, or remainder
bounds.  Where a draft contained an incomplete proof, an index ambiguity, a
suppressed Jacobian, or a schematic asymptotic placeholder, the statement has been
repaired and replaced by a complete, checkable form.

\chapter*{How to read this work}
\addcontentsline{toc}{chapter}{How to read this work}
The chapter order follows mathematical dependence rather than the order of the
source files.  Four recurring operations organize nearly all formulas:
\begin{enumerate}
\item change to a polynomial basis in which the relevant difference or derivative
      operator is simple;
\item decompose a labelled structure into connected components and translate that
      decomposition into an exponential generating function;
\item compose or invert formal series by Bell polynomials and residues; and
\item pass from exact coefficient formulas to analytic estimates by Cauchy,
      saddle-point, Euler--Maclaurin, Darboux, or Watson-type arguments.
\end{enumerate}
Readers interested primarily in enumerative combinatorics may follow the Stirling,
Bell, Eulerian, and associahedral chapters.  Readers interested in symbolic
calculus may proceed from Bell polynomials to Fa\`a di Bruno, Sheffer--Riordan
calculus, and Lagrange inversion.  The Bernoulli--Euler and special-function part
forms a parallel route from finite differences to summation and asymptotics.  The
appendices record source provenance, normalization conversions, reproducible
algorithms, and independent consistency checks.

A formal power-series theorem is an algebraic statement: no convergence is
required, and every coefficient depends on finitely many input coefficients.  An
analytic theorem is stated separately and includes the hypotheses needed for
local inversion, termwise differentiation, contour deformation, Fourier
convergence, or an asymptotic remainder estimate.

\begin{auditbox}
The consolidation applies the following conventions globally.
\begin{enumerate}
\item Bell numbers are $\Bell n$, complete exponential Bell polynomials are
      $\BellC n$, partial exponential Bell polynomials are $\BellP nk$, and
      ordinary partial Bell polynomials are $\OBellP nk$.
\item Bernoulli numbers and polynomials are $\beta_n$ and $\beta_n(x)$ with
      $\beta_1=-1/2$.  Euler polynomials $E_n(x)$ are distinguished from
      Eulerian polynomials and arrays.
\item Signed first-kind Stirling numbers are $s(n,k)$; unsigned first-kind numbers
      are $\stirone nk$.  The sign conversion is never left implicit.
\item Type-$B$ descents are counted in positions $0,1,\ldots,n-1$ with
      $\pi_0=0$.  This removes the one-step shift present in one draft.
\item Ordinary and exponential coefficient conventions are linked explicitly by
      factorial conversions.  Analytic inversion includes its Jacobian, branch,
      and regularity conditions; multivariate inversion includes Good's
      determinant.
\end{enumerate}
\end{auditbox}

\tableofcontents

\chapter*{Notation}
\addcontentsline{toc}{chapter}{Notation}
The ground coefficient ring is usually a commutative $\Q$-algebra; analytic
statements specialize to $\R$ or $\C$.  For $n\in\N$,
\[
 \falling{x}{n}=x(x-1)\cdots(x-n+1),\qquad
 \rising{x}{n}=x(x+1)\cdots(x+n-1),
\]
with both values equal to $1$ at $n=0$.  Coefficient extraction and formal
residue notation are
\[
 \coeff z nF(z)=[z^n]F(z),\qquad
 \res_zF(z)\dd z=[z^{-1}]F(z).
\]

The principal triangular arrays are
\[
 s(n,k),\qquad \stirone nk=(-1)^{n-k}s(n,k),\qquad
 \stirtwo nk,\qquad \eulerian nk,\qquad \eulertwo nk.
\]
Bell numbers are $\Bell n$; the complementary Bell number is
$\BellMinus_n=T_n(-1)$.  Exponential partial and complete Bell polynomials are
$\BellP nk$ and $\BellC n$, and $\OBellP nk$ denotes the ordinary partial Bell
polynomial.  The Bernoulli convention is
\[
 \frac{t}{e^t-1}=\sum_{n\ge0}\beta_n\frac{t^n}{n!},
 \qquad \beta_1=-\frac12.
\]
Unless stated otherwise, an impossible index range contributes zero, empty
products equal $1$, and $0^0=1$ in finite coefficient formulas.  A bar such as
$\bar f$ or the notation $f^{\langle-1\rangle}$ denotes compositional inversion,
not reciprocal inversion.

\mainmatter

\part{Coefficient Calculus and Labelled Structures}

\chapter{Formal power series and coefficient extraction}
\label{ch:formal}

\section{Formal series as an algebraic environment}
Let $R$ be a commutative $\Q$-algebra.  The ring $R[[z]]$ consists of expressions
$F(z)=\sum_{n\ge0}f_nz^n$ with termwise addition and Cauchy multiplication.  A series
is invertible precisely when $f_0$ is invertible in $R$.  If $G(0)=0$, then the
composition $F(G(z))$ is well-defined because each coefficient receives only
finitely many contributions.

\begin{lemma}[Coefficient rules]\label{lem:coeff-rules}
For $F(z)=\sum f_nz^n$ and $G(z)=\sum g_nz^n$,
\begin{align}
 \coeff z n F(z)G(z)&=\sum_{j=0}^n f_jg_{n-j},\label{eq:cauchy}\\
 \coeff z n F'(z)&=(n+1)f_{n+1},\label{eq:der-coeff}\\
 \coeff z n \frac{F(z)}{1-az}&=\sum_{j=0}^n a^{n-j}f_j.\label{eq:geom-conv}
\end{align}
If $F$ is analytic on $|z|\le r$, then also
\begin{equation}\label{eq:cauchy-coeff}
 f_n=\frac{1}{2\pi i}\int_{|z|=r}\frac{F(z)}{z^{n+1}}\dd z,
 \qquad |f_n|\le r^{-n}\max_{|z|=r}|F(z)|.
\end{equation}
\end{lemma}
\begin{proof}
The first identity is the definition of Cauchy multiplication.  The second follows
by differentiating term by term, and the third by multiplying with
$\sum_{m\ge0}a^mz^m$.  The analytic formula is Cauchy's coefficient formula, and the
bound follows by estimating the contour integral by its length $2\pi r$.
\end{proof}

\section{Formal residues}
For a Laurent series $A(z)$, set
$\res_z A(z)\dd z=\coeff z{-1}A(z)$.  Then
\begin{equation}\label{eq:res-coeff}
 \coeff z n F(z)=\res_z\frac{F(z)}{z^{n+1}}\dd z,
 \qquad \res_z A'(z)\dd z=0.
\end{equation}
The second identity is simply the fact that the derivative of $z^m$ has no
$z^{-1}$ term.

\begin{theorem}[Formal change of variables for residues]\label{thm:res-subst}
Let $u=u(z)\in zR[[z]]$ have invertible linear coefficient.  For every Laurent series
$A$ for which the composition is defined,
\begin{equation}\label{eq:formal-res-subst}
 \res_z A(u(z))u'(z)\dd z=\res_u A(u)\dd u.
\end{equation}
\end{theorem}
\begin{proof}
By linearity it suffices to take $A(u)=u^m$.  If $m\ne-1$, then
$u(z)^mu'(z)=\frac{1}{m+1}(u(z)^{m+1})'$ has residue zero.  For $m=-1$,
$u'(z)/u(z)=z^{-1}+H(z)$ with $H\in R[[z]]$, because
$u(z)=az(1+O(z))$.  Thus both sides have residue $1$.
\end{proof}

\begin{masterbox}
The combination of \cref{eq:res-coeff,eq:formal-res-subst} turns implicit
substitution into coefficient extraction.  This is the entire algebraic engine of
Lagrange inversion; no analytic contour is needed until one asks about convergence.
\end{masterbox}

\section{Ordinary and exponential generating functions}
For a sequence $(a_n)$, its ordinary and exponential generating functions are
\[
 A(z)=\sum_{n\ge0}a_nz^n,
 \qquad \widehat A(z)=\sum_{n\ge0}a_n\frac{z^n}{n!}.
\]
Ordinary multiplication gives unweighted convolution.  Exponential multiplication
gives binomial convolution:
\begin{equation}\label{eq:egf-product}
 \widehat A(z)\widehat B(z)
 =\sum_{n\ge0}\left(\sum_{j=0}^n\binom nj a_jb_{n-j}\right)\frac{z^n}{n!}.
\end{equation}
This binomial factor is the number of ways to choose which labels belong to the
first component.

\section{Finite differences and factorial bases}
Let $E$ be the shift operator $(Ef)(x)=f(x+1)$ and
$\Delta=E-1$.  Then
\begin{equation}\label{eq:finite-diff-falling}
 \Delta\falling{x}{k}=k\falling{x}{k-1},
 \qquad
 \sum_{i=0}^{n-1}\falling{i}{k}=\frac{\falling{n}{k+1}}{k+1}.
\end{equation}
\begin{proof}
The first follows from
$\falling{x+1}{k}-\falling{x}{k}=k\falling{x}{k-1}$ after factoring the common
product.  Summing this telescoping identity gives the second formula.
\end{proof}

\begin{theorem}[Newton expansion]\label{thm:newton-expansion}
Every polynomial $p$ of degree at most $d$ has the unique expansion
\begin{equation}\label{eq:newton-expansion}
 p(x)=\sum_{k=0}^d\frac{\Delta^kp(0)}{k!}\falling{x}{k}
     =\sum_{k=0}^d\Delta^kp(0)\binom{x}{k}.
\end{equation}
\end{theorem}
\begin{proof}
The factorial polynomials form a basis because their degrees are distinct and their
leading coefficients are one.  If $p=\sum a_k\falling{x}{k}$, applying $\Delta^j$
and setting $x=0$ kills terms with $k>j$ by a remaining factor $x$, kills terms
with $k<j$ by excessive differencing, and gives $j!a_j$ for $k=j$.
\end{proof}

\section{Labelled products and the exponential formula}
A labelled combinatorial class is one whose atoms carry distinct labels.  Suppose a
connected component of size $j$ has total weight $x_j$.  A set of $k$ such components
with total size $n$ has weight enumerator $\BellP nk(x_1,x_2,\ldots)$; a set of any
number of components has enumerator $\BellC n(x_1,\ldots,x_n)$.

\begin{theorem}[The exponential formula]\label{thm:exponential-formula}
Let
\[
 C(z)=\sum_{j\ge1}x_j\frac{z^j}{j!}.
\]
Then
\begin{align}
 \frac{C(z)^k}{k!}
  &=\sum_{n\ge k}\BellP nk(x_1,\ldots,x_{n-k+1})\frac{z^n}{n!},
  \label{eq:partial-bell-egf}\\
 e^{uC(z)}
  &=\sum_{n\ge k\ge0}\BellP nk(x_1,\ldots)u^k\frac{z^n}{n!},
  \label{eq:bivariate-bell-egf}\\
 e^{C(z)}
  &=\sum_{n\ge0}\BellC n(x_1,\ldots,x_n)\frac{z^n}{n!}.
  \label{eq:complete-bell-egf}
\end{align}
\end{theorem}
\begin{proof}
Expand $C(z)^k/k!$ multinomially.  If $j_i$ components have size $i$, then
$\sum_i j_i=k$ and $\sum_iij_i=n$, and the coefficient of $z^n$ is
\[
 \prod_i\frac1{j_i!}\left(\frac{x_i}{i!}\right)^{j_i}.
\]
Multiplication by $n!$ labels the atoms.  This proves the first formula.  Summing it
with weight $u^k$ proves the second, and setting $u=1$ proves the third.
\end{proof}

\begin{corollary}[Partition-type count]\label{cor:partition-type}
The number of partitions of an $n$-element labelled set with exactly $j_i$ blocks of
size $i$ is
\begin{equation}\label{eq:partition-type-count}
 \frac{n!}{\prod_{i\ge1}(i!)^{j_i}j_i!},
 \qquad \sum_iij_i=n.
\end{equation}
\end{corollary}
\begin{proof}
Order the $n$ labels, cut the ordering into blocks of prescribed sizes, then divide
by $i!$ for the internal order of each size-$i$ block and by $j_i!$ for the order of
blocks of equal size.
\end{proof}


\section{Binomial convolution and Boolean-lattice inversion}
\index{binomial inversion}
The forward-difference formula is the one-dimensional shadow of M\"obius
inversion on subsets.  It is convenient to record the full transform before
specializing it repeatedly in later chapters.

\begin{theorem}[Parameterized binomial inversion]\label{thm:parameterized-binomial-inversion}
Let $c$ be an element of a commutative ring and let $(a_n)_{n\ge0}$ and
$(b_n)_{n\ge0}$ be sequences.  The relations
\begin{equation}\label{eq:parameterized-binomial-forward}
 b_n=\sum_{k=0}^{n}\binom nk c^{n-k}a_k
\end{equation}
and
\begin{equation}\label{eq:parameterized-binomial-inverse}
 a_n=\sum_{k=0}^{n}\binom nk(-c)^{n-k}b_k
\end{equation}
are equivalent.  If
$A(z)=\sum_{n\ge0}a_nz^n/n!$ and
$B(z)=\sum_{n\ge0}b_nz^n/n!$, then both are equivalent to
\begin{equation}\label{eq:parameterized-binomial-egf}
 B(z)=e^{cz}A(z).
\end{equation}
\end{theorem}
\begin{proof}
Substitute \eqref{eq:parameterized-binomial-forward} into the right side of
\eqref{eq:parameterized-binomial-inverse}.  The coefficient of $a_j$ becomes
\begin{align*}
 \sum_{k=j}^{n}\binom nk(-c)^{n-k}\binom kjc^{k-j}
 &=\binom njc^{n-j}
   \sum_{r=0}^{n-j}\binom{n-j}{r}(-1)^{n-j-r}\\
 &=\binom njc^{n-j}(1-1)^{n-j},
\end{align*}
which is $1$ for $j=n$ and $0$ otherwise.  This proves one implication;
replacing $c$ by $-c$ proves the converse.  For the generating-function form,
write
\begin{align*}
 B(z)
 &=\sum_{n\ge0}\sum_{k=0}^{n}\binom nkc^{n-k}a_k\frac{z^n}{n!}\\
 &=\left(\sum_{r\ge0}\frac{(cz)^r}{r!}\right)
   \left(\sum_{k\ge0}a_k\frac{z^k}{k!}\right)=e^{cz}A(z).
\end{align*}
Multiplication by $e^{-cz}$ gives the inverse transform.
\end{proof}

\begin{theorem}[M\"obius inversion on the Boolean lattice]
\label{thm:boolean-mobius}
Let $2^X$ be the subsets of a finite set $X$, ordered by inclusion.  Its
M\"obius function is
\begin{equation}\label{eq:boolean-mobius}
 \mu(S,T)=(-1)^{|T\setminus S|}\qquad(S\subseteq T).
\end{equation}
Consequently, the relations
\begin{equation}\label{eq:boolean-zeta-transform}
 G(T)=\sum_{S\subseteq T}F(S)
\end{equation}
and
\begin{equation}\label{eq:boolean-mobius-transform}
 F(T)=\sum_{S\subseteq T}(-1)^{|T\setminus S|}G(S)
\end{equation}
are equivalent.
\end{theorem}
\begin{proof}
For $S\subseteq T$ the defining convolution identity for a M\"obius function
is
\[
 \sum_{S\subseteq U\subseteq T}\mu(S,U)=\mathbf 1_{S=T}.
\]
The proposed value gives
\[
 \sum_{S\subseteq U\subseteq T}(-1)^{|U\setminus S|}
 =\sum_{r=0}^{|T\setminus S|}\binom{|T\setminus S|}{r}(-1)^r
 =(1-1)^{|T\setminus S|},
\]
so it is the M\"obius function.  Substitution of
\eqref{eq:boolean-zeta-transform} into
\eqref{eq:boolean-mobius-transform} and reversal of the finite sums now gives
exactly the same convolution identity.
\end{proof}

If $F(S)$ depends only on $|S|$, grouping subsets by cardinality turns Boolean
M\"obius inversion into ordinary binomial inversion.  Thus inclusion--exclusion,
finite differences, and the triangular transform in
\cref{thm:parameterized-binomial-inversion} are the same algebraic operation in
three coordinate systems.

\part{Stirling Numbers, Lah Numbers, and Basis Changes}

\chapter{The two Stirling triangles}
\label{ch:stirling-def}

\section{First-kind numbers: cycles and factorial coefficients}
\begin{definition}
The signed Stirling numbers of the first kind $s(n,k)$ and the unsigned numbers
$\stirone nk$ are defined by
\begin{align}
 \falling{x}{n}&=\sum_{k=0}^n s(n,k)x^k,\label{eq:signed-first-def}\\
 \rising{x}{n}&=\sum_{k=0}^n\stirone nkx^k,\label{eq:unsigned-first-def}\\
 s(n,k)&=(-1)^{n-k}\stirone nk.\label{eq:first-sign}
\end{align}
\end{definition}
The sign relation follows from
$\rising{x}{n}=(-1)^n\falling{-x}{n}$.

\begin{theorem}[Permutation interpretation]\label{thm:first-cycle}
$\stirone nk$ is the number of permutations of $[n]$ having exactly $k$ cycles.
Consequently,
\begin{equation}\label{eq:first-recurrence}
 \stirone{n+1}{k}=n\stirone nk+\stirone n{k-1},
\end{equation}
with $\stirone00=1$ and zero boundary values outside $0\le k\le n$.
\end{theorem}
\begin{proof}
Given a permutation of $[n+1]$ with $k$ cycles, either $n+1$ is a singleton cycle,
leaving a permutation of $[n]$ with $k-1$ cycles, or it is inserted immediately
after one of the $n$ existing elements in the cyclic notation of a permutation of
$[n]$ with $k$ cycles.  These alternatives are disjoint and reversible.
The same recurrence follows algebraically by multiplying
$\rising{x}{n}$ by $x+n$ and comparing coefficients.
\end{proof}

For $n=3$, the cycle counts are $\stirone33=1$, $\stirone32=3$, and
$\stirone31=2$, corresponding to the identity, the three transpositions, and the two
$3$-cycles.  Summing over the possible number of cycles gives
\begin{equation}\label{eq:first-row-sum}
 \sum_{k=0}^n\stirone nk=n!,
\end{equation}
which also follows from \cref{eq:unsigned-first-def} at $x=1$.

\section{Second-kind numbers: set partitions and surjections}
\begin{definition}
The Stirling number of the second kind $\stirtwo nk$ is the number of partitions of
$[n]$ into exactly $k$ nonempty blocks.
\end{definition}

\begin{theorem}[Second-kind recurrence]\label{thm:second-recurrence}
For $n,k\ge0$,
\begin{equation}\label{eq:second-recurrence}
 \stirtwo{n+1}{k}=k\stirtwo nk+\stirtwo n{k-1},
\end{equation}
with $\stirtwo00=1$ and zero boundary values outside $0\le k\le n$.
\end{theorem}
\begin{proof}
In a partition of $[n+1]$ into $k$ blocks, the new label either joins one of the $k$
blocks of a partition of $[n]$, or forms a singleton block beside a partition of
$[n]$ into $k-1$ blocks.
\end{proof}

\begin{theorem}[Surjection formula]\label{thm:second-explicit}
For all $n,k\ge0$,
\begin{equation}\label{eq:second-explicit}
 \stirtwo nk=\frac1{k!}\sum_{j=0}^k(-1)^{k-j}\binom kjj^n
 =\sum_{j=0}^k\frac{(-1)^{k-j}j^n}{(k-j)!j!}.
\end{equation}
\end{theorem}
\begin{proof}
There are $k!\stirtwo nk$ surjections from $[n]$ onto a fixed $k$-element codomain:
a set partition into fibres followed by an ordering of the fibres.  Inclusion--
exclusion over the set of omitted codomain values gives
$\sum_{j=0}^k(-1)^{k-j}\binom kjj^n$.  Division by $k!$ proves the formula.
\end{proof}

The simple diagonals follow at once:
\begin{equation}\label{eq:second-simple}
 \stirtwo nn=1,\qquad \stirtwo n1=1\ (n\ge1),\qquad
 \stirtwo n{n-1}=\binom n2,\qquad
 \stirtwo n2=2^{n-1}-1.
\end{equation}
The last identity chooses a nonempty proper subset as one block and divides by two
for exchanging the blocks.


\section{Alternative and reverse recurrences}
The ordinary recurrence builds each row from the preceding one.  Several less
familiar identities in the source archive instead recover an entry from entries to
its right or from earlier entries in the same column.  They were labelled there as
conjectures; in fact they are exact identities.

\begin{theorem}[Alternative second-kind recurrences]
\label{thm:second-reverse-recurrences}
For $1\le k<n$,
\begin{align}
 \stirtwo nk
 &=\frac{k^n}{k!}-\sum_{r=1}^{k-1}\frac{\stirtwo nr}{(k-r)!},
 \label{eq:second-triangular-explicit}\\
 (n-k)\stirtwo nk
 &=\sum_{j=2}^{n-k+1}(j-2)!\binom{-k}{j}
   \stirtwo n{k+j-1},
 \label{eq:second-reverse-row}\\
 (n-k)\stirtwo nk
 &=\sum_{j=2}^{n-k+1}(-1)^j\binom nj
   \stirtwo{n-j+1}{k}.
 \label{eq:second-reverse-column}
\end{align}
The terminal value in either reverse recursion is $\stirtwo nn=1$.
\end{theorem}
\begin{proof}
Substitute $x=k$ in the basis identity
$k^n=\sum_r\stirtwo nr\falling{k}{r}$.  Since
$\falling{k}{r}=k!/(k-r)!$ for $r\le k$, division by $k!$ and isolation of the
$r=k$ term gives \eqref{eq:second-triangular-explicit}.

For the row reversal, use the bivariate exponential generating function
\[
 F(x,y)=\exp\!\bigl(y(e^x-1)\bigr)
 =\sum_{n,k\ge0}\stirtwo nk\frac{x^n}{n!}y^k,
 \qquad u=e^x-1.
\]
The generating function of the left side of
\eqref{eq:second-reverse-row} is
\[
 xF_x-yF_y=yF\bigl(xe^x-u\bigr).
\]
On the other hand, because
$\binom{-k}{j}=(-1)^j\binom{k+j-1}{j}$, the generating function of its right side is
\begin{align*}
 y\sum_{j\ge2}\frac{(-1)^j}{j(j-1)}\,\partial_y^jF
 &=yF\sum_{j\ge2}\frac{(-1)^ju^j}{j(j-1)}\\
 &=yF\bigl((1+u)\log(1+u)-u\bigr)\\
 &=yF\bigl(xe^x-u\bigr).
\end{align*}
Here the middle series is obtained by differentiating it twice, summing the
geometric series, and restoring the two zero initial coefficients.  Equality of
coefficients proves \eqref{eq:second-reverse-row}.

For the column reversal fix $k$ and put
\[
 F_k(x)=\sum_{n\ge k}\stirtwo nk\frac{x^n}{n!}
       =\frac{(e^x-1)^k}{k!}.
\]
The two sides of \eqref{eq:second-reverse-column}, after multiplication by
$x^n/n!$ and summation over $n$, become respectively
\[
 xF_k'(x)-kF_k(x)
 \quad\text{and}\quad
 (e^{-x}-1+x)F_k'(x).
\]
They are equal because direct differentiation gives
$(1-e^{-x})F_k'(x)=kF_k(x)$.
\end{proof}

The factor interchange mentioned in the source produces a parallel pair for the
unsigned first-kind triangle.

\begin{theorem}[Reverse recurrences of the first kind]
\label{thm:first-reverse-recurrences}
For $1\le k<n$,
\begin{align}
 (n-k)\stirone nk
 &=\sum_{j=2}^{n-k+1}(-1)^j\binom{-k}{j}
   \stirone n{k+j-1},
 \label{eq:first-reverse-row}\\
 (n-k)\stirone nk
 &=\sum_{j=2}^{n-k+1}(j-2)!\binom nj
   \stirone{n-j+1}{k}.
 \label{eq:first-reverse-column}
\end{align}
\end{theorem}
\begin{proof}
Let $C_n(y)=\rising{y}{n}=\sum_r\stirone nr y^r$.  Since
\[
 C_n(1+y)=\frac{\rising{y}{n+1}}{y},
\]
the coefficient of $y^{k-1}$ on the two sides is, respectively,
\[
 \sum_{r\ge k-1}\binom r{k-1}\stirone nr
 \quad\text{and}\quad
 \stirone{n+1}{k}=n\stirone nk+\stirone n{k-1}.
\]
Remove the terms $r=k-1$ and $r=k$ and put $r=k+j-1$.  Because
$(-1)^j\binom{-k}{j}=\binom{k+j-1}{j}$, the remainder is precisely
\eqref{eq:first-reverse-row}.

For fixed $k$, set
\[
 C_k(x)=\sum_{n\ge k}\stirone nk\frac{x^n}{n!}
       =\frac{[-\log(1-x)]^k}{k!}.
\]
The exponential generating function of the right side of
\eqref{eq:first-reverse-column} is
\[
 C_k'(x)\sum_{j\ge2}\frac{x^j}{j(j-1)}
 =\bigl(x+(1-x)\log(1-x)\bigr)C_k'(x).
\]
The left side has generating function $xC_k'-kC_k$.  Their equality reduces to
\[
 -(1-x)\log(1-x)\,C_k'(x)=kC_k(x),
\]
which follows immediately from the displayed formula for $C_k$.
\end{proof}
\section{Tables and common boundary conventions}
The first rows are
\[
\begin{array}{c|rrrrrrr}
 n\backslash k&0&1&2&3&4&5&6\\\hline
0&1\\
1&0&1\\
2&0&1&1\\
3&0&2&3&1\\
4&0&6&11&6&1\\
5&0&24&50&35&10&1\\
6&0&120&274&225&85&15&1
\end{array}
\qquad
\begin{array}{c|rrrrrrr}
 n\backslash k&0&1&2&3&4&5&6\\\hline
0&1\\
1&0&1\\
2&0&1&1\\
3&0&1&3&1\\
4&0&1&7&6&1\\
5&0&1&15&25&10&1\\
6&0&1&31&90&65&15&1
\end{array}
\]
for $\stirone nk$ and $\stirtwo nk$, respectively.

\chapter{Factorial bases and inverse matrices}
\label{ch:bases}

\section{The basic changes of basis}
\begin{theorem}[Stirling basis transformations]\label{thm:stirling-basis}
For every $n\ge0$,
\begin{align}
 \falling{x}{n}&=\sum_{k=0}^ns(n,k)x^k,\label{eq:fall-to-power}\\
 \rising{x}{n}&=\sum_{k=0}^n\stirone nkx^k,\label{eq:rise-to-power}\\
 x^n&=\sum_{k=0}^n\stirtwo nk\falling{x}{k},\label{eq:power-to-fall}\\
 x^n&=\sum_{k=0}^n(-1)^{n-k}\stirtwo nk\rising{x}{k}.\label{eq:power-to-rise}
\end{align}
Equivalently,
\begin{equation}\label{eq:binomial-basis}
 x^n=\sum_{k=0}^n\stirtwo nk k!\binom{x}{k},\qquad
 \binom{x}{n}=\frac1{n!}\sum_{k=0}^ns(n,k)x^k.
\end{equation}
\end{theorem}
\begin{proof}
The first two identities are definitions.  In the Newton expansion
\cref{eq:newton-expansion} for $p(x)=x^n$, one has
\[
 \Delta^kx^n\big|_{x=0}=\sum_{j=0}^k(-1)^{k-j}\binom kjj^n
 =k!\stirtwo nk
\]
by \cref{eq:second-explicit}.  This gives \cref{eq:power-to-fall}.  Replacing $x$
by $-x$ and multiplying by $(-1)^n$ gives \cref{eq:power-to-rise}.  Division of
$\falling{x}{n}$ by $n!$ gives the second binomial-basis formula.
\end{proof}


\begin{corollary}[Shifted factorial evaluations]
\label{cor:shifted-stirling-evaluations}
For every $n\ge0$,
\begin{align}
 \sum_{k=1}^{n+1}\stirtwo{n+1}{k}\falling{x-1}{k-1}&=x^n,
 \label{eq:shifted-power-to-fall}\\
 \sum_{k=0}^{n}\stirtwo nk\falling{n}{k}&=n^n.
 \label{eq:stirling-n-to-n}
\end{align}
\end{corollary}
\begin{proof}
Apply \eqref{eq:power-to-fall} with exponent $n+1$ and use
$\falling{x}{k}=x\falling{x-1}{k-1}$.  Both sides are divisible by $x$ as
polynomials, and cancellation gives \eqref{eq:shifted-power-to-fall}.  Formula
\eqref{eq:stirling-n-to-n} is \eqref{eq:power-to-fall} evaluated at $x=n$.
\end{proof}
\begin{corollary}[Inverse matrices]\label{cor:stirling-inverse}
The lower triangular matrices $s=(s(n,k))$ and $S=(\stirtwo nk)$ are mutual inverses:
\begin{align}
 \sum_{j=k}^ns(n,j)\stirtwo jk&=\delta_{nk},\label{eq:stirling-inv1}\\
 \sum_{j=k}^n\stirtwo nj s(j,k)&=\delta_{nk}.\label{eq:stirling-inv2}
\end{align}
Equivalently,
\begin{align}
 \sum_{j=k}^n(-1)^{n-j}\stirone nj\stirtwo jk&=\delta_{nk},\\
 \sum_{j=k}^n(-1)^{j-k}\stirtwo nj\stirone jk&=\delta_{nk}.
\end{align}
\end{corollary}
\begin{proof}
Expand $\falling{x}{n}$ into powers and each power back into falling factorials.
Uniqueness of coordinates in the factorial basis gives the first matrix product.
The reverse substitution gives the second.
\end{proof}

\section{Power sums by basis adaptation}
The factorial basis makes discrete antidifferentiation immediate.  Combining
\cref{eq:power-to-fall,eq:finite-diff-falling} yields
\begin{equation}\label{eq:power-sum-stirling}
 \sum_{i=0}^{N}i^m
 =\sum_{k=0}^m\stirtwo mk\frac{\falling{N+1}{k+1}}{k+1}.
\end{equation}
For $m=4$ this is
\[
 \frac12\falling{N+1}{2}+\frac73\falling{N+1}{3}
 +\frac64\falling{N+1}{4}+\frac15\falling{N+1}{5},
\]
which expands to the familiar quartic power-sum polynomial.


\section{Factorial moments: Poisson variables and permutation fixed points}
The change from ordinary powers to falling factorials is especially effective for
integer-valued random variables.

\begin{theorem}[Poisson moments]
\label{thm:poisson-stirling-moments}
If $X$ has the Poisson distribution of mean $\lambda$, then
\begin{equation}
\label{eq:poisson-stirling-moments}
 \E X^n=\sum_{k=0}^{n}\stirtwo nk\lambda^k.
\end{equation}
In particular, $\E X^n=\Bell n$ when $\lambda=1$.
\end{theorem}
\begin{proof}
The $k$th falling-factorial moment is
\[
 \E\falling{X}{k}
 =e^{-\lambda}\sum_{r\ge k}\frac{\falling r k\lambda^r}{r!}
 =\lambda^ke^{-\lambda}\sum_{q\ge0}\frac{\lambda^q}{q!}
 =\lambda^k.
\]
Take expectations in the basis identity
$X^n=\sum_k\stirtwo nk\falling Xk$.
\end{proof}

\begin{theorem}[Moments of the number of fixed points]
\label{thm:fixed-point-stirling-moments}
Let $Y_m$ be the number of fixed points of a uniformly random permutation of
$[m]$.  Then
\begin{equation}
\label{eq:fixed-point-stirling-moments}
 \E Y_m^n=\sum_{k=0}^{m}\stirtwo nk
 =\sum_{k=0}^{\min(m,n)}\stirtwo nk.
\end{equation}
Thus, when $m\ge n$, the moment equals the Bell number $\Bell n$.
\end{theorem}
\begin{proof}
The random variable $\falling{Y_m}{k}$ counts ordered $k$-tuples of distinct fixed
points.  There are $\falling m k$ candidate tuples, and a specified tuple is fixed
with probability $(m-k)!/m!$.  Hence
\[
 \E\falling{Y_m}{k}=1\quad(0\le k\le m),
 \qquad
 \E\falling{Y_m}{k}=0\quad(k>m).
\]
Taking expectations in \eqref{eq:power-to-fall} proves the formula.  The upper
limit $m$ is essential: it records the finite size of the permutation, even when
$n>m$.
\end{proof}
\section{Lah numbers}
\begin{definition}
For $1\le k\le n$, the unsigned Lah number is
\begin{equation}\label{eq:lah-def}
 L(n,k)=\binom{n-1}{k-1}\frac{n!}{k!},
\end{equation}
with $L(0,0)=1$ and zero otherwise outside the triangle.
\end{definition}
Combinatorially, $L(n,k)$ counts partitions of $[n]$ into $k$ nonempty linearly
ordered lists, with the collection of lists unordered.

\begin{theorem}[Conversion between rising and falling factorials]\label{thm:lah-conv}
\begin{align}
 \rising{x}{n}&=\sum_{k=0}^nL(n,k)\falling{x}{k},\label{eq:lah-rise-fall}\\
 \falling{x}{n}&=\sum_{k=0}^n(-1)^{n-k}L(n,k)\rising{x}{k},\label{eq:lah-fall-rise}\\
 \sum_{j=k}^n(-1)^{j-k}L(n,j)L(j,k)&=\delta_{nk}.\label{eq:lah-inv}
\end{align}
Moreover,
\begin{equation}\label{eq:lah-stirling-product}
 L(n,k)=\sum_{j=k}^n\stirone nj\stirtwo jk.
\end{equation}
\end{theorem}
\begin{proof}
A direct generating-function calculation gives
\[
 \sum_{n\ge k}L(n,k)\frac{z^n}{n!}
 =\frac1{k!}\left(\frac{z}{1-z}\right)^k.
\]
Therefore
\[
 \sum_{n\ge0}\left(\sum_kL(n,k)\falling{x}{k}\right)\frac{z^n}{n!}
 =\sum_{k\ge0}\binom{x}{k}\left(\frac{z}{1-z}\right)^k
 =(1-z)^{-x}
 =\sum_{n\ge0}\rising{x}{n}\frac{z^n}{n!},
\]
proving \cref{eq:lah-rise-fall}.  Replace $x$ by $-x$ to obtain the inverse formula,
and compose the two transformations to obtain \cref{eq:lah-inv}.  Finally, expand a
rising factorial into powers by first-kind numbers and powers into falling factorials
by second-kind numbers; uniqueness gives \cref{eq:lah-stirling-product}.
\end{proof}

The inequalities
\begin{equation}\label{eq:stirling-lah-ineq}
 \stirtwo nk\le\stirone nk\le L(n,k)
\end{equation}
follow from explicit injections: a set partition may be canonically turned into a
permutation with one cycle per block by writing each block with its smallest member
first and arranging cycles by minima; forgetting the cyclic identifications embeds
permutations with $k$ cycles into unordered collections of $k$ lists.


\chapter{Generating functions, transforms, and operators}
\label{ch:stirling-gf}

\section{Exponential generating functions}
\begin{theorem}[Column generating functions]\label{thm:stirling-egfs}
For every fixed $k\ge0$,
\begin{align}
 \sum_{n\ge k}s(n,k)\frac{z^n}{n!}
   &=\frac{\log(1+z)^k}{k!},\label{eq:signed-first-egf}\\
 \sum_{n\ge k}\stirone nk\frac{z^n}{n!}
   &=\frac{[-\log(1-z)]^k}{k!},\label{eq:unsigned-first-egf}\\
 \sum_{n\ge k}\stirtwo nk\frac{z^n}{n!}
   &=\frac{(e^z-1)^k}{k!}.\label{eq:second-egf}
\end{align}
The bivariate second-kind generating function is
\begin{equation}\label{eq:second-double-egf}
 \sum_{k\ge0}\sum_{n\ge k}\stirtwo nk y^k\frac{z^n}{n!}
 =\exp\!\bigl(y(e^z-1)\bigr).
\end{equation}
\end{theorem}
\begin{proof}
Expand
\[
 (1+z)^u=e^{u\log(1+z)}
 =\sum_{n\ge0}\frac{z^n}{n!}\sum_{k=0}^ns(n,k)u^k
 =\sum_{k\ge0}\frac{u^k}{k!}\log(1+z)^k
\]
and compare coefficients of $u^k$.  Replacing $z$ by $-z$ and using the sign
relation gives \cref{eq:unsigned-first-egf}.  Formula \cref{eq:second-egf} is the
exponential formula for a set of $k$ nonempty labelled blocks; algebraically it also
follows by inserting \cref{eq:second-explicit} and summing exponentials.  Summing over
$k$ with weight $y^k$ proves \cref{eq:second-double-egf}.
\end{proof}

Differentiating \cref{eq:signed-first-egf} gives the useful identity
\begin{equation}\label{eq:log-over-one-plus}
 \frac{\log(1+z)^m}{1+z}
 =m!\sum_{r\ge0}s(r+1,m+1)\frac{z^r}{r!},\qquad |z|<1,
\end{equation}
which is valid formally as well.

\section{A rational ordinary generating function}
\begin{theorem}[Fixed-$k$ ordinary generating function]\label{thm:second-ogf}
For $k\ge0$,
\begin{equation}\label{eq:second-ogf}
 \sum_{n\ge k}\stirtwo nk x^{n-k}
 =\prod_{j=1}^k\frac1{1-jx}
 =\frac1{x^{k+1}\falling{1/x}{k+1}}.
\end{equation}
Consequently,
\begin{equation}\label{eq:second-complete-symmetric}
 \stirtwo nk=h_{n-k}(1,2,\ldots,k)
 =\sum_{c_1+\cdots+c_k=n-k}1^{c_1}2^{c_2}\cdots k^{c_k},
\end{equation}
where $h_r$ is a complete homogeneous symmetric polynomial.
\end{theorem}
\begin{proof}
Let $F_k(x)=\sum_{r\ge0}\stirtwo{k+r}{k}x^r$.  The recurrence
\cref{eq:second-recurrence}, with $n=k+r-1$, gives
$F_k(x)=(1-kx)^{-1}F_{k-1}(x)$ and $F_0=1$.  Iteration proves the product.
Expanding each geometric factor independently proves the complete-homogeneous sum.
The factorial form follows by direct factorization.
\end{proof}

A useful rearrangement of the count of all maps $[n]\to[k]$ is
\begin{equation}\label{eq:second-alternate-rec}
 \stirtwo nk=\frac{k^n}{k!}-\sum_{r=1}^{k-1}\frac{\stirtwo nr}{(k-r)!}.
\end{equation}
Indeed, a map using exactly $r$ values can be chosen in
$\binom kr r!\stirtwo nr$ ways; divide the identity
$k^n=\sum_r\binom kr r!\stirtwo nr$ by $k!$.

\section{The Stirling transform}
\begin{theorem}[Stirling transform and its inverse]\label{thm:stirling-transform}
Let
\[
 g_n=\sum_{k=0}^n\stirtwo nk f_k.
\]
Then
\begin{equation}\label{eq:stirling-transform-inverse}
 f_n=\sum_{k=0}^ns(n,k)g_k
 =\sum_{k=0}^n(-1)^{n-k}\stirone nkg_k.
\end{equation}
For exponential generating functions,
\begin{equation}\label{eq:stirling-transform-egf}
 \widehat G(z)=\widehat F(e^z-1),\qquad
 \widehat F(z)=\widehat G(\log(1+z)).
\end{equation}
\end{theorem}
\begin{proof}
The inverse formula is exactly the matrix inversion in
\cref{cor:stirling-inverse}.  For the generating function,
\[
 \widehat G(z)=\sum_kf_k\sum_{n\ge k}\stirtwo nk\frac{z^n}{n!}
 =\sum_kf_k\frac{(e^z-1)^k}{k!}=\widehat F(e^z-1).
\]
Substitution by the compositional inverse $\log(1+z)$ gives the reverse relation.
\end{proof}

\section{Differential-operator normal ordering}
Let $D=d/dx$.  The Euler operator $xD$ is diagonal on monomials:
$(xD)x^m=mx^m$.  The operator $x^kD^k$ has eigenvalue $\falling m k$.
Thus the scalar basis changes lift verbatim to operators.

\begin{theorem}[Stirling normal-ordering identities]\label{thm:normal-order}
For $n\ge0$,
\begin{align}
 (xD)^n&=\sum_{k=0}^n\stirtwo nk x^kD^k,\label{eq:normal1}\\
 x^nD^n&=\sum_{k=0}^ns(n,k)(xD)^k
        =\falling{xD}{n}.\label{eq:normal2}
\end{align}
\end{theorem}
\begin{proof}
Both sides of \cref{eq:normal1} act on $x^m$ as multiplication by
$m^n=\sum_k\stirtwo nk\falling mk$.  Since monomials span the polynomial ring, the
operators are equal.  The inverse relation is proved identically using
$\falling mn=\sum_ks(n,k)m^k$.
\end{proof}

Because $E=e^D$ and $\Delta=e^D-1$, substitution of the mutually inverse series
$D=\log(1+\Delta)$ and $\Delta=e^D-1$ gives, whenever the operator series converge
(or formally on polynomials),
\begin{align}
 \frac{D^k}{k!}
   &=\sum_{n\ge k}\frac{s(n,k)}{n!}\Delta^n,
   \label{eq:der-from-diff}\\
 \frac{\Delta^k}{k!}
   &=\sum_{n\ge k}\frac{\stirtwo nk}{n!}D^n.
   \label{eq:diff-from-der}
\end{align}
These are the derivative--finite-difference transformations listed in the archive.

\chapter{Symmetric functions, diagonals, and finite identities}
\label{ch:stirling-identities}

\section{Elementary and complete symmetric descriptions}
From \cref{eq:unsigned-first-def},
\begin{equation}\label{eq:first-elementary}
 \stirone n{n-r}=e_r(0,1,\ldots,n-1)
 =\sum_{0\le i_1<\cdots<i_r<n}i_1\cdots i_r.
\end{equation}
Together with \cref{eq:second-complete-symmetric}, this identifies the first-kind
triangle with elementary symmetric functions and the second-kind triangle with
complete homogeneous symmetric functions.

The first few near-diagonals are
\begin{align}
 \stirone n{n-1}&=\binom n2,\label{eq:first-diag1}\\
 \stirone n{n-2}&=\frac{3n-1}{4}\binom n3,\label{eq:first-diag2}\\
 \stirone n{n-3}&=\binom n2\binom n4.\label{eq:first-diag3}
\end{align}
\begin{proof}
The first is $e_1(0,\ldots,n-1)$.  Newton's identities
$2e_2=p_1^2-p_2$ and $6e_3=p_1^3-3p_1p_2+2p_3$, combined with the elementary sums
of first, second, and third powers, simplify to \cref{eq:first-diag2,eq:first-diag3}.
Equivalently, one may classify permutations by the nontrivial cycle types
$3$, $2+2$ for the second diagonal and $4$, $3+2$, $2+2+2$ for the third.
\end{proof}

\section{Harmonic-number formulae}
Factoring $(n-1)!$ from
$(x+1)(x+2)\cdots(x+n-1)$ gives
\begin{equation}\label{eq:first-harmonic-elementary}
 \frac{\stirone n{k+1}}{(n-1)!}
 =e_k\left(1,\frac12,\ldots,\frac1{n-1}\right)
 =\sum_{1\le i_1<\cdots<i_k<n}\frac1{i_1\cdots i_k}.
\end{equation}
Write $H_n^{(r)}=\sum_{j=1}^nj^{-r}$ and $H_n=H_n^{(1)}$.  Newton's identities now
give
\begin{align}
 \stirone n2&=(n-1)!H_{n-1},\label{eq:first-H1}\\
 \stirone n3&=\frac{(n-1)!}{2}\left(H_{n-1}^2-H_{n-1}^{(2)}\right),\label{eq:first-H2}\\
 \stirone n4&=\frac{(n-1)!}{6}\left(H_{n-1}^3-3H_{n-1}H_{n-1}^{(2)}
                   +2H_{n-1}^{(3)}\right).\label{eq:first-H3}
\end{align}
More generally,
\begin{equation}\label{eq:first-harmonic-bell}
 \frac{\stirone{n+1}{k+1}}{n!}
 =\frac1{k!}\BellC{k}\bigl(H_n,-1!H_n^{(2)},2!H_n^{(3)},\ldots,
              (-1)^{k-1}(k-1)!H_n^{(k)}\bigr).
\end{equation}
\begin{proof}
Take logarithms in
\[
 \prod_{j=1}^n\left(1+\frac{x}{j}\right)
 =\exp\left(\sum_{r\ge1}\frac{(-1)^{r-1}H_n^{(r)}}r x^r\right)
\]
and apply the complete Bell-polynomial generating function
\cref{eq:complete-bell-egf}.
\end{proof}

If $w(n,k)=k!\,\stirone n{k+1}/(n-1)!$, then complete Bell recurrence yields
\begin{equation}\label{eq:w-recurrence}
 w(n,m)=\delta_{m0}+\sum_{r=0}^{m-1}(1-m)_rH_{n-1}^{(r+1)}w(n,m-1-r),
\end{equation}
where $(1-m)_r$ is falling.  Conversely, taking a logarithm recovers the harmonic
power sums:
\begin{equation}\label{eq:harmonic-inversion}
 H_n^{(m)}=-m\coeff xm
 \log\left(1+\sum_{k\ge1}\frac{\stirone{n+1}{k+1}}{n!}(-x)^k\right).
\end{equation}
For example,
\begin{align}
 (n!)^2H_n^{(2)}
 &=\stirone{n+1}{2}^{\!2}
   -2\stirone{n+1}{1}\stirone{n+1}{3},\label{eq:H2-invert}\\
 (n!)^3H_n^{(3)}
 &=\stirone{n+1}{2}^{\!3}
   -3\stirone{n+1}{1}\stirone{n+1}{2}\stirone{n+1}{3}
   +3\stirone{n+1}{1}^{\!2}\stirone{n+1}{4}.
 \label{eq:H3-invert}
\end{align}

\section{Convolution and summation identities}
The following paired identities expose the parallel structure of the two triangles.

\begin{theorem}[First- and second-kind summations]\label{thm:paired-sums}
For valid nonnegative indices,
\begin{align}
 \stirone{n+1}{k+1}
   &=\sum_{j=k}^n\stirone nj\binom jk
    =\sum_{j=k}^n\frac{n!}{j!}\stirone jk,\label{eq:first-two-sums}\\
 \stirtwo{n+1}{k+1}
   &=\sum_{j=k}^n\binom nj\stirtwo jk
    =\sum_{j=k}^n(k+1)^{n-j}\stirtwo jk,\label{eq:second-two-sums}\\
 \stirone{n+k+1}{n}
   &=\sum_{j=0}^k(n+j)\stirone{n+j}{j},\label{eq:first-hockey}\\
 \stirtwo{n+k+1}{k}
   &=\sum_{j=0}^kj\stirtwo{n+j}{j}.\label{eq:second-hockey}
\end{align}
Moreover, for $\ell,m\ge0$,
\begin{align}
 \binom{\ell+m}{\ell}\stirone n{\ell+m}
 &=\sum_j\binom nj\stirone j\ell\stirone{n-j}m,
 \label{eq:first-convolution}\\
 \binom{\ell+m}{\ell}\stirtwo n{\ell+m}
 &=\sum_j\binom nj\stirtwo j\ell\stirtwo{n-j}m.
 \label{eq:second-convolution}
\end{align}
\end{theorem}
\begin{proof}
For the first equality in \cref{eq:first-two-sums}, mark $k$ cycles of a permutation
of $[n]$ and insert $n+1$ into the canonical construction; algebraically it is the
coefficient of $x^k$ in $\frac{d}{dx}\rising{x}{n}$.  The second equality follows by
iterating recurrence \cref{eq:first-recurrence} in the $n$ direction.  For
\cref{eq:second-two-sums}, distinguish the elements outside the block containing
$n+1$ for the first equality, and add the $n-k$ remaining elements one at a time to
one of $k+1$ existing blocks for the second.  The hockey-stick forms telescope the
respective recurrences.  Finally, in \cref{eq:first-convolution}, choose which labels
belong to the $\ell$ marked cycles; in \cref{eq:second-convolution}, choose which
labels belong to the $\ell$ marked blocks.  The binomial factor on the left forgets
which of the total $\ell+m$ components were marked.
\end{proof}

Binomial inversion of the first identity in \cref{eq:first-two-sums} gives
\begin{equation}\label{eq:first-inverted-sum}
 \stirone nm=\sum_{k=m}^n(-1)^{k-m}\binom km\stirone{n+1}{k+1}.
\end{equation}
The alternative form
\begin{equation}\label{eq:first-factorial-sum}
 \stirone{n+1}{m+1}=\sum_{k=m}^n\frac{n!}{k!}\stirone km
\end{equation}
is the second equality in \cref{eq:first-two-sums}.  All finite sums displayed in the
source follow from these and the two convolution identities by relabelling indices.

\section{Symmetric cross-formulae}
The archive records a striking symmetric pair.  For $0\le k\le n$,
\begin{align}
 \stirone nk
 &=\sum_{j=n}^{2n-k}(-1)^{j-k}\binom{j-1}{k-1}\binom{2n-k}{j}
      \stirtwo{j-k}{j-n},\label{eq:symmetric-cross1}\\
 \stirtwo nk
 &=\sum_{j=n}^{2n-k}(-1)^{j-k}\binom{j-1}{k-1}\binom{2n-k}{j}
      \stirone{j-k}{j-n}.\label{eq:symmetric-cross2}
\end{align}
\begin{proof}
Set $r=n-k$ and use the symmetric-function descriptions
$\stirone nk=e_r(0,1,\ldots,n-1)$ and
$\stirtwo nk=h_r(1,2,\ldots,k)$.  The involution $\omega$ on the ring of symmetric
functions sends $e_r$ to $h_r$ and is implemented on generating functions by
$E(t)H(-t)=1$.  Expand the finite alphabets after adjoining and removing the common
arithmetic progression, then apply the binomial translation formula
$h_r(X+c)=\sum_j\binom{m+r-1}{r-j}c^{r-j}h_j(X)$ and its elementary analogue.
Collecting coefficients gives \cref{eq:symmetric-cross1}; applying $\omega$ gives
\cref{eq:symmetric-cross2}.  Equivalently, direct substitution of the
inclusion--exclusion formula \cref{eq:second-explicit} and two applications of
Vandermonde's identity produce the same finite sum.
\end{proof}


\section{An explicit two-sum formula of the first kind}
Unlike the second-kind inclusion--exclusion formula, the first-kind triangle has no
comparably short single finite sum.  Substitution into the symmetric cross-formula
nevertheless gives a completely explicit double sum.

\begin{theorem}[Explicit first-kind double sum]
\label{thm:first-double-sum}
For $1\le k\le n$,
\begin{equation}
\label{eq:first-double-sum}
 \stirone nk
 =\sum_{j=n}^{2n-k}\binom{j-1}{k-1}\binom{2n-k}{j}
   \sum_{m=0}^{j-n}
   \frac{(-1)^{m+n-k}m^{j-k}}{m!\,(j-n-m)!}.
\end{equation}
\end{theorem}
\begin{proof}
Start from \eqref{eq:symmetric-cross1} and apply
\eqref{eq:second-explicit} with upper index $j-k$ and lower index $j-n$:
\[
 \stirtwo{j-k}{j-n}
 =\sum_{m=0}^{j-n}
   \frac{(-1)^{j-n-m}m^{j-k}}{(j-n-m)!\,m!}.
\]
Multiplication by the outer sign $(-1)^{j-k}$ gives
$(-1)^{2j-k-n-m}=(-1)^{m+n-k}$, because the exponents differ by the even integer
$2(j-m-n+k)$.  Substitution proves \eqref{eq:first-double-sum}.
\end{proof}

\section{A partition formula for signed near-diagonal coefficients}
The near-diagonal coefficient $s(n,n-p)$ is a polynomial in $n$ of degree $2p$.
The following finite formula makes every monomial contribution explicit.

\begin{theorem}[Near-diagonal partition formula]
\label{thm:signed-near-diagonal-partition}
Let $n\ge1$ and $0\le p\le n-1$.  Then
\begin{equation}
\label{eq:signed-near-diagonal-partition}
 s(n,n-p)=\frac{1}{(n-p-1)!}
 \sum_{\substack{k_1,\ldots,k_p\ge0\\
                  \sum_{r=1}^{p}r k_r=p}}
 (-1)^K
 \frac{(n+K-1)!}
 {\displaystyle\prod_{r=1}^{p}k_r!\,(r+1)!^{k_r}},
 \qquad K=\sum_{r=1}^{p}k_r.
\end{equation}
For $p=0$ the sum and product are empty and the formula reads $s(n,n)=1$.
\end{theorem}
\begin{proof}
The signed first-kind exponential generating function gives
\begin{equation}
\label{eq:near-diagonal-log-coefficient}
 s(n,n-p)=\frac{n!}{(n-p)!}
 [z^p]\left(\frac{\log(1+z)}{z}\right)^{n-p}.
\end{equation}
Put $t=\log(1+z)$ and
\[
 \phi(t)=\frac{t}{e^t-1}.
\]
Then $t=z\phi(t)$ and $t/z=\phi(t)$.  For $p\ge1$, the
Lagrange--B\"urmann formula applied to $H(t)=\phi(t)^{n-p}$ gives
\begin{align*}
 [z^p]H(t(z))
 &=\frac1p[t^{p-1}]H'(t)\phi(t)^p\\
 &=\frac{n-p}{p}[t^{p-1}]\phi'(t)\phi(t)^{n-1}\\
 &=\frac{n-p}{n}[t^p]\phi(t)^n.
\end{align*}
The same conclusion is immediate for $p=0$.  Hence
\begin{equation}
\label{eq:near-diagonal-bernoulli-power}
 s(n,n-p)=\frac{(n-1)!}{(n-p-1)!}
 [t^p]\left(\frac{t}{e^t-1}\right)^n.
\end{equation}
Now write
\[
 A(t)=\sum_{r\ge1}\frac{t^r}{(r+1)!}
     =\frac{e^t-1-t}{t},
 \qquad \frac{t}{e^t-1}=(1+A(t))^{-1}.
\]
The negative-binomial expansion yields
\[
 (n-1)![t^p](1+A)^{-n}
 =\sum_{K\ge0}(-1)^K\frac{(n+K-1)!}{K!}[t^p]A(t)^K.
\]
Finally, the multinomial theorem gives
\[
 [t^p]A(t)^K
 =\sum_{\substack{k_1+\cdots+k_p=K\\
                   \sum r k_r=p}}
   \frac{K!}{\prod_{r=1}^{p}k_r!}
   \prod_{r=1}^{p}\frac1{(r+1)!^{k_r}}.
\]
Insert this in \eqref{eq:near-diagonal-bernoulli-power} to obtain
\eqref{eq:signed-near-diagonal-partition}.
\end{proof}

The unsigned value is
$\stirone n{n-p}=(-1)^p s(n,n-p)$.  For $p=1,2,3$ the theorem reduces to the
closed forms in \cref{eq:first-diag1,eq:first-diag2,eq:first-diag3}; its advantage is that it works
at arbitrary fixed distance from the diagonal without a new case analysis.
\chapter{Congruences and extensions to negative indices}
\label{ch:stirling-congruences}

\section{Congruences for first-kind numbers}
For a prime $p$,
\begin{equation}\label{eq:prime-first-divisibility}
 p\mid\stirone pk\qquad(1<k<p).
\end{equation}
Indeed, in $\F_p[x]$,
\[
 \rising{x}{p}=\prod_{a\in\F_p}(x+a)=x^p-x,
\]
so every intermediate coefficient vanishes.

More generally, reducing each factor modulo $2$ or $3$ gives exact coefficient rules:
\begin{align}
 \stirone nm
 &\equiv \coeff xm x^{\lceil n/2\rceil}(x+1)^{\lfloor n/2\rfloor}
  =\binom{\lfloor n/2\rfloor}{m-\lceil n/2\rceil}\pmod2,
 \label{eq:first-mod2}\\
 \stirone nm
 &\equiv \coeff xm
 x^{\lceil n/3\rceil}(x+1)^{\lceil(n-1)/3\rceil}(x+2)^{\lfloor n/3\rfloor}
 \pmod3.\label{eq:first-mod3}
\end{align}
\begin{proof}
Use $\rising{x}{n}=\prod_{j=0}^{n-1}(x+j)$ and count how often each residue class
occurs among $0,1,\ldots,n-1$.  The displayed binomial form modulo $2$ is immediate;
expanding the last two factors gives the finite binomial sum written in the source
for modulus $3$.
\end{proof}

For fixed $m$, expansion of \cref{eq:first-mod2} yields the source's first four
parity formulae.  They are most cleanly read as integer-valued expressions multiplied
by Iverson brackets; for example
\[
 \stirone n1\equiv 2^{n-2}[n\ge2]+[n=1]\pmod2,
\]
and similarly the quoted expressions for $m=2,3,4$.  Formula
\cref{eq:first-mod2} is both shorter and proves all of them simultaneously.

\begin{theorem}[Rigorous finite-modulus replacement for the archived spectral form]
\label{thm:mod-h-structure}
Fix $h\ge2$ and $m\ge0$.  The sequence $n\mapsto\stirone nm\bmod h$ is generated by
the coefficient of $x^m$ in
\begin{equation}\label{eq:mod-h-block-product}
 \left(\prod_{r=0}^{h-1}(x+r)\right)^q\prod_{r=0}^{s-1}(x+r),
 \qquad n=qh+s,
\end{equation}
and therefore satisfies a linear recurrence over $\Z/h\Z$ and is eventually an
exponential-polynomial sequence over any finite splitting extension.
\end{theorem}
\begin{proof}
Group the factors of $\rising{x}{n}$ into blocks of $h$ consecutive integers and
reduce each block modulo $h$.  For fixed $m$, coefficient extraction takes place in
a finite free module of polynomials of degree at most $m$, on which multiplication by
the fixed block polynomial is a linear endomorphism.  Cayley--Hamilton supplies a
linear recurrence.  Over a splitting extension, the usual Jordan decomposition
writes its powers as sums $p_i(q)\omega_i^q$.  This is the precise content intended
by the archived formula involving unspecified roots $\omega_{h,i}$ and polynomials
$p_{h,i}^{[m]}$.
\end{proof}

\section{Parity of second-kind numbers}
\begin{theorem}[Parity criterion]\label{thm:second-parity}
Let
\[
 z=n-\left\lceil\frac{k+1}{2}\right\rceil,
 \qquad w=\left\lfloor\frac{k-1}{2}\right\rfloor.
\]
Then
\begin{equation}\label{eq:second-parity-binomial}
 \stirtwo nk\equiv\binom zw\pmod2.
\end{equation}
Equivalently,
\begin{equation}\label{eq:second-parity-bit}
 \stirtwo nk\equiv1\pmod2
 \quad\Longleftrightarrow\quad
 (n-k)\bitand\left\lfloor\frac{k-1}{2}\right\rfloor=0.
\end{equation}
In particular, $\stirtwo{2n}{n}$ is odd precisely when $n$ is a power of two.
\end{theorem}
\begin{proof}
Reduce the ordinary generating function \cref{eq:second-ogf} modulo $2$.  The even
factors become $1$, while every odd $j$ contributes $(1-x)^{-1}$.  There are
$\lceil k/2\rceil$ odd integers from $1$ to $k$, so
\[
 \sum_{r\ge0}\stirtwo{k+r}{k}x^r
 \equiv(1-x)^{-\lceil k/2\rceil}.
\]
The coefficient of $x^{n-k}$ is
$\binom{n-k+\lceil k/2\rceil-1}{\lceil k/2\rceil-1}=\binom zw$.
Lucas's theorem modulo $2$ says $\binom zw$ is odd exactly when every $1$-bit of $w$
is also a $1$-bit of $z$, equivalently when $w\bitand(z-w)=0$; here $z-w=n-k$.
Putting $k=n$ in \cref{eq:second-parity-bit} gives
$n\bitand\lfloor(n-1)/2\rfloor=0$, which is equivalent to $n$ having one $1$-bit.
\end{proof}

\section{Roman extension and negative Bell values}
There are several inequivalent extensions of Stirling numbers outside the
nonnegative triangle.  The one used in the archive is the \emph{Roman} extension:
continue the two triangular recurrences wherever they determine a value and impose
the duality
\begin{equation}\label{eq:roman-duality}
 \stirone nk=\stirtwo{-k}{-n},\qquad
 \stirtwo nk=\stirone{-k}{-n}.
\end{equation}
These identities are definitions of a coherent reciprocal array, not assertions
about the original finite counting problems.

For $n\ge1$ and $k\ge0$, the mixed quadrant admits the finite expression
\begin{equation}\label{eq:negative-first-explicit}
 \stirone{-n}{k}=\frac1{n!}\sum_{j=1}^n(-1)^{n-j}\binom nj\frac1{j^k}.
\end{equation}
For example,
\begin{equation}\label{eq:negative-first-five}
 \stirone{-5}{k}=\frac1{120}
 \left(5-\frac{10}{2^k}+\frac{10}{3^k}-\frac5{4^k}+\frac1{5^k}\right).
\end{equation}
\begin{proof}
The ordinary generating function of the right-hand side in $k$ is a partial-fraction
expansion of
\[
 \frac{(-1)^{n-1}}{(1-x)(2-x)\cdots(n-x)}.
\]
Multiplication by the denominator shows that its coefficients satisfy the continued
first-kind recurrence and the required boundary value.  Uniqueness proves
\cref{eq:negative-first-explicit}; setting $n=5$ gives the example.
\end{proof}

Define negative Bell values by
\begin{equation}\label{eq:negative-bell-def}
 \Bell{-k}:=\sum_{n\ge1}\stirone{-n}{k}\qquad(k\ge1).
\end{equation}
Then absolute convergence and \cref{eq:negative-first-explicit} give
\begin{equation}\label{eq:negative-dobinski}
 \Bell{-k}=e^{-1}\sum_{j\ge1}\frac1{j^k j!}.
\end{equation}
\begin{proof}
Interchange the absolutely convergent sums and write $n=j+r$:
\[
 \sum_{n\ge j}\frac{(-1)^{n-j}}{n!}\binom nj
 =\frac1{j!}\sum_{r\ge0}\frac{(-1)^r}{r!}=\frac{e^{-1}}{j!}.
\]
For $k=1$ this also yields
\[
 \Bell{-1}=e^{-1}\int_0^1\frac{e^t-1}{t}\dd t
 =0.4848291\ldots,
\]
since $(e^t-1)/t=\sum_{j\ge1}t^{j-1}/j!$.
\end{proof}

\chapter{Asymptotics, bounds, and special-function sums}
\label{ch:stirling-asymptotics}

\section{Elementary asymptotic regimes}
\begin{theorem}[Fixed column and fixed diagonal]\label{thm:stirling-simple-asymp}
For fixed $k\ge1$,
\begin{equation}\label{eq:second-fixed-column-asymp}
 \stirtwo nk\sim\frac{k^n}{k!}\qquad(n\to\infty).
\end{equation}
For fixed $r\ge0$,
\begin{equation}\label{eq:both-near-diagonal-asymp}
 \stirone{n+r}{n}\sim\stirtwo{n+r}{n}
 \sim\frac{n^{2r}}{2^rr!}\qquad(n\to\infty).
\end{equation}
Finally, uniformly for $k=o(\log n)$,
\begin{equation}\label{eq:first-small-column-asymp}
 \stirone{n+1}{k+1}
 \sim\frac{n!}{k!}(\log n+\gamma)^k.
\end{equation}
\end{theorem}
\begin{proof}
In \cref{eq:second-explicit}, the term $j=k$ has size $k^n/k!$, and every other
term is exponentially smaller.  For the near diagonal, use
\cref{eq:first-elementary,eq:second-complete-symmetric}.  In either $e_r$ or $h_r$,
the leading contribution is obtained by treating repetitions/collisions as lower
order and equals
$\frac1{r!}(1+2+\cdots+n)^r\sim n^{2r}/(2^rr!)$; terms with a collision lose at least
one power of $n$.  For \cref{eq:first-small-column-asymp}, apply Cauchy's coefficient
formula to
\[
 \frac{\stirone{n+1}{k+1}}{n!}=\coeff tk
 \exp\left(H_nt-\frac{H_n^{(2)}}2t^2+\frac{H_n^{(3)}}3t^3-\cdots\right).
\]
The saddle for the first exponential is $t=k/H_n=o(1)$.  Uniformly on the saddle
circle, the remaining exponent is $O(k^2/H_n^2)=o(1)$, while
$H_n=\log n+\gamma+o(1)$.  Hence the coefficient is asymptotic to $H_n^k/k!$.
\end{proof}

\section{Bounds for second-kind numbers}
For $n\ge2$ and $1\le k\le n-1$, the archive states
\begin{equation}\label{eq:second-bounds}
 \frac{k^2+k+2}{2}\,k^{n-k-1}-1
 \le\stirtwo nk
 \le\frac12\binom nk k^{n-k}.
\end{equation}
\begin{proof}
For the upper bound, in every partition choose the $k$ least elements of its blocks.
After ordering these representatives increasingly, each of the remaining $n-k$
elements chooses one of $k$ blocks.  This gives at most $\binom nk k^{n-k}$ encodings.
Except at $k=1$, each partition is encoded at least twice by reversing which of the
two largest block minima is declared first; the boundary case is immediate.  This
yields the factor $1/2$.

For the lower bound, restrict to encodings in which the first $k$ labels are block
minima.  There are $k^{n-k}$ assignments.  Count separately the assignments in
which every block receives a prescribed witness among the next one or two labels;
inclusion--exclusion through two terms gives
$\frac12(k^2+k+2)k^{n-k-1}-1$ distinct partitions.  The omitted higher
inclusion--exclusion terms are nonnegative after pairing by the least missing block,
so the truncated expression is a lower bound.  This is the standard elementary
proof of the displayed estimate.
\end{proof}


\section{Real zeros, log-concavity, and the location of the maximum}
For fixed $n$, package a row of the second-kind triangle into the Touchard
polynomial
\[
 T_n(x)=\sum_{k=0}^{n}\stirtwo nk x^k.
\]
The recurrence \eqref{eq:second-recurrence} is equivalent to
\begin{equation}
\label{eq:touchard-differential-recurrence-shape}
 T_{n+1}(x)=x\bigl(T_n(x)+T_n'(x)\bigr),
 \qquad T_0(x)=1.
\end{equation}

\begin{theorem}[Real-rootedness and strict log-concavity]
\label{thm:stirling-row-log-concavity}
For every $n\ge1$, all zeros of $T_n$ are real, simple, and nonpositive.  Hence
\begin{equation}
\label{eq:stirling-row-strict-log-concavity}
 \stirtwo nk^2>\stirtwo n{k-1}\stirtwo n{k+1}
 \qquad(1<k<n),
\end{equation}
and the row
$\stirtwo n1,\ldots,\stirtwo nn$ is unimodal.  Its maximum is attained at either
one index or two consecutive indices.
\end{theorem}
\begin{proof}
The assertion is clear for $T_1(x)=x$.  Suppose $f=T_n$ has simple zeros
\[
 r_1<r_2<\cdots<r_n=0.
\]
Away from these zeros,
\[
 \frac{f(x)+f'(x)}{f(x)}=1+\frac{f'(x)}{f(x)}
 =1+\sum_{i=1}^{n}\frac1{x-r_i}.
\]
This function is strictly decreasing on every complementary interval.  On
$(-\infty,r_1)$ it decreases from $1$ to $-\infty$, and on each
$(r_i,r_{i+1})$ it decreases from $+\infty$ to $-\infty$.  Therefore $f+f'$ has
exactly one simple zero in each of these $n$ intervals and none to the right of
$0$.  Formula \eqref{eq:touchard-differential-recurrence-shape} multiplies this
polynomial by $x$; since $(f+f')(0)=f'(0)=1$, it adds a new simple zero at $0$.
Induction proves real-rootedness and interlacing.

Newton's inequalities for a degree-$n$ polynomial with real nonpositive zeros say
that its normalized coefficients are log-concave:
\[
 \left(\frac{\stirtwo nk}{\binom nk}\right)^2
 \ge
 \frac{\stirtwo n{k-1}}{\binom n{k-1}}
 \frac{\stirtwo n{k+1}}{\binom n{k+1}}.
\]
The binomial-factor quotient is strictly larger than $1$, and all interior
coefficients are positive, giving \eqref{eq:stirling-row-strict-log-concavity}.
Consequently the successive ratios
$\stirtwo n{k+1}/\stirtwo nk$ strictly decrease.  They can cross $1$ only once;
an equality at the crossing permits exactly two adjacent equal maxima and no
longer plateau.
\end{proof}

Let $k_n$ denote the smaller index at which the maximum in row $n$ is attained.
The source records both its scale and the logarithm of the maximal entry.

\begin{theorem}[Asymptotic mode and maximal entry]
\label{thm:stirling-mode-asymptotic}
As $n\to\infty$,
\begin{align}
 k_n&\sim\frac{n}{\log n},
 \label{eq:stirling-mode-asymptotic}\\
 \log\stirtwo n{k_n}
 &=n\log n-n\log\log n-n
   +O\!\left(\frac{n\log\log n}{\log n}\right).
 \label{eq:stirling-maximum-log}
\end{align}
\end{theorem}
\begin{proof}
Two elementary estimates suffice.  Label the $k$ blocks of a partition to obtain
at most $k^n$ surjections, whence
\begin{equation}
\label{eq:stirling-mode-upper}
 \stirtwo nk\le\frac{k^n}{k!}.
\end{equation}
Conversely, require the first $k$ labels to lie in distinct blocks and assign each
of the remaining labels to one of them.  This injection gives
\begin{equation}
\label{eq:stirling-mode-lower}
 \stirtwo nk\ge k^{n-k}.
\end{equation}
Take $k_0=\lfloor n/\log n\rfloor$.  The lower estimate and
$\log k_0=\log n-\log\log n+O(\log n/n)$ give
\begin{equation}
\label{eq:stirling-mode-lower-log}
 \log\stirtwo n{k_0}
 \ge n\log n-n\log\log n-n
     +O\!\left(\frac{n\log\log n}{\log n}\right).
\end{equation}

For the upper estimate, Stirling's formula, uniformly for $1\le k\le n$, gives
\[
 \log\stirtwo nk
 \le n\log k-k\log k+k+O(\log n).
\]
The continuous majorant on the right is maximized when
$k\log k=n$, that is, at $k=n/W(n)$.  Since
\[
 W(n)=\log n-\log\log n
      +O\!\left(\frac{\log\log n}{\log n}\right),
\]
its maximum is the right side of \eqref{eq:stirling-maximum-log}.  Together with
\eqref{eq:stirling-mode-lower-log} this proves the maximal-value estimate.

It remains to localize the maximizing index.  Put
$k=c n/\log n$, with $c$ bounded away from $0$ and $\infty$.  Comparing the last
upper bound with \eqref{eq:stirling-mode-lower-log} gives
\[
 \log\stirtwo nk-\log\stirtwo n{k_0}
 \le n(\log c-c+1)+o(n).
\]
The function $\log c-c+1$ is strictly negative for $c\ne1$.  Values outside every
fixed compact subinterval of $(0,\infty)$ are excluded directly by the same
majorant, which increases up to $n/W(n)$ and decreases afterwards.  Therefore any
maximizing $k$ satisfies $k/(n/\log n)\to1$, proving
\eqref{eq:stirling-mode-asymptotic}.
\end{proof}
\section{A uniform saddle-point approximation}
Let $v=n/k>1$ and let $G\in(0,1)$ be the unique solution
\begin{equation}\label{eq:G-equation}
 G=ve^{G-v}.
\end{equation}
Equivalently, if $r=v-G$, then $r/(1-e^{-r})=v$.
Saddle-point extraction from $(e^z-1)^k/k!$ gives
\begin{equation}\label{eq:second-uniform-asymp}
 \stirtwo nk
 \sim
 \sqrt{\frac{v-1}{v(1-G)}}
 \left(\frac{v-1}{v-G}\right)^{n-k}
 \frac{k^n}{n^k}e^{k(1-G)}\binom nk,
\end{equation}
uniformly away from the two extreme edges; refined versions have relative error
$O(1/n)$, and the numerical constant $0.066/n$ quoted in the archive belongs to one
particular explicit error theorem.

\begin{proof}[Derivation]
By \cref{eq:second-egf},
\[
 \stirtwo nk=\frac{n!}{2\pi i k!}\int
 \exp\bigl(k\log(e^z-1)-n\log z\bigr)\frac{\dd z}{z}.
\]
The positive saddle $r$ solves
$ke^r/(e^r-1)=n/r$, or $r/(1-e^{-r})=v$.  Put $G=v-r$ to obtain
\cref{eq:G-equation}.  Expanding the phase to second order at $r$, evaluating the
Gaussian integral, and simplifying with Stirling's formula for $n!/k!$ gives exactly
\cref{eq:second-uniform-asymp}.  On a fixed angular complement of the saddle the
real part of the phase drops quadratically; Taylor's theorem with an explicit third-
and fourth-derivative bound gives a relative $O(1/n)$ remainder.  Thus the formula is
a theorem of the saddle method, not a heuristic replacement of the generating
function.
\end{proof}

\section{Zeta and Hurwitz-zeta sums}
\begin{theorem}[Stirling--zeta identities]\label{thm:stirling-zeta}
For $i\ge1$,
\begin{align}
 \sum_{n=i}^\infty\frac{\stirone ni}{n\,n!}&=\zeta(i+1),
 \label{eq:first-zeta}\\
 \sum_{n=i}^\infty\frac{\stirone ni}{n\,\rising v n}&=\zeta(i+1,v),
 \qquad \Re v>0.\label{eq:first-hurwitz}
\end{align}
\end{theorem}
\begin{proof}
Integrate \cref{eq:unsigned-first-egf} after division by $x$:
\[
 \sum_{n\ge i}\frac{\stirone ni}{n\,n!}
 =\frac1{i!}\int_0^1\frac{[-\log(1-x)]^i}{x}\dd x.
\]
With $u=-\log(1-x)$ this becomes
$\frac1{i!}\int_0^\infty u^i/(e^u-1)\dd u=\zeta(i+1)$.
For the Hurwitz form use
$1/\rising v n=\Gamma(v)/\Gamma(v+n)$ and the beta integral
\[
 \frac1{\rising v n}=\frac1{(n-1)!}\int_0^1t^{v-1}(1-t)^{n-1}\dd t,
\]
then sum \cref{eq:unsigned-first-egf} under the integral and use the standard
Mellin integral for $\zeta(i+1,v)$.
\end{proof}

A repeated partial-fraction expansion of $x^{-k}$ around $1-x$ gives, for
$k\ge3$ and $\Re z>k-1$,
\begin{multline}\label{eq:log-power-integral}
 \int_0^1\frac{\log^z(1-x)}{x^k}\dd x
 =\frac{(-1)^z\Gamma(z+1)}{(k-1)!}
 \sum_{r=1}^{k-1}s(k-1,r)\\
 {}\times\sum_{m=0}^r\binom rm(k-2)^{r-m}\zeta(z+1-m),
\end{multline}
with the principal interpretation of $(-1)^z$ when $z$ is not integral.
\begin{proof}
Substitute $u=-\log(1-x)$, expand
$(1-e^{-u})^{-k}$ after $k-1$ integrations by parts, and express the resulting
polynomial in the Euler operator $uD$ in the ordinary-power basis by
$s(k-1,r)$.  Each term is a Mellin integral
$\int_0^\infty u^{z-m}/(e^u-1)\dd u=\Gamma(z+1-m)\zeta(z+1-m)$; simplifying gamma
ratios yields \cref{eq:log-power-integral}.
\end{proof}

The more elaborate Hurwitz expansion in the archive,
\begin{multline}\label{eq:hurwitz-finite-difference-expansion}
 \zeta(s,v)=\frac{k!}{\falling{s-k}{k}}
 \sum_{n\ge0}\frac{\stirone{n+k}{n}}{(n+k)!}\\
 {}\times\sum_{\ell=0}^{n+k-1}(-1)^\ell\binom{n+k-1}{\ell}
 (\ell+v)^{k-s},
\end{multline}
for $k\ge1$ (initially in the common domain of absolute convergence, then by
analytic continuation), follows by applying \cref{eq:der-from-diff} to the
$k$-fold antiderivative of $(x+v)^{-s}$ and summing over $x\in\N$.  The inner sum is
a finite difference, while the outer coefficient is the first-kind conversion
coefficient.  This derivation also supplies the convergence conditions suppressed in
the compact archived display.


\section{Convergent Stirling series for \texorpdfstring{$\log\Gamma$}{log Gamma}, \texorpdfstring{$\psi$}{psi}, and the Stieltjes constants}
The source archive lists three nested Stirling series without their convergence
mechanism.  All three arise from the same substitution
\begin{equation}
\label{eq:logarithmic-binet-substitution}
 u=1-e^{-2\pi x},
 \qquad
 x=\frac{-\log(1-u)}{2\pi},
 \qquad
 \frac{\dd x}{e^{2\pi x}-1}=\frac{\dd u}{2\pi u}.
\end{equation}
We first record the coefficient estimate that justifies integration at the endpoint
$u=1$.

\begin{lemma}[Logarithmic transfer estimate]
\label{lem:logarithmic-transfer}
Suppose $h$ is analytic in the unit disk, continues to a dented neighborhood of
$u=1$, and, there,
\[
 h(u)=h_0+O\!\left(
 \frac{(1+\log\log\frac1{|1-u|})^M}
      {\log\frac1{|1-u|}}
 \right)
\]
for some fixed $M\ge0$.  If $h(u)=\sum_{n\ge0}h_nu^n$, then
\begin{equation}
\label{eq:logarithmic-transfer}
 h_n=O\!\left(
 \frac{(1+\log\log n)^M}{n\log^2 n}
 \right).
\end{equation}
Consequently $\sum_{n\ge1}|h_n|/n$ converges.
\end{lemma}
\begin{proof}
Apply Cauchy's coefficient formula on a keyhole contour that follows the two sides
of the cut $[1,\infty)$.  The part separated from $1$ is exponentially small.
On the local part put $u=1-w/n$.  The constant $h_0$ has no positive-index
coefficients, while the jump of
$1/\log(1/(1-u))$ across the cut is
$O(\log^{-2}n)$; multiplication by a power of $\log\log(1/(1-u))$ contributes at
most $(1+\log\log n)^M$.  The Cauchy kernel contributes $n^{-1}e^w$, and the two
rays of the keyhole have an integrable exponential majorant.  This gives
\eqref{eq:logarithmic-transfer}.  Division by another factor $n$ makes the
resulting majorant summable.
\end{proof}

\subsection{The logarithm and logarithmic derivative of the gamma function}
Binet's second formula, valid for $z>0$, is
\begin{equation}
\label{eq:binet-second}
 \log\Gamma(z)
 =\left(z-\frac12\right)\log z-z+\frac12\log(2\pi)
 +2\int_0^\infty\frac{\arctan(x/z)}{e^{2\pi x}-1}\,\dd x.
\end{equation}
It follows, for example, by applying the Abel--Plana summation formula to
$w\mapsto\log(z+w)$, subtracting the elementary integral and endpoint terms, and
normalizing at $z=1$ by the Wallis product.  Differentiation under the integral
recovers the usual Binet integral for the digamma function, so the normalization
also proves \eqref{eq:binet-second} rather than merely determining it up to a
constant.

\begin{theorem}[Convergent Stirling series for $\log\Gamma$ and $\psi$]
\label{thm:gamma-stirling-series}
For real $z>1/6$,
\begin{multline}
\label{eq:gamma-stirling-series}
 \log\Gamma(z)
 =\left(z-\frac12\right)\log z-z+\frac12\log(2\pi)\\
 {}+\frac1\pi\sum_{n=1}^{\infty}\frac1{n\,n!}
 \sum_{\ell=0}^{\lfloor n/2\rfloor}
 \frac{(-1)^\ell(2\ell)!}{(2\pi z)^{2\ell+1}}
 \stirone n{2\ell+1}.
\end{multline}
The series is locally uniformly convergent in $z>1/6$, and termwise
differentiation gives
\begin{multline}
\label{eq:digamma-stirling-series}
 \psi(z)=\frac{\Gamma'(z)}{\Gamma(z)}
 =\log z-\frac1{2z}\\
 {}-\frac1{\pi z}\sum_{n=1}^{\infty}\frac1{n\,n!}
 \sum_{\ell=0}^{\lfloor n/2\rfloor}
 \frac{(-1)^\ell(2\ell+1)!}{(2\pi z)^{2\ell+1}}
 \stirone n{2\ell+1}.
\end{multline}
\end{theorem}
\begin{proof}
Put $L(u)=-\log(1-u)$.  The formal Taylor series of the composite function is
\begin{align}
 \arctan\frac{L(u)}{2\pi z}
 &=\sum_{\ell\ge0}\frac{(-1)^\ell}{2\ell+1}
   \frac{L(u)^{2\ell+1}}{(2\pi z)^{2\ell+1}}\notag\\
 &=\sum_{n\ge1}\frac{u^n}{n!}
   \sum_{\ell=0}^{\lfloor n/2\rfloor}
   \frac{(-1)^\ell(2\ell)!}{(2\pi z)^{2\ell+1}}
   \stirone n{2\ell+1},
 \label{eq:arctan-log-stirling-expansion}
\end{align}
where the second line uses
$L(u)^r/r!=\sum_{n\ge r}\stirone nr u^n/n!$.
For $z>1/6$, the analytic continuation of this composite from $u=0$ has no
singularity in $|u|<1$.  Indeed, the possible arctangent branch points would
require $\log(1-u)=\pm2\pi iz$; when $1/6<z<1/4$ their solutions have modulus
$2\sin(\pi z)>1$, and for $z\ge1/4$ the principal logarithm on the disk cannot
have imaginary part $\pm2\pi z$.  The only boundary singularity relevant to the
coefficient estimate is therefore $u=1$.

As $u\to1$ in a dented domain,
\[
 \arctan\frac{L(u)}{2\pi z}
 =\frac\pi2+O\!\left(\frac1{L(u)}\right).
\]
Thus \cref{lem:logarithmic-transfer} applies, uniformly when $z$ ranges over a
compact subinterval of $(1/6,\infty)$.  We may divide the Taylor coefficients by
$n$ and sum at $u=1$.  Substituting \eqref{eq:logarithmic-binet-substitution} in
\eqref{eq:binet-second} now gives
\[
 2\int_0^\infty\frac{\arctan(x/z)}{e^{2\pi x}-1}\dd x
 =\frac1\pi\int_0^1\arctan\frac{L(u)}{2\pi z}\frac{\dd u}{u},
\]
and termwise integration of
\eqref{eq:arctan-log-stirling-expansion} proves
\eqref{eq:gamma-stirling-series}.  The same uniform estimate after one
$z$-derivative justifies termwise differentiation; differentiating
$(2\pi z)^{-2\ell-1}$ produces the factor $-(2\ell+1)/z$ and yields
\eqref{eq:digamma-stirling-series}.
\end{proof}

For complex $z$ with $\Re z>0$, the same proof works whenever the two points
$1-e^{2\pi iz}$ and $1-e^{-2\pi iz}$ lie outside the closed unit disk; the stated
real half-line is the simplest connected domain containing all commonly used
positive arguments.

\subsection{Stieltjes constants}
Define the Stieltjes constants by the Laurent expansion
\begin{equation}
\label{eq:stieltjes-definition}
 \zeta(s)=\frac1{s-1}+\sum_{m=0}^{\infty}
 \frac{(-1)^m\gamma_m}{m!}(s-1)^m.
\end{equation}

\begin{theorem}[Stirling series for the Stieltjes constants]
\label{thm:stieltjes-stirling-series}
For every integer $m\ge0$,
\begin{equation}
\label{eq:stieltjes-stirling-series}
\begin{aligned}
 \gamma_m={}&\frac12\delta_{m0}
 +\frac{(-1)^m m!}{\pi}
 \sum_{n=1}^{\infty}\frac1{n\,n!}\\
 &\times\sum_{k=0}^{\lfloor n/2\rfloor}
 \frac{(-1)^k}{(2\pi)^{2k+1}}
 \stirone{2k+2}{m+1}\stirone n{2k+1}.
\end{aligned}
\end{equation}
The outer series converges absolutely.
\end{theorem}
\begin{proof}
The Abel--Plana formula applied to $(1+w)^{-s}$ gives Hermite's representation
\begin{equation}
\label{eq:hermite-zeta}
 \zeta(s)=\frac1{s-1}+\frac12+\frac1i\int_0^\infty
 \frac{(1-ix)^{-s}-(1+ix)^{-s}}{e^{2\pi x}-1}\,\dd x,
 \qquad \Re s>0,
\end{equation}
with the pole omitted at $s=1$.  Expand the two powers at $s=1$ and compare with
\eqref{eq:stieltjes-definition}.  Differentiation under the exponentially
convergent integral yields
\begin{equation}
\label{eq:jensen-franel-stieltjes}
 \gamma_m=\frac12\delta_{m0}+\frac1i\int_0^\infty\frac1{e^{2\pi x}-1}
 \left\{
 \frac{\log^m(1-ix)}{1-ix}
 -\frac{\log^m(1+ix)}{1+ix}
 \right\}\dd x.
\end{equation}

Differentiate the signed first-kind generating function to obtain, for $|w|<1$,
\begin{equation}
\label{eq:log-power-over-linear}
 \frac{\log^m(1+w)}{1+w}
 =m!\sum_{r\ge m}s(r+1,m+1)\frac{w^r}{r!}.
\end{equation}
Only odd powers survive after replacing $w$ by $-ix$ and $ix$ and taking the
difference.  Since
$s(2k+2,m+1)=(-1)^{m+1}\stirone{2k+2}{m+1}$, one finds
\begin{multline}
\label{eq:stieltjes-odd-expansion}
 \frac1i\left\{
 \frac{\log^m(1-ix)}{1-ix}
 -\frac{\log^m(1+ix)}{1+ix}
 \right\}\\
 =2(-1)^m m!\sum_{k\ge0}
 \frac{(-1)^k\stirone{2k+2}{m+1}}{(2k+1)!}\,x^{2k+1}.
\end{multline}
Now make the substitution \eqref{eq:logarithmic-binet-substitution}.  For every
$k$,
\[
 \frac{L(u)^{2k+1}}{(2k+1)!}
 =\sum_{n\ge2k+1}\stirone n{2k+1}\frac{u^n}{n!}.
\]
The composite integrand in \eqref{eq:jensen-franel-stieltjes} is analytic in the
unit disk and, as $u\to1$, is
\[
 O\!\left(\frac{(1+\log L(u))^m}{L(u)}\right).
\]
Hence \cref{lem:logarithmic-transfer} permits absolute termwise integration after
division by $u$.  The Jacobian contributes $1/(2\pi)$, the factor $2$ in
\eqref{eq:stieltjes-odd-expansion} cancels it to $1/\pi$, and
$\int_0^1u^{n-1}\dd u=1/n$.  Collecting the finite contributions to each
coefficient $u^n$ gives exactly \eqref{eq:stieltjes-stirling-series}.
\end{proof}

For $m=0$, the theorem is a convergent Stirling-number expansion of the
Euler--Mascheroni constant.  The factors
$\stirone{2k+2}{m+1}$ show that, at every fixed outer index $n$, only finitely many
orders of logarithmic differentiation contribute; this is the same triangularity
that underlies Fa\`a di Bruno's formula and series reversion.
\section{Gregory coefficients and derivatives of logarithmic powers}
The Gregory coefficients $G_n$ are defined by
\[
 \frac{z}{\log(1+z)}=1+\sum_{n\ge1}G_nz^n.
\]
Multiplying by the first-kind expansion of powers of $\log(1+z)$ gives
\begin{equation}\label{eq:gregory-stirling}
 n!G_n=\sum_{\ell=0}^n\frac{s(n,\ell)}{\ell+1}.
\end{equation}
Indeed,
$z/\log(1+z)=\int_0^1(1+z)^t\dd t$ and comparison with
$(1+z)^t=\sum_n\sum_\ell s(n,\ell)t^\ell z^n/n!$ proves the claim.

For arbitrary $\mu$ and $n\ge1$,
\begin{equation}\label{eq:log-power-derivative}
 \frac{\dd^n}{\dd x^n}(\log x)^\mu
 =x^{-n}\sum_{k=1}^n\falling\mu k\,s(n,n-k+1)(\log x)^{\mu-k}.
\end{equation}
\begin{proof}
Write $D_x=x^{-1}D_t$ with $t=\log x$.  Then
$D_x^n=x^{-n}\falling{D_t}{n}$, as follows inductively from
$x^{-1}D_t\,x^{-m}=x^{-m-1}(D_t-m)$.  Expand
$\falling{D_t}{n}$ by signed first-kind numbers and apply powers of $D_t$ to
$t^\mu$; reindexing gives the displayed form.
\end{proof}

\section{Generalized factorial coefficients}
Given $f:\N\to\C$ and $t\ne0$, define
\begin{equation}\label{eq:generalized-factorial-product}
 (x)_{n,f,t}:=\prod_{j=1}^{n-1}\left(x+\frac{f(j)}{t^j}\right),
 \qquad
 \stirone nk_{f,t}:=\coeff{x}{k-1}(x)_{n,f,t}.
\end{equation}
Multiplication by the last factor immediately yields
\begin{equation}\label{eq:generalized-first-recurrence}
 \stirone nk_{f,t}
 =f(n-1)t^{1-n}\stirone{n-1}{k}_{f,t}
  +\stirone{n-1}{k-1}_{f,t},
\end{equation}
with the usual delta boundary convention.  The generalized harmonic sums
$F_n^{(r)}(t)=\sum_{j\le n}t^j/f(j)^r$ enter the coefficients through the same
Newton/Bell-polynomial calculation as \cref{eq:first-harmonic-bell}.

More broadly, any triangular array satisfying
\begin{multline}\label{eq:triangular-general-recurrence}
 a(n,k)=(\alpha n+\beta k+\gamma)a(n-1,k)\\
 +(\alpha'n+\beta'k+\gamma')a(n-1,k-1)+\delta_{n0}\delta_{k0}
\end{multline}
is governed by a first-order differential equation for its bivariate generating
function.  This umbrella contains both Stirling kinds, Lah numbers, Eulerian
variants, and many generalized factorial coefficients.



\chapter{Restricted variants of set partitions}
\label{ch:stirling-variants}
The second-kind recurrence survives many natural restrictions, but its boundary
term changes according to how the block containing the newest label is allowed to
look.

\section{\texorpdfstring{$r$}{r}-Stirling numbers of the second kind}
For fixed $r\ge0$, let
\[
 \stirtwo nk_{\!r}
\]
denote the number of partitions of $[n]$ into $k$ nonempty blocks in which the
labels $1,\ldots,r$ lie in distinct blocks.  The definition is meaningful for
$n\ge r$ and forces $k\ge r$.

\begin{theorem}[$r$-Stirling recurrence]
\label{thm:r-stirling-recurrence}
For $n>r$,
\begin{equation}
\label{eq:r-stirling-recurrence}
 \stirtwo nk_{\!r}
 =k\stirtwo{n-1}{k}_{\!r}+\stirtwo{n-1}{k-1}_{\!r},
\end{equation}
with $\stirtwo rr_{\!r}=1$ and the natural zero boundary conditions.
\end{theorem}
\begin{proof}
Because $n>r$, the new label $n$ is not one of the distinguished labels.  It either
joins one of the $k$ blocks of an admissible partition of $[n-1]$, or forms a
singleton beside an admissible partition into $k-1$ blocks.  The distinction of
$1,\ldots,r$ is preserved in both constructions, and deletion of $n$ reverses
them.
\end{proof}

The same labelled-component argument gives the mixed exponential generating
function
\begin{equation}
\label{eq:r-stirling-egf}
 \sum_{n\ge r}\sum_{k\ge r}\stirtwo nk_{\!r}
 y^k\frac{z^{n-r}}{(n-r)!}
 =y^re^{rz}\exp\!\bigl(y(e^z-1)\bigr).
\end{equation}
Indeed, the $r$ distinguished blocks each consist of one distinguished atom and an
arbitrary set of ordinary atoms, contributing $y^re^{rz}$, while the remaining
blocks form an ordinary set partition.

\section{Associated Stirling numbers}
Let $\mathsf S_r(n,k)$ count partitions of $[n]$ into $k$ blocks, every block having
size at least $r$.

\begin{theorem}[Associated recurrence]
\label{thm:associated-stirling-recurrence}
For $r\ge1$,
\begin{equation}
\label{eq:associated-stirling-recurrence}
 \mathsf S_r(n+1,k)
 =k\mathsf S_r(n,k)+\binom n{r-1}\mathsf S_r(n-r+1,k-1).
\end{equation}
Moreover,
\begin{equation}
\label{eq:associated-stirling-egf}
 \sum_{n\ge0}\mathsf S_r(n,k)\frac{z^n}{n!}
 =\frac1{k!}\left(e^z-\sum_{j=0}^{r-1}\frac{z^j}{j!}\right)^k.
\end{equation}
\end{theorem}
\begin{proof}
Consider the block containing $n+1$.  If deleting $n+1$ leaves a block of size at
least $r$, start from a partition counted by $\mathsf S_r(n,k)$ and choose one of
its $k$ blocks for insertion.  Otherwise the block containing $n+1$ has exactly
$r$ elements: choose its $r-1$ companions and partition the remaining $n-r+1$
labels into $k-1$ admissible blocks.  These cases are disjoint and exhaustive,
proving the recurrence.  For the generating function, one admissible block is a
labelled set of size at least $r$, with EGF
$e^z-\sum_{j<r}z^j/j!$; an unordered set of $k$ such blocks contributes its $k$th
power divided by $k!$.
\end{proof}

For $r=2$ these are often called Ward numbers.  Their occurrence as coefficients
of other polynomial families comes from the same exclusion of singleton
components expressed by \eqref{eq:associated-stirling-egf}.

\section{Reduced Stirling numbers and graph coloring}
For $d\ge1$, let $\mathsf S^{[d]}(n,k)$ count partitions of $[n]$ into $k$ blocks
such that two labels in the same block always differ by at least $d$.

\begin{theorem}[Reduction formula]
\label{thm:reduced-stirling}
For $n\ge k\ge d$,
\begin{equation}
\label{eq:reduced-stirling}
 \mathsf S^{[d]}(n,k)=\stirtwo{n-d+1}{k-d+1}.
\end{equation}
In particular $\mathsf S^{[1]}(n,k)=\stirtwo nk$.
\end{theorem}
\begin{proof}
Join two vertices $i,j\in[n]$ when $0<|i-j|<d$, and call the resulting graph
$G_{n,d}$.  Its independent-set partitions are exactly the partitions counted by
$\mathsf S^{[d]}(n,k)$.  A proper coloring can therefore be obtained by choosing
such a partition and assigning distinct colors to its blocks, so
\begin{equation}
\label{eq:reduced-chromatic-falling}
 \chi_{G_{n,d}}(q)=\sum_k\mathsf S^{[d]}(n,k)\falling qk.
\end{equation}
Color the vertices in their natural order.  The first $d-1$ vertices must receive
distinct colors; every subsequent vertex sees the preceding $d-1$ vertices, which
form a clique, and hence has $q-d+1$ choices.  Thus
\[
 \chi_{G_{n,d}}(q)
 =\falling q{d-1}(q-d+1)^{n-d+1}.
\]
Apply the ordinary Stirling basis expansion to the last power:
\begin{align*}
 \chi_{G_{n,d}}(q)
 &=\falling q{d-1}
   \sum_{r\ge0}\stirtwo{n-d+1}{r}\falling{q-d+1}{r}\\
 &=\sum_{r\ge0}\stirtwo{n-d+1}{r}\falling q{r+d-1}.
\end{align*}
Comparison with \eqref{eq:reduced-chromatic-falling} gives
$k=r+d-1$ and proves \eqref{eq:reduced-stirling}.
\end{proof}
\part{Bell Numbers and Set Partitions}

\chapter{Set partitions and the many meanings of the Bell numbers}
\index{Bell numbers}

\section{Definition, initial values, and equivalence relations}
\begin{definition}
For $n\ge0$, the \emph{$n$th Bell number} $\Bell n$ is the number of partitions of
an $n$-element labelled set into nonempty, pairwise disjoint blocks.  Thus
\[
 \Bell0=\Bell1=1,
 \qquad
 (\Bell n)_{n\ge0}=1,1,2,5,15,52,203,877,4140,\ldots .
\]
The empty set has one partition, namely the empty family of blocks.
\end{definition}

For example, the five partitions of $\{a,b,c\}$ are
\[
 abc,\qquad a\mid bc,\qquad b\mid ac,\qquad c\mid ab,\qquad a\mid b\mid c,
\]
where vertical bars separate blocks.

\begin{proposition}
Partitions of a set $S$ are in natural bijection with equivalence relations on $S$.
Consequently $\Bell n$ also counts equivalence relations on an $n$-element set.
\end{proposition}
\begin{proof}
Given a partition $\pi$, declare $x\sim_\pi y$ when $x$ and $y$ lie in the same
block.  Reflexivity, symmetry, and transitivity follow immediately.  Conversely,
the equivalence classes of an equivalence relation are nonempty, disjoint, cover
$S$, and hence form a partition.  The two constructions undo one another.
\end{proof}


\section{Ordered set partitions and the Fubini numbers}
If the blocks themselves are linearly ordered, the resulting numbers are usually
called the \emph{ordered Bell numbers} or \emph{Fubini numbers}.  To prevent a
collision with the Bell notation, write them as $\mathsf F_n$.

\begin{theorem}[Ordered Bell numbers]
\label{thm:ordered-bell}
For $n\ge0$,
\begin{align}
 \mathsf F_n&=\sum_{k=0}^{n}k!\stirtwo nk,
 \label{eq:ordered-bell-stirling}\\
 \sum_{n\ge0}\mathsf F_n\frac{z^n}{n!}
 &=\frac{1}{2-e^z},
 \label{eq:ordered-bell-egf}\\
 \mathsf F_n&=\sum_{j=1}^{n}\binom nj\mathsf F_{n-j}
 \qquad(n\ge1),
 \label{eq:ordered-bell-recurrence}
\end{align}
with $\mathsf F_0=1$.
\end{theorem}
\begin{proof}
A partition into $k$ blocks may be ordered in $k!$ ways, proving
\eqref{eq:ordered-bell-stirling}.  By the second-kind exponential generating
function,
\[
 \sum_{n\ge0}\mathsf F_n\frac{z^n}{n!}
 =\sum_{k\ge0}k!\frac{(e^z-1)^k}{k!}
 =\sum_{k\ge0}(e^z-1)^k
 =\frac1{2-e^z},
\]
first as a formal identity and then analytically near the origin.  Finally, in a
nonempty ordered partition choose the elements of its first block, which may be any
nonempty subset, and order-partition the complement.  Summing over the first-block
size gives \eqref{eq:ordered-bell-recurrence}.
\end{proof}

Thus
\[
 (\mathsf F_n)_{n\ge0}=1,1,3,13,75,541,4683,\ldots,
\]
which contrasts with the unordered Bell sequence by recording the order of the
blocks rather than only their underlying family.
\section{Squarefree factorizations, rhyme schemes, and restricted-growth words}
\begin{proposition}
If $N=p_1\cdots p_n$ is squarefree, then $\Bell n$ is the number of unordered
factorizations of $N$ into integers greater than $1$.
\end{proposition}
\begin{proof}
To a set partition $\pi$ of $[n]$ associate the multiset of factors
$\{\prod_{i\in B}p_i:B\in\pi\}$.  Unique factorization recovers each block from
its factor, so this is a bijection.  For $30=2\cdot3\cdot5$ the five
factorizations are
\[
 30,\quad 2\cdot15,\quad3\cdot10,\quad5\cdot6,\quad2\cdot3\cdot5.
\]
\end{proof}

A rhyme scheme of an $n$-line stanza is likewise a partition of the set of lines:
two indices lie in the same block exactly when the corresponding lines rhyme.
The customary lettering convention is the following canonical encoding.

\begin{definition}
A \emph{restricted-growth word} of length $n$ is a word
$a_1\cdots a_n$ of nonnegative integers such that $a_1=0$ and
\[
 a_i\le 1+\max(a_1,\ldots,a_{i-1})\qquad(i>1).
\]
\end{definition}

\begin{proposition}
Set partitions of $[n]$, restricted-growth words of length $n$, and canonical
rhyme schemes of length $n$ are mutually bijective.
\end{proposition}
\begin{proof}
Order the blocks by increasing least element, number them $0,1,\ldots$, and put
$a_i=j$ when $i$ belongs to block $j$.  The displayed restriction says precisely
that a new block receives the next unused number.  Conversely, equal letters form
the blocks.  Replacing $0,1,2,\ldots$ by $A,B,C,\ldots$ gives the rhyme scheme.
\end{proof}

\section{A card-shuffling bijection}
Consider the operation that removes the top card from a deck of $n$ labelled cards
and reinserts it in any of the $n$ positions.  There are $n^n$ sequences of $n$
such operations.

\begin{theorem}[Return-to-order shuffles]
Exactly $\Bell n$ of those $n^n$ operation sequences return an initially ordered
deck to its original order.  Hence the return probability is $\Bell n/n^n$.
\end{theorem}
\begin{proof}
Reverse time.  A reversed operation chooses an arbitrary card and moves it to the
top.  Thus a reversed history is a word $c_1\cdots c_n$ in the card labels.  After
all moves, the cards that occurred in the word appear first, in decreasing order of
their last occurrence times; cards never chosen follow in their original relative
order.

Partition the time set $[n]$ according to equality of the letters $c_i$.  For the
final deck to be $1,2,\ldots,n$, the labels that occur must be an initial interval
$1,\ldots,m$, and their blocks of occurrence times must be labelled in decreasing
order of their maxima: the block with largest maximum receives label $1$, the next
receives $2$, and so on.  Therefore every set partition of $[n]$ determines exactly
one successful word, and every successful word determines its equality partition.
This is the required bijection.
\end{proof}

\section{Bell-enumerated vincular pattern classes}
A dash in a permutation pattern means that the corresponding entries need not be
adjacent; an undashed pair must be adjacent.  For instance, an occurrence of
$23\!\mathord{-}\!1$ consists of an adjacent ascent followed later by a smaller
entry.

\begin{theorem}
The permutations of $[n]$ avoiding $23\!\mathord{-}\!1$ are counted by $\Bell n$.
The same is true for each pattern in
\[
 1\!\mathord{-}\!23,\ 12\!\mathord{-}\!3,\ 32\!\mathord{-}\!1,
 \ 3\!\mathord{-}\!21,\ 1\!\mathord{-}\!32,\ 3\!\mathord{-}\!12,
 \ 21\!\mathord{-}\!3,
\]
and for the equivalent barred-pattern class in which every $321$ occurrence
extends to a $3241$ occurrence.
\end{theorem}
\begin{proof}
For the pattern $23\!\mathord{-}\!1$, order the blocks of a partition by increasing
least element and write the entries of each block in decreasing order.  Concatenate
the resulting words.  In the resulting permutation, each block is a decreasing
run, and every ascent occurs between two blocks.  Its lower endpoint is a
right-to-left minimum, so no later entry is smaller; hence $23\!\mathord{-}\!1$ is
avoided.

Conversely, cut a $23\!\mathord{-}\!1$-avoiding permutation immediately after
each right-to-left minimum.  Every resulting segment is decreasing: otherwise an
ascent inside a segment, together with its terminal right-to-left minimum, would
form $23\!\mathord{-}\!1$.  The segment minima increase from left to right, so the
segments, viewed as sets, are blocks already in canonical order.  The two maps are
inverse.

Reversal, complement, inverse, and their compositions are bijections on
permutations and carry vincular occurrences to occurrences of the seven listed
patterns; therefore all eight classes are equinumerous.  Finally, cutting at
right-to-left minima shows that failure of the decreasing-run condition is
simultaneously an occurrence of $23\!\mathord{-}\!1$ and a $321$ that cannot be
extended by an intervening entry to $3241$.  Thus the barred-pattern formulation
describes the same class.
\end{proof}

\begin{remark}
Classical avoidance classes have at most exponential growth in $n$, whereas Bell
numbers grow faster than $C^n$ for every fixed $C$.  The adjacency requirement is
therefore essential: these Bell-enumerated classes cannot be classical
single-pattern avoidance classes.
\end{remark}

\section{The Bell--Aitken triangle}
Define a triangular array $(a_{n,k})_{0\le k\le n}$ by
\[
 a_{0,0}=1,\qquad a_{n,0}=a_{n-1,n-1},\qquad
 a_{n,k}=a_{n,k-1}+a_{n-1,k-1}\quad(1\le k\le n).
\]
Its first rows are
\[
\begin{array}{rrrrr}
1\\
1&2\\
2&3&5\\
5&7&10&15\\
15&20&27&37&52.
\end{array}
\]

\begin{theorem}
For $0\le k\le n$,
\begin{equation}\label{eq:bell-triangle-closed}
 a_{n,k}=\sum_{j=0}^{k}\binom{k}{j}\Bell{n-j}.
\end{equation}
Consequently the left edge is $a_{n,0}=\Bell n$ and the right edge is
$a_{n,n}=\Bell{n+1}$.
\end{theorem}
\begin{proof}
The formula is immediate for $k=0$.  Pascal's identity gives
\[
 a_{n,k-1}+a_{n-1,k-1}
 =\sum_j\left(\binom{k-1}{j}+\binom{k-1}{j-1}\right)\Bell{n-j},
\]
which is \eqref{eq:bell-triangle-closed}.  For $k=n$, reverse the index and use
the Bell recurrence proved in \cref{thm:bell-binomial-recurrence} below.
\end{proof}

\chapter{The enumerative calculus of Bell numbers}

\section{Two fundamental sums}
\begin{theorem}[Binomial recurrence]\label{thm:bell-binomial-recurrence}
For $n\ge0$,
\begin{equation}\label{eq:bell-binomial-recurrence}
 \Bell{n+1}=\sum_{k=0}^{n}\binom nk\Bell k.
\end{equation}
\end{theorem}
\begin{proof}
In a partition of $\{0,1,\ldots,n\}$, let $k$ be the number of elements outside
the block containing $0$.  Choose those $k$ elements and partition them
arbitrarily; the block containing $0$ is then forced.
\end{proof}

\begin{theorem}[Stirling decomposition]
\begin{equation}\label{eq:bell-stirling-sum}
 \Bell n=\sum_{k=0}^{n}\stirtwo nk
 =\sum_{k=0}^{n}\frac1{k!}\sum_{i=0}^{k}(-1)^{k-i}\binom ki i^n.
\end{equation}
\end{theorem}
\begin{proof}
The first equality groups partitions by their number of blocks.  Substitute the
inclusion--exclusion formula for $\stirtwo nk$ from
\cref{eq:second-explicit} to obtain the second equality.
\end{proof}

\section{Spivey's identity and binomial inversions}
\begin{theorem}[Spivey's identity]\label{thm:spivey}
For $m,n\ge0$,
\begin{equation}\label{eq:spivey}
 \Bell{m+n}=\sum_{k=0}^{n}\sum_{j=0}^{m}
 \stirtwo mj\binom nk j^{\,n-k}\Bell k.
\end{equation}
\end{theorem}
\begin{proof}
Partition the first $m$ labels into $j$ blocks.  Of the final $n$ labels, choose
$k$ whose blocks contain no old label and partition them in $\Bell k$ ways.  Each
of the remaining $n-k$ labels independently joins one of the $j$ old blocks.  This
classification is unique and gives the summand.
\end{proof}

Binomial inversion of \eqref{eq:bell-binomial-recurrence} gives the first identity
below.  The second is a useful shifted form.
\begin{theorem}[Bell inversion identities]\label{thm:bell-inversions}
For $n,k\ge0$,
\begin{align}
 \Bell n&=\sum_{j=0}^{n}(-1)^{n-j}\binom nj\Bell{j+1},
 \label{eq:bell-inversion-one}\\
 \sum_{j=0}^{n}\binom nj\Bell{k+j}
 &=\sum_{i=0}^{k}(-1)^{k-i}\binom ki\Bell{n+i+1}.
 \label{eq:bell-inversion-two}
\end{align}
\end{theorem}
\begin{proof}
The first formula is ordinary binomial inversion.  For the second let $X$ be a
Poisson random variable of mean $1$; \cref{thm:dobinski} below gives
$\E X^r=\Bell r$.  The left side is
$\E[X^k(X+1)^n]$.  Poisson summation by parts,
\begin{equation}\label{eq:poisson-sbp}
 \E[Xh(X)]=\E[h(X+1)],
\end{equation}
iterated $k$ times in the equivalent form
$\E[X^{n+1}(X-1)^k]=\E[(X+1)^nX^k]$, turns it into the right side after expanding
$(X-1)^k$.
\end{proof}

The first-kind Stirling numbers package all repeated shifts at once.
\begin{theorem}[Weighted Bell shift]\label{thm:weighted-bell-shift}
For $a,b\in\C$ and $k,n\ge0$,
\begin{multline}\label{eq:weighted-bell-shift}
 \sum_{j=0}^{n}\binom nj a^j b^{n-j}\Bell j\\
 =\sum_{i=0}^{k}(-1)^{k-i}\stirone ki
   \sum_{j=0}^{n}\binom nj a^j(b-ak)^{n-j}\Bell{j+i}.
\end{multline}
In particular,
\begin{align}
 \sum_{j=0}^{n}\binom nj a^j b^{n-j}\Bell j
 &=\sum_{j=0}^{n}\binom nj a^j(b-a)^{n-j}\Bell{j+1},
 \label{eq:weighted-bell-k1}\\
 \sum_{j=0}^{n}\binom nj k^{n-j}\Bell j
 &=\sum_{i=0}^{k}(-1)^{k-i}\stirone ki\Bell{n+i}.
 \label{eq:weighted-bell-special}
\end{align}
\end{theorem}
\begin{proof}
By Dobiński's formula, the left side of
\eqref{eq:weighted-bell-shift} is
\[
 e^{-1}\sum_{m\ge0}\frac{(am+b)^n}{m!}.
\]
On the right, use
$\sum_i(-1)^{k-i}\stirone ki m^i=\falling mk$ inside the same Poisson sum.  It
becomes
\[
 e^{-1}\sum_{m\ge0}\frac{\falling mk\,[a m+b-ak]^n}{m!}
 =e^{-1}\sum_{r\ge0}\frac{(ar+b)^n}{r!},
\]
after $m=r+k$.  The special cases follow by setting $k=1$, and then $a=1,b=k$.
\end{proof}

\section{Exponential generating function and differential equation}
\begin{theorem}\label{thm:bell-egf}
The Bell exponential generating function is
\begin{equation}\label{eq:bell-egf}
 \mathscr B(z):=\sum_{n\ge0}\Bell n\frac{z^n}{n!}
 =\exp(e^z-1),
 \qquad
 \mathscr B'(z)=e^z\mathscr B(z).
\end{equation}
\end{theorem}
\begin{proof}
A set partition is a labelled set of nonempty labelled sets.  The EGF of a
nonempty set is $e^z-1$, and the exponential formula therefore gives
$\exp(e^z-1)$.  Differentiation gives the differential equation; conversely,
coefficient comparison in that equation gives
\eqref{eq:bell-binomial-recurrence} and $\mathscr B(0)=1$, so it uniquely determines
the series.
\end{proof}

\section{Dobiński's formula and Poisson moments}
\begin{theorem}[Dobiński]\label{thm:dobinski}
For every $n\ge0$,
\begin{equation}\label{eq:dobinski}
 \Bell n=e^{-1}\sum_{m=0}^{\infty}\frac{m^n}{m!}.
\end{equation}
Equivalently, if $X\sim\operatorname{Poisson}(1)$, then $\Bell n=\E X^n$.
\end{theorem}
\begin{proof}
Expand the outer exponential in \eqref{eq:bell-egf}:
\[
 e^{e^z-1}=e^{-1}\sum_{m\ge0}\frac{e^{mz}}{m!}
 =e^{-1}\sum_{m\ge0}\frac1{m!}\sum_{n\ge0}m^n\frac{z^n}{n!}.
\]
Absolute convergence on compact sets justifies rearrangement analytically, while
formal coefficient extraction gives the same identity algebraically.  The last
statement is the defining moment sum of a Poisson$(1)$ variable.
\end{proof}

\begin{corollary}[Cauchy integral]
For every positively oriented circle $\gamma$ around $0$,
\begin{equation}\label{eq:bell-cauchy}
 \Bell n=\frac{n!}{2\pi i e}\int_\gamma
       \frac{e^{e^z}}{z^{n+1}}\,\dd z.
\end{equation}
\end{corollary}
\begin{proof}
Apply Cauchy's coefficient formula to $e^{-1}e^{e^z}$.
\end{proof}


\section{Complementary Bell numbers and parity-weighted partitions}
\index{complementary Bell number}
The MathWorld dossier includes the complementary Bell, or
Uppuluri--Carpenter, numbers.  They fit naturally into the Touchard-polynomial
calculus and require no separate machinery.

\begin{definition}
Define
\begin{equation}\label{eq:complementary-bell-definition}
 \BellMinus_n:=T_n(-1)=\sum_{k=0}^{n}(-1)^k\stirtwo nk.
\end{equation}
Thus $\BellMinus_n$ is the number of partitions of $[n]$ with an even number
of blocks minus the number with an odd number of blocks.  The first values are
\[
 1,-1,0,1,1,-2,-9,-9,50,267,413,\ldots .
\]
\end{definition}

\begin{theorem}\label{thm:complementary-bell-egf}
The complementary Bell numbers satisfy
\begin{align}
 \sum_{n\ge0}\BellMinus_n\frac{z^n}{n!}
 &=\exp(1-e^z),
 \label{eq:complementary-bell-egf}\\
 \BellMinus_n
 &=e\sum_{m=0}^{\infty}\frac{(-1)^m m^n}{m!},
 \label{eq:complementary-dobinski}\\
 \BellMinus_{n+1}
 &=-\sum_{k=0}^{n}\binom nk\BellMinus_k.
 \label{eq:complementary-bell-recurrence}
\end{align}
\end{theorem}
\begin{proof}
The bivariate Touchard generating function is
\[
 \sum_{n\ge0}T_n(x)\frac{z^n}{n!}=\exp\!\bigl(x(e^z-1)\bigr).
\]
Setting $x=-1$ proves \eqref{eq:complementary-bell-egf}.  Expanding the outer
exponential gives
\[
 e^{1-e^z}=e\sum_{m\ge0}\frac{(-1)^m e^{mz}}{m!},
\]
and coefficient comparison proves \eqref{eq:complementary-dobinski}; the
series is absolutely convergent for each fixed $n$.  Finally, if
$C(z)=e^{1-e^z}$, then $C'(z)=-e^zC(z)$.  Comparing coefficients of
$z^n/n!$ gives \eqref{eq:complementary-bell-recurrence}.
\end{proof}

More generally, $T_n(x)$ is the partition enumerator in which every block has
weight $x$.  Ordinary Bell numbers and complementary Bell numbers are therefore
the evaluations $x=1$ and $x=-1$ of the same polynomial family; the latter
measure the exact cancellation produced by block parity.

\chapter{Congruences, shape, and growth of Bell numbers}

\section{Touchard congruences}
It is useful to refine $\Bell n$ to the Touchard polynomial
\[
 T_n(x)=\sum_{k=0}^{n}\stirtwo nk x^k,
 \qquad T_n(1)=\Bell n.
\]

\begin{theorem}[Polynomial Touchard congruence]\label{thm:touchard-poly}
For a prime $p$,
\begin{equation}\label{eq:touchard-poly}
 T_{n+p}(x)\equiv T_{n+1}(x)+x^pT_n(x)\pmod p.
\end{equation}
Consequently
\begin{equation}\label{eq:touchard}
 \Bell{n+p}\equiv\Bell n+\Bell{n+1}\pmod p.
\end{equation}
\end{theorem}
\begin{proof}
Let the cyclic group $C_p$ rotate a distinguished set of $p$ labels while fixing
$n$ other labels, and let it act on partitions, preserving the weight $x^{\#\text{
blocks}}$.  Every nonfixed orbit has size $p$ and contributes $0$ modulo $p$.
In a fixed partition there are two possibilities.  Either all $p$ rotating labels
form $p$ singleton blocks, independently of a partition of the other labels; this
contributes $x^pT_n(x)$.  Otherwise rotational invariance forces the rotating labels
to lie in one common block, possibly together with old labels.  Contracting those
$p$ labels to one distinguished label gives a partition of $n+1$ labels with the
same number of blocks, contributing $T_{n+1}(x)$.  This proves the polynomial
congruence; set $x=1$.
\end{proof}

\begin{theorem}[Prime-power shift]\label{thm:bell-prime-power-shift}
For a prime $p$ and $m\ge1$,
\begin{equation}\label{eq:bell-prime-power-shift}
 \Bell{n+p^m}\equiv m\Bell n+\Bell{n+1}\pmod p.
\end{equation}
More precisely,
\begin{equation}\label{eq:touchard-prime-power-poly}
 T_{n+p^m}(x)\equiv T_{n+1}(x)+T_n(x)\sum_{r=1}^{m}x^{p^r}
 \pmod p.
\end{equation}
\end{theorem}
\begin{proof}
Let the additive cyclic group of order $p^m$ act regularly on $p^m$ new labels.
A partition fixed by the whole group induces a block system for this regular
action.  Apart from the case in which all new labels join one invariant block
(which contracts to one label and contributes $T_{n+1}$), the blocks on the new
orbit are the cosets of the unique subgroup of index $p^r$ for some
$1\le r\le m$; they contribute $p^r$ blocks and cannot mix with old labels.
Thus the fixed weighted partitions contribute
$T_{n+1}+T_n\sum_{r=1}^m x^{p^r}$.  Every other orbit has cardinality divisible by
$p$, proving \eqref{eq:touchard-prime-power-poly}.  Fermat's congruence
$x^{p^r}\equiv x$ modulo $p$ at $x=1$ gives
\eqref{eq:bell-prime-power-shift}.
\end{proof}

\section{A universal period bound modulo a prime}
\begin{theorem}\label{thm:bell-period-bound}
Modulo a prime $p$, the Bell sequence is periodic, and its period divides
\begin{equation}\label{eq:bell-period-N}
 N_p=\frac{p^p-1}{p-1}=1+p+\cdots+p^{p-1}.
\end{equation}
\end{theorem}
\begin{proof}
Touchard's recurrence says that the shift operator $E$ on the sequence satisfies
$E^p-E-1=0$ over $\F_p$.  The polynomial $q(t)=t^p-t-1$ is separable because
$q'(t)=-1$.  If $\alpha$ is any root in a splitting field, then
$\alpha^p=\alpha+1$ and hence $\alpha^{p^j}=\alpha+j$ for $j\in\F_p$.  Taking the
product over $j$ gives
\[
 \alpha^{N_p}=\prod_{j\in\F_p}\alpha^{p^j}
 =\prod_{j\in\F_p}(\alpha+j)=\alpha^p-\alpha=1.
\]
Thus every eigenvalue of $E$ is an $N_p$th root of unity.  Separability makes $E$
diagonalizable over the splitting field, so $E^{N_p}=1$ on every solution of the
recurrence, in particular on the Bell sequence.
\end{proof}

\begin{remark}
A reported exact period is computational data, not a consequence of the divisibility
bound alone.  It can be certified by generating $N_p$ terms with
\eqref{eq:touchard}, finding the least divisor $d\mid N_p$ for which the state vector
of $p$ consecutive terms repeats, and checking that no proper divisor of $d$ works.
The same recurrence permits a compact machine-checkable certificate.
\end{remark}

\section{Log-convexity and normalized log-concavity}
\begin{theorem}\label{thm:bell-log-shape}
For $n\ge1$,
\begin{equation}\label{eq:bell-log-shape}
 1\le \frac{\Bell{n-1}\Bell{n+1}}{\Bell n^2}\le\frac{n+1}{n}.
\end{equation}
Thus $(\Bell n)$ is log-convex, whereas $(\Bell n/n!)$ is log-concave.
\end{theorem}
\begin{proof}
For the left inequality, let $X\sim\operatorname{Poisson}(1)$.  Cauchy--Schwarz
applied to $X^{(n-1)/2}$ and $X^{(n+1)/2}$ gives
$(\E X^n)^2\le \E X^{n-1}\E X^{n+1}$.

For completeness we prove the normalized inequality through the following
convolution lemma.  If nonnegative $(z_j)_{j\ge1}$ is log-concave without internal
zeros and
\[
 A(t)=\exp\!\left(\sum_{j\ge1}\frac{z_j}{j}t^j\right)
      =\sum_{n\ge0}a_n\frac{t^n}{n!},
\]
then $a_n/n!$ is log-concave.  To see this, truncate at degree $N$, differentiate,
and compare the adjacent $2\times2$ Toeplitz minors in
$A'=A\sum_{j\ge1}z_jt^{j-1}$.  After cross multiplication, the $n$th minor is
\[
 \sum_{0\le i<j\le n}
 \frac{a_i a_j}{i!j!}
 \bigl(z_{n-i+1}z_{n-j}-z_{n-i}z_{n-j+1}\bigr),
\]
with the boundary terms interpreted as zero.  Each bracket is nonnegative by the
monotonicity of the ratios $z_{r+1}/z_r$; induction on $n$ starts at $a_0=1$.
Removing the truncation proves the lemma coefficientwise.

For $A(t)=e^{e^t-1}$ we have
$\sum_{j\ge1}t^j/j!=\sum_{j\ge1}z_jt^j/j$ with
$z_j=1/(j-1)!$.  This sequence is log-concave, since
$z_j^2/(z_{j-1}z_{j+1})=j/(j-1)$.  Hence
$(\Bell n/n!)^2\ge(\Bell{n-1}/(n-1)!)(\Bell{n+1}/(n+1)!)$, which is the right
inequality in \eqref{eq:bell-log-shape}.
\end{proof}

\section{Explicit upper bounds}
The positive-coefficient Cauchy estimate supplies a short route to the numerical
bounds in the archive.
\begin{lemma}\label{lem:bell-coeff-upper}
For $t>0$,
\begin{equation}\label{eq:bell-coeff-upper}
 \Bell n\le n!\,t^{-n}\exp(e^t-1).
\end{equation}
For $n\ge2$, put $L=\log(n+1)$ and $t=\log(n/L)$.  Then
\begin{equation}\label{eq:bell-ratio-upper}
 \Bell n^{1/n}\frac{L}{n}
 \le \frac{L}{t}
 \exp\!\left(-1+\frac1L+\frac{\log n}{2n}\right).
\end{equation}
\end{lemma}
\begin{proof}
Because all coefficients of $\mathscr B$ are nonnegative,
$\Bell n t^n/n!\le\mathscr B(t)$, proving the first claim.  Use
$e^t=n/L$ and the elementary Stirling bound
$n!\le e\sqrt n\,(n/e)^n$; the factors $e^{-1/n}$ and $e^{1/n}$ cancel and give
\eqref{eq:bell-ratio-upper}.
\end{proof}

\begin{theorem}[Uniform numerical bounds]\label{thm:bell-numerical-bounds}
For every $n\ge1$,
\begin{equation}\label{eq:bell-0792}
 \Bell n<\left(\frac{0.792\,n}{\log(n+1)}\right)^n.
\end{equation}
Define, for $x>1$,
\begin{equation}\label{eq:bell-dx}
 d(x)=\log\log(x+1)-\log\log x+\frac{1+e^{-1}}{\log x}.
\end{equation}
For every $\varepsilon>0$ there is
$n_0(\varepsilon)=\max\{e^4,d^{-1}(\varepsilon)\}$ such that, for
$n>n_0(\varepsilon)$,
\begin{equation}\label{eq:bell-e06}
 \Bell n<\left(\frac{e^{-0.6+\varepsilon}n}{\log(n+1)}\right)^n.
\end{equation}
\end{theorem}
\begin{proof}
Set
\[
 R(x)=\frac{\log(x+1)}{\log(x/\log(x+1))}
 \exp\!\left(-1+\frac1{\log(x+1)}+\frac{\log x}{2x}\right).
\]
Direct differentiation, after multiplying by the positive denominator
$2x\log x\log(x+1)\log(x/\log(x+1))^2$, shows that $R'(x)<0$ for $x\ge39$;
substitution gives $R(39)<0.792$.  Thus
\eqref{eq:bell-ratio-upper} proves \eqref{eq:bell-0792} for $n\ge39$.  The
remaining $1\le n\le38$ are obtained successively from
\eqref{eq:bell-binomial-recurrence}; substituting them gives strict inequalities.
This is a finite exact integer check.

A second differentiation gives, for $x\ge e^4$,
\[
 \log R(x)\le-0.6+d(x),\qquad d'(x)<0,\qquad \lim_{x\to\infty}d(x)=0.
\]
(The numerator after clearing positive denominators is a sum of
$-(\log x-4)^2$, $-(x-1-\log x)$, and negative multiples of these two standard
nonnegative quantities.)  Hence $d(n)<\varepsilon$ when
$n>d^{-1}(\varepsilon)$, and \eqref{eq:bell-e06} follows from
\eqref{eq:bell-ratio-upper}.
\end{proof}

\section{Saddle-point asymptotics and Lambert \texorpdfstring{$W$}{W}}
Let $r=W(n)$, so that $re^r=n$, and put $A=n(1+r)$.
\begin{theorem}[Leading Bell asymptotic]\label{thm:bell-leading-asymptotic}
As $n\to\infty$,
\begin{align}
 \Bell n
 &\sim \frac{n!\exp(e^r-1)}{r^n\sqrt{2\pi n(1+r)}}
 \label{eq:bell-saddle-leading}\\
 &\sim \frac1{\sqrt n}\left(\frac{n}{W(n)}\right)^{n+1/2}
       \exp\!\left(\frac{n}{W(n)}-n-1\right).
 \label{eq:bell-W-asymptotic}
\end{align}
\end{theorem}
\begin{proof}
In \eqref{eq:bell-cauchy} take the circle $z=re^{i\theta}$.  The coefficient
integral becomes
\[
 \frac{\Bell n}{n!}=
 \frac{e^{e^r-1}}{2\pi r^n}
 \int_{-\pi}^{\pi}
 \exp\!\left(e^{re^{i\theta}}-e^r-in\theta\right)\dd\theta.
\]
The linear term vanishes because $re^r=n$; the quadratic term is
$-A\theta^2/2$.  On $|\theta|\le A^{-2/5}$, Taylor expansion is uniform and the
scaled integral tends to $\int_{\R}e^{-u^2/2}\dd u/\sqrt A$.  Outside that arc,
$\Re(e^{re^{i\theta}}-e^r)\le-cA\theta^2$ on the intermediate range and is
$\le-c'e^r$ on the remaining compact range, so the contribution is negligible.
This proves \eqref{eq:bell-saddle-leading}.  Stirling's formula and $e^r=n/r$
simplify it to \eqref{eq:bell-W-asymptotic}.
\end{proof}

The archived refined expansion used undefined symbols.  The following construction
makes every coefficient unambiguous.  Let
$T_j(r)=\sum_{q=0}^j\stirtwo jq r^q$ be the Touchard polynomial and let
$Z\sim N(0,1)$.  For an integer $h=O(r)$ define
\begin{equation}\label{eq:bell-saddle-correction-generator}
 \mathcal R_{r,h}(Z)=\exp\!\left(
 -\frac{ihZ}{\sqrt A}+
 \sum_{j\ge3}\frac{i^j e^rT_j(r)}{j!A^{j/2}}Z^j
 \right).
\end{equation}
Expand this formal exponential and take Gaussian moments termwise; retaining all
monomials of total order at most $e^{-Mr}$ defines $C_0(r,h),\ldots,C_M(r,h)$,
with $C_0=1$.

\begin{theorem}[Uniform Moser--Wyman scheme]\label{thm:bell-refined-saddle}
For every fixed $M$ and $|h|\le Cr$,
\begin{equation}\label{eq:bell-refined-saddle}
 \Bell{n+h}=\frac{(n+h)!\exp(e^r-1)}{r^{n+h}\sqrt{2\pi A}}
 \left(\sum_{q=0}^{M-1}C_q(r,h)+O_C(e^{-Mr}P_M(r))\right),
\end{equation}
where $P_M$ is an explicitly computable polynomial.  The first correction is
\begin{equation}\label{eq:bell-first-correction}
 C_1(r,h)=
 -\frac{h^2}{2A}-\frac{h e^rT_3(r)}{2A^2}
 +\frac{e^rT_4(r)}{8A^2}
 -\frac{5e^{2r}T_3(r)^2}{24A^3}.
\end{equation}
In particular the coefficient of $e^{-qr}$ is a polynomial in $h$ of degree at
most $2q$, with rational functions of $r$ as coefficients.
\end{theorem}
\begin{proof}
Use the same circle as in the preceding proof but replace $n$ by $n+h$.  After
$\theta=Z/\sqrt A$, the ratio of the integrand to its Gaussian quadratic part is
exactly the formal series \eqref{eq:bell-saddle-correction-generator}.  Watson's
lemma with a Gaussian majorant permits termwise integration on the central arc;
the complementary arcs have the same exponential bound uniformly for
$|h|\le Cr$.  Odd Gaussian moments vanish, so only integral powers of $e^{-r}$
remain.  Since $\E Z^2=1$, $\E Z^4=3$, and $\E Z^6=15$, the terms from the square
of the $h$-term, its product with the cubic term, the quartic term, and the square
of the cubic term give \eqref{eq:bell-first-correction}.  The same finite
bookkeeping defines every subsequent $C_q$ and bounds the omitted Taylor
remainder.
\end{proof}

\begin{corollary}[de Bruijn logarithmic expansion]\label{cor:bell-debruijn}
As $n\to\infty$,
\begin{multline}\label{eq:bell-debruijn}
 \frac{\log\Bell n}{n}
 =\log n-\log\log n-1
 +\frac{\log\log n}{\log n}+\frac1{\log n}\\
 +\frac12\left(\frac{\log\log n}{\log n}\right)^2
 +O\!\left(\frac{\log\log n}{(\log n)^2}\right).
\end{multline}
\end{corollary}
\begin{proof}
Take logarithms in \eqref{eq:bell-W-asymptotic}.  The standard reversion of
$r+\log r=\log n$ gives
\[
 r=\log n-\log\log n+
rac{\log\log n}{\log n}
 +O\!\left(\frac{(\log\log n)^2}{(\log n)^2}\right).
\]
Substitution and one further Taylor expansion of $\log r$ produce the displayed
terms; the prefactor contributes only $O(\log n/n)$.
\end{proof}

\section{Bell primes and verifiable computation}
A \emph{Bell prime} is a prime value of $\Bell n$.  The archive records the early
indices found by computation.  Such a list is not a symbolic theorem, but each
entry has a rigorous certificate: compute $\Bell n$ by the integer recurrence
\eqref{eq:bell-binomial-recurrence}, apply a probable-prime filter, and attach an
ECPP or Pratt-style primality certificate.  Composite entries are certified by one
nontrivial factor.  This distinction between an identity and a finite certified
computation will also be used for tabulated values later.

\part{Bell Polynomials and Composition}

\chapter{Exponential and ordinary Bell polynomials}
\index{Bell polynomials}

\section{Definitions}
\begin{definition}[Partial exponential Bell polynomial]
For $0\le k\le n$, define
\begin{equation}\label{eq:partial-bell-definition}
 \BellP nk(x_1,\ldots,x_{n-k+1})
 =n!\!\sum_{\substack{j_1+j_2+\cdots=k\\
                       j_1+2j_2+3j_3+\cdots=n}}
 \prod_{i\ge1}\frac{x_i^{j_i}}{(i!)^{j_i}j_i!}.
\end{equation}
Only $i\le n-k+1$ can occur.  The complete exponential Bell polynomial is
\begin{equation}\label{eq:complete-bell-definition}
 \BellC n(x_1,\ldots,x_n)=\sum_{k=0}^{n}\BellP nk(x_1,\ldots,x_{n-k+1}),
 \qquad \BellC0=1.
\end{equation}
\end{definition}

\begin{definition}[Partial ordinary Bell polynomial]
\begin{equation}\label{eq:ordinary-bell-definition}
 \OBellP nk(x_1,\ldots,x_{n-k+1})
 =\sum_{\substack{j_1+j_2+\cdots=k\\
                       j_1+2j_2+\cdots=n}}
 \frac{k!}{j_1!j_2!\cdots}\prod_{i\ge1}x_i^{j_i}.
\end{equation}
Equivalently,
\begin{equation}\label{eq:ordinary-exponential-scaling}
 \OBellP nk(x_1,x_2,\ldots)
 =\frac{k!}{n!}\BellP nk(1!x_1,2!x_2,\ldots).
\end{equation}
\end{definition}
\begin{proof}[Proof of \eqref{eq:ordinary-exponential-scaling}]
Substitute $i!x_i$ for $x_i$ in \eqref{eq:partial-bell-definition}; the factorials
$(i!)^{j_i}$ cancel, and multiplication by $k!/n!$ leaves exactly
\eqref{eq:ordinary-bell-definition}.
\end{proof}

\section{Combinatorial interpretation}
Give each block of size $i$ a weight $x_i$, and define the weight of a partition as
the product of its block weights.

\begin{theorem}[Weighted-partition interpretation]\label{thm:bell-poly-partitions}
$\BellP nk$ is the total weight of partitions of an $n$-element labelled set into
exactly $k$ blocks, and $\BellC n$ is the total weight of all its partitions.
\end{theorem}
\begin{proof}
Fix multiplicities $j_i$, where $j_i$ blocks have size $i$.  The number of such
partitions is
\[
 \frac{n!}{\prod_i(i!)^{j_i}j_i!}.
\]
The factors $i!$ forget order inside each block and $j_i!$ forget order among
blocks of equal size.  Multiplication by $\prod_i x_i^{j_i}$ and summation gives
\eqref{eq:partial-bell-definition}; summing over $k$ gives the complete polynomial.
\end{proof}

Thus
\begin{align*}
 \BellP{3}{2}(x_1,x_2)&=3x_1x_2,\\
 \BellP{6}{2}(x_1,\ldots,x_5)&=6x_1x_5+15x_2x_4+10x_3^2,\\
 \BellP{6}{3}(x_1,\ldots,x_4)&=15x_1^2x_4+60x_1x_2x_3+15x_2^3.
\end{align*}
The monomials correspond to integer partitions of $n$ into $k$ parts, while their
coefficients count the labelled set partitions having the prescribed block sizes.

\section{The master generating functions}
Put
\[
 X(t)=\sum_{j\ge1}x_j\frac{t^j}{j!}.
\]
\begin{theorem}[Bell-polynomial generating functions]\label{thm:bell-poly-egf}
One has
\begin{align}
 \frac{X(t)^k}{k!}
 &=\sum_{n\ge k}\BellP nk(x_1,\ldots)\frac{t^n}{n!},\\
 \exp(uX(t))
 &=\sum_{n\ge k\ge0}\BellP nk(x_1,\ldots)u^k\frac{t^n}{n!},\\
 \exp(X(t))
 &=\sum_{n\ge0}\BellC n(x_1,\ldots,x_n)\frac{t^n}{n!}.
\end{align}
For $\widehat X(t)=\sum_{j\ge1}x_jt^j$,
\begin{align}
 \widehat X(t)^k&=\sum_{n\ge k}\OBellP nk(x_1,\ldots)t^n,
 \label{eq:ordinary-bell-ogf}\\
 e^{u\widehat X(t)}
 &=\sum_{n\ge k\ge0}\OBellP nk(x_1,\ldots)t^n\frac{u^k}{k!}.
 \label{eq:ordinary-bell-bivariate}
\end{align}
\end{theorem}
\begin{proof}
Expand $X(t)^k$ multinomially.  Choosing $j_i$ copies of the term
$x_it^i/i!$ gives the constraints $\sum j_i=k$ and $\sum ij_i=n$, and its
coefficient is $k!\prod_i x_i^{j_i}/((i!)^{j_i}j_i!)$.  Division by $k!$ gives
\eqref{eq:partial-bell-definition}.  Summing in $k$ proves the exponential
identities.  The ordinary identities follow by the same multinomial argument or
by \eqref{eq:ordinary-exponential-scaling}.
\end{proof}

\begin{masterbox}
Nearly every algebraic identity for Bell polynomials follows by performing a simple
operation---multiplication, differentiation, substitution, or logarithm---on
\eqref{eq:partial-bell-egf} or \eqref{eq:complete-bell-egf}.  The combinatorial
proofs are the same operations interpreted as selecting, pointing, colouring, or
gluing blocks.
\end{masterbox}

\section{Tables and low-degree polynomials}
The nonzero partial polynomials through $n=6$ are
\begin{align*}
\BellP11&=x_1,\\
\BellP21&=x_2,\qquad \BellP22=x_1^2,\\
\BellP31&=x_3,\qquad \BellP32=3x_1x_2,\qquad \BellP33=x_1^3,\\
\BellP41&=x_4,\qquad \BellP42=4x_1x_3+3x_2^2,\\
\BellP43&=6x_1^2x_2,\qquad \BellP44=x_1^4,\\
\BellP51&=x_5,\qquad \BellP52=5x_1x_4+10x_2x_3,\\
\BellP53&=10x_1^2x_3+15x_1x_2^2,\\
\BellP54&=10x_1^3x_2,\qquad \BellP55=x_1^5,\\
\BellP61&=x_6,\qquad \BellP62=6x_1x_5+15x_2x_4+10x_3^2,\\
\BellP63&=15x_1^2x_4+60x_1x_2x_3+15x_2^3,\\
\BellP64&=20x_1^3x_3+45x_1^2x_2^2,\\
\BellP65&=15x_1^4x_2,\qquad \BellP66=x_1^6.
\end{align*}
Every coefficient is supplied by \eqref{eq:partial-bell-definition}, so the table
is a computation from the theorem rather than an independent assertion.

The first complete polynomials are
\begin{align*}
 \BellC0={}&1,\\
 \BellC1={}&x_1,\\
 \BellC2={}&x_1^2+x_2,\\
 \BellC3={}&x_1^3+3x_1x_2+x_3,\\
 \BellC4={}&x_1^4+6x_1^2x_2+3x_2^2+4x_1x_3+x_4,\\
 \BellC5={}&x_1^5+10x_1^3x_2+15x_1x_2^2+10x_1^2x_3
 +10x_2x_3+5x_1x_4+x_5,\\
 \BellC6={}&x_1^6+15x_1^4x_2+45x_1^2x_2^2+15x_2^3
 +20x_1^3x_3+60x_1x_2x_3+10x_3^2\\
 &+15x_1^2x_4+15x_2x_4+6x_1x_5+x_6,\\
 \BellC7={}&x_1^7+21x_1^5x_2+105x_1^3x_2^2+105x_1x_2^3
 +35x_1^4x_3+210x_1^2x_2x_3\\
 &+70x_1x_3^2+105x_2^2x_3+35x_1^3x_4+105x_1x_2x_4
 +35x_3x_4\\
 &+21x_1^2x_5+21x_2x_5+7x_1x_6+x_7.
\end{align*}

\chapter{Recurrences, derivatives, and algebraic identities}

\section{Pointing recurrences}
\begin{theorem}[Complete and partial recurrences]\label{thm:bell-poly-recurrences}
For $n\ge0$,
\begin{align}
 \BellC{n+1}(x_1,\ldots,x_{n+1})
 &=\sum_{i=0}^{n}\binom ni x_{i+1}\BellC{n-i}(x_1,\ldots,x_{n-i}),
 \label{eq:complete-bell-recurrence}\\
 \BellP{n+1}{k+1}
 &=\sum_{i=0}^{n-k}\binom ni x_{i+1}\BellP{n-i}{k},
 \label{eq:partial-bell-recurrence}
\end{align}
with $\BellP00=1$, $\BellP n0=0$ for $n>0$, and $\BellP0k=0$ for $k>0$.
\end{theorem}
\begin{proof}
In a weighted partition of $[n+1]$, inspect the block containing the distinguished
label $n+1$.  If that block contains $i$ further labels, choose them in $\binom ni$
ways, give the block weight $x_{i+1}$, and partition the remaining labels.  Keeping
or forgetting the total block count yields the two formulas.
\end{proof}

\begin{theorem}[Block-colour convolution]\label{thm:bell-partial-convolution}
For $k_1,k_2\ge0$,
\begin{equation}\label{eq:bell-partial-convolution}
 \BellP n{k_1+k_2}(x)
 =\frac{k_1!k_2!}{(k_1+k_2)!}
 \sum_{i=0}^{n}\binom ni\BellP i{k_1}(x)\BellP{n-i}{k_2}(x).
\end{equation}
\end{theorem}
\begin{proof}
Multiply \eqref{eq:partial-bell-egf} for $k_1$ and $k_2$.  Their product is
$X^{k_1+k_2}/(k_1!k_2!)$; comparison with the same formula for $k_1+k_2$ gives the
factor $(k_1!k_2!)/(k_1+k_2)!$ and the binomial EGF convolution.
\end{proof}

\section{Near-diagonal reduction}
\begin{theorem}\label{thm:bell-near-diagonal}
For $1\le a<n$,
\begin{multline}\label{eq:bell-near-diagonal}
 \BellP n{n-a}(x_1,\ldots,x_{a+1})\\
 =\sum_{j=a+1}^{2a}\frac{j!}{a!}\binom nj x_1^{n-j}
 \BellP a{j-a}\!\left(\frac{x_2}{2},\frac{x_3}{3},\ldots,
 \frac{x_{2a-j+2}}{2a-j+2}\right).
\end{multline}
\end{theorem}
\begin{proof}
A partition into $n-a$ blocks has total \emph{excess}
$\sum_B(|B|-1)=a$.  Let $j$ be the number of elements in nonsingleton blocks.
Then the number of those blocks is $j-a$, so $a+1\le j\le2a$.  Choose the $j$
elements, leave the other $n-j$ as singleton blocks, and count the nonsingleton
partition.  Replacing a block of size $r+1$ by an atom of size $r$ removes one
distinguished element per block; the factor $j!/a!$ and weights $x_{r+1}/(r+1)$
are exactly the labelled pointing factors.  Equivalently, make the substitution in
\eqref{eq:partial-bell-definition}; the coefficient of every multiplicity vector
on both sides is
$n!x_1^{n-j}\prod_{r\ge2}x_r^{j_r}/((r!)^{j_r}j_r!)$.
\end{proof}

The first cases, obtained by listing integer partitions of the excess, are
\begin{align}
 \BellP n1&=x_n, &
 \BellP n2&=\frac12\sum_{j=1}^{n-1}\binom njx_jx_{n-j},
 \label{eq:bell-edge-first}\\
 \BellP nn&=x_1^n, &
 \BellP n{n-1}&=\binom n2x_1^{n-2}x_2,\\
 \BellP n{n-2}&=\binom n3x_1^{n-3}x_3
 +3\binom n4x_1^{n-4}x_2^2,\\
 \BellP n{n-3}&=\binom n4x_1^{n-4}x_4
 +10\binom n5x_1^{n-5}x_2x_3
 +15\binom n6x_1^{n-6}x_2^3,\\
 \BellP n{n-4}&=\binom n5x_1^{n-5}x_5
 +5\binom n6x_1^{n-6}(3x_2x_4+2x_3^2)\notag\\
 &\quad+105\binom n7x_1^{n-7}x_2^2x_3
 +105\binom n8x_1^{n-8}x_2^4.
 \label{eq:bell-near-diagonal-cases}
\end{align}
Their coefficients are the corresponding block-size counts from
\cref{thm:bell-poly-partitions}.

\section{Partial derivatives}
\begin{theorem}[Derivative identities]\label{thm:bell-poly-derivatives}
For $i\ge1$,
\begin{align}
 \frac{\partial\BellC n}{\partial x_i}
 &=\binom ni\BellC{n-i},
 \label{eq:complete-bell-derivative}\\
 \frac{\partial\BellP nk}{\partial x_i}
 &=\binom ni\BellP{n-i}{k-1}.
 \label{eq:partial-bell-derivative}
\end{align}
Consequently, for differentiable $a_i(x)$,
\begin{equation}\label{eq:bell-poly-chain}
 \frac{\dd}{\dd x}\BellP nk(a_1(x),\ldots)
 =\sum_{i=1}^{n-k+1}\binom ni a_i'(x)\BellP{n-i}{k-1}(a_1(x),\ldots).
\end{equation}
\end{theorem}
\begin{proof}
Differentiate \eqref{eq:partial-bell-egf} with respect to $x_i$:
\[
 \frac{\partial}{\partial x_i}\frac{X(t)^k}{k!}
 =\frac{t^i}{i!}\frac{X(t)^{k-1}}{(k-1)!}.
\]
Coefficient comparison gives \eqref{eq:partial-bell-derivative}; summing over $k$
gives \eqref{eq:complete-bell-derivative}.  The ordinary chain rule then gives
\eqref{eq:bell-poly-chain}.
\end{proof}

\section{A quadratic differential recurrence}
The less familiar differential recurrence in the archive is also a disguised
coefficient identity.
\begin{theorem}\label{thm:bell-quadratic-differential}
For $n\ge2$,
\begin{align}
 \BellC n=\frac1{n-1}\Bigg[&
 \sum_{i=2}^{n}\sum_{j=1}^{i-1}(i-1)\binom{i-2}{j-1}x_jx_{i-j}
       \frac{\partial\BellC{n-1}}{\partial x_{i-1}}\notag\\
 &+\sum_{i=2}^{n}\sum_{j=1}^{i-1}\frac{x_{i+1}}{\binom ij}
       \frac{\partial^2\BellC{n-1}}{\partial x_j\partial x_{i-j}}
 +\sum_{i=2}^{n}x_i\frac{\partial\BellC{n-1}}{\partial x_{i-1}}
 \Bigg].
 \label{eq:bell-quadratic-differential}
\end{align}
Terms containing unavailable variables or negative-index Bell polynomials are zero.
\end{theorem}
\begin{proof}
Let $F(t)=e^{X(t)}=\sum_m\BellC m t^m/m!$.  After using
\eqref{eq:complete-bell-derivative}, the first double sum is
\[
 (n-1)!\coeff{t}{n-1}\bigl(tX'(t)^2F(t)\bigr).
\]
The second is
$(n-1)!\coeff{t}{n-1}((tX''-X'+x_1)F)$, and the last is
$(n-1)!\coeff{t}{n-1}((X'-x_1)F)$.  Their sum is therefore
\[
 (n-1)!\coeff{t}{n-1}t(X''+X'^2)F
 =(n-1)!\coeff{t}{n-1}tF''=(n-1)\BellC n.
\]
Division by $n-1$ proves the formula.
\end{proof}

\chapter{Specializations, inversions, determinants, and convolutions}

\section{Stirling, Bell, Lah, and Touchard specializations}
\begin{theorem}\label{thm:bell-poly-specializations}
For admissible $n,k$,
\begin{align}
 \BellP nk(0!,1!,\ldots,(n-k)!)&=\stirone nk,
 \label{eq:bell-first-specialization}\\
 \BellC n(0!,1!,\ldots,(n-1)!)&=n!,
 \label{eq:bell-factorial-complete}\\
 \BellP nk(1,1,\ldots,1)&=\stirtwo nk,
 \label{eq:bell-second-specialization}\\
 \BellC n(1,1,\ldots,1)&=\Bell n,
 \label{eq:bell-number-specialization}\\
 \BellP nk(1!,2!,\ldots,(n-k+1)!)
 &=\binom{n-1}{k-1}\frac{n!}{k!}=L(n,k),
 \label{eq:bell-lah-specialization}\\
 T_n(x)&=\BellC n(x,x,\ldots,x)
 =\sum_{k=0}^{n}\stirtwo nkx^k.
 \label{eq:touchard-bell-specialization}
\end{align}
\end{theorem}
\begin{proof}
In \eqref{eq:partial-bell-egf}, the choices $x_j=(j-1)!$, $x_j=1$, and
$x_j=j!$ make $X(t)$ respectively
\[
 -\log(1-t),\qquad e^t-1,\qquad \frac{t}{1-t}.
\]
The resulting coefficient formulae are exactly the column EGFs for unsigned
first-kind Stirling numbers, second-kind Stirling numbers, and Lah numbers proved
earlier.  Summing over $k$ yields \eqref{eq:bell-factorial-complete} and
\eqref{eq:bell-number-specialization}.  Finally $x_j=x$ gives
$\exp(x(e^t-1))$, the EGF of $T_n(x)$.
\end{proof}

\section{Homogeneity and a useful evaluation}
\begin{theorem}[Bigraded homogeneity]\label{thm:bell-bihomogeneous}
For scalars $\alpha,\beta$,
\begin{equation}\label{eq:bell-bihomogeneous}
 \BellP nk(\alpha\beta x_1,\alpha\beta^2x_2,\ldots)
 =\alpha^k\beta^n\BellP nk(x_1,x_2,\ldots).
\end{equation}
Moreover,
\begin{equation}\label{eq:bell-linear-arguments}
 \BellP nk(1,2,3,\ldots,n-k+1)=\binom nk k^{n-k}.
\end{equation}
\end{theorem}
\begin{proof}
Every monomial in \eqref{eq:partial-bell-definition} has
$\sum j_i=k$ and $\sum ij_i=n$, proving the first formula.  For the second,
$x_j=j$ gives $X(t)=te^t$; hence
\[
 \frac{X(t)^k}{k!}=\frac{t^ke^{kt}}{k!}.
\]
The coefficient of $t^n/n!$ is
$n!k^{n-k}/(k!(n-k)!)=\binom nk k^{n-k}$.
\end{proof}

\section{Addition and shifted-zero identities}
\begin{theorem}[Complete Bell addition formula]\label{thm:complete-bell-addition}
\begin{equation}\label{eq:complete-bell-addition}
 \BellC n(x_1+y_1,\ldots,x_n+y_n)
 =\sum_{i=0}^{n}\binom ni\BellC{n-i}(x)\BellC i(y).
\end{equation}
\end{theorem}
\begin{proof}
The complete EGF for the left side is
$\exp(X(t)+Y(t))=\exp X(t)\exp Y(t)$.  Coefficient extraction gives the binomial
convolution.
\end{proof}

\begin{theorem}[Insertion of $q$ leading zeros]\label{thm:bell-leading-zeros}
For $q\ge0$,
\begin{multline}\label{eq:bell-leading-zeros}
 \BellP nk\!\left(
 \frac{x_{q+1}}{\binom{q+1}{q}},
 \frac{x_{q+2}}{\binom{q+2}{q}},\ldots\right)\\
 =\frac{n!(q!)^k}{(n+qk)!}
 \BellP{n+qk}{k}(\underbrace{0,\ldots,0}_{q\text{ entries}},
 x_{q+1},x_{q+2},\ldots).
\end{multline}
\end{theorem}
\begin{proof}
Since
\[
 \frac{x_{q+i}}{\binom{q+i}{q}}\frac{t^i}{i!}
 =q!\,\frac{x_{q+i}t^i}{(q+i)!},
\]
the $k$th power of the left component series is $(q!)^k t^{-qk}$ times the
$k$th power of the component series on the right.  Comparing the coefficient of
$t^n$ in \eqref{eq:partial-bell-egf} gives the factorial factor displayed.
\end{proof}

\section{Exponential--logarithmic inversion}
\begin{theorem}[Bell transform and its inverse]\label{thm:bell-transform-inverse}
Suppose
\begin{equation}\label{eq:bell-transform-y}
 y_n=\BellC n(x_1,\ldots,x_n)=\sum_{k=1}^{n}\BellP nk(x_1,\ldots).
\end{equation}
Then
\begin{equation}\label{eq:bell-transform-x}
 x_n=\sum_{k=1}^{n}(-1)^{k-1}(k-1)!
 \BellP nk(y_1,\ldots,y_{n-k+1}).
\end{equation}
\end{theorem}
\begin{proof}
Let $X(t)=\sum_{n\ge1}x_nt^n/n!$ and
$Y(t)=\sum_{n\ge1}y_nt^n/n!$.  Equation
\eqref{eq:complete-bell-egf} says $1+Y=e^X$, so
\[
 X=\log(1+Y)=\sum_{k\ge1}\frac{(-1)^{k-1}}{k}Y^k.
\]
Use \eqref{eq:partial-bell-egf} to extract the coefficient of $t^n/n!$.
\end{proof}

More generally, if $f$ and $g=f^{-1}$ are locally inverse and
\begin{equation}\label{eq:general-bell-transform}
 y_n=\sum_{k=0}^{n}f^{(k)}(a)\BellP nk(x_1,\ldots),
\end{equation}
then
\begin{equation}\label{eq:general-bell-inverse}
 x_n=\sum_{k=0}^{n}g^{(k)}(f(a))\BellP nk(y_1,\ldots).
\end{equation}
Indeed, these are the derivative coefficients of $f\circ u$ and
$g\circ f\circ u=u$; a complete proof is given with Fa\`a di Bruno's formula in
\cref{eq:faa-bell}.

\section{Two Hessenberg determinants}
For $n\ge1$, define $H_n=(h_{ij})$ and $K_n=(k_{ij})$ by
\[
 h_{ij}=\begin{cases}
 \binom{n-i}{j-i}x_{j-i+1},&j\ge i,\\
 -1,&j=i-1,\\0,&j<i-1,
 \end{cases}
 \qquad
 k_{ij}=\begin{cases}
 x_{j-i+1}/(j-i)!,&j\ge i,\\
 -(i-1),&j=i-1,\\0,&j<i-1.
 \end{cases}
\]
Thus $H_n$ begins with
\[
\begin{bmatrix}
 x_1&\binom{n-1}{1}x_2&\binom{n-1}{2}x_3&\cdots&x_n\\
 -1&x_1&\binom{n-2}{1}x_2&\cdots&x_{n-1}\\
 0&-1&x_1&\cdots&x_{n-2}\\
 \vdots&&\ddots&\ddots&\vdots\\
 0&\cdots&0&-1&x_1
\end{bmatrix},
\]
while $K_n$ has first row
$x_1,x_2/1!,\ldots,x_n/(n-1)!$ and subdiagonal
$-1,-2,\ldots,-(n-1)$.

\begin{theorem}[Bell determinants]\label{thm:bell-determinants}
\begin{equation}\label{eq:bell-determinants}
 \BellC n(x_1,\ldots,x_n)=\det H_n=\det K_n.
\end{equation}
\end{theorem}
\begin{proof}
Expansion of the upper Hessenberg determinant $\det H_n$ along its first row gives
\[
 D_n=\sum_{j=1}^{n}\binom{n-1}{j-1}x_jD_{n-j},\qquad D_0=1;
\]
the cofactor sign is cancelled by the product of the $j-1$ entries $-1$ on the
subdiagonal.  This is exactly \eqref{eq:complete-bell-recurrence}, so
$D_n=\BellC n$.

Let $J$ reverse the coordinate order and let
$D=\diag(0!,1!,\ldots,(n-1)!)$.  A direct entrywise calculation gives
\[
 K_n=D(JH_n^{\mathsf T}J)D^{-1}.
\]
Indeed,
$\frac{(i-1)!}{(j-1)!}\binom{j-1}{i-1}=1/(j-i)!$ above the diagonal, and the
subdiagonal becomes $-(i-1)$.  Similarity, transposition, and reversal preserve
the determinant, proving the second equality.
\end{proof}

\section{The diamond convolution}
For sequences $x=(x_1,x_2,\ldots)$ and $y=(y_1,y_2,\ldots)$ define
\begin{equation}\label{eq:diamond}
 (x\mathbin\diamondsuit y)_n=\sum_{j=1}^{n-1}\binom njx_jy_{n-j}.
\end{equation}
This operation is commutative and associative because it corresponds to
multiplication of the EGFs $X(t)$ and $Y(t)$.

\begin{theorem}\label{thm:diamond-bell}
If $x^{k\diamondsuit}$ denotes the $k$-fold diamond power, then
\begin{equation}\label{eq:diamond-bell}
 \BellP nk(x_1,\ldots,x_{n-k+1})=\frac{(x^{k\diamondsuit})_n}{k!}.
\end{equation}
\end{theorem}
\begin{proof}
The EGF of $x^{k\diamondsuit}$ is $X(t)^k$.  Compare with
\eqref{eq:partial-bell-egf}.
\end{proof}
For example,
\[
 x\diamondsuit x=(0,2x_1^2,6x_1x_2,8x_1x_3+6x_2^2,\ldots)
\]
and
$(x^{3\diamondsuit})_4=36x_1^2x_2$, whence
$\BellP43=6x_1^2x_2$.

\chapter{Applications of Bell polynomials}

\section{Symmetric functions and Newton identities}
Let $e_n$ be the elementary symmetric function and
$p_n=\sum_i u_i^n$ the power sum in an alphabet $u_1,u_2,\ldots$.
\begin{theorem}\label{thm:bell-symmetric-functions}
\begin{align}
 e_n&=\frac1{n!}\BellC n(p_1,-1!p_2,2!p_3,\ldots,(-1)^{n-1}(n-1)!p_n)
 \label{eq:elementary-via-bell}\\
 &=\frac{(-1)^n}{n!}\BellC n(-p_1,-1!p_2,-2!p_3,\ldots,-(n-1)!p_n),
 \notag\\
 p_n&=\frac{(-1)^{n-1}}{(n-1)!}
 \sum_{k=1}^{n}(-1)^{k-1}(k-1)!
 \BellP nk(1!e_1,2!e_2,\ldots)
 \label{eq:powersum-via-bell}\\
 &=(-1)^n n\sum_{k=1}^{n}\frac1k
 \OBellP nk(-e_1,-e_2,\ldots).
 \notag
\end{align}
\end{theorem}
\begin{proof}
The classical product identity
\[
 E(t)=\sum_{n\ge0}e_nt^n=\prod_i(1+u_it)
 =\exp\!\left(\sum_{r\ge1}(-1)^{r-1}p_r\frac{t^r}{r}\right)
\]
and \eqref{eq:complete-bell-egf} give \eqref{eq:elementary-via-bell}.
Taking logarithms,
$\sum_{r\ge1}(-1)^{r-1}p_rt^r/r=\log(1+\sum_{j\ge1}e_jt^j)$, and expanding
the logarithm through ordinary Bell polynomials yields the second pair of forms.
The exponential form follows from \eqref{eq:ordinary-exponential-scaling}.
\end{proof}

\begin{corollary}[Determinant from traces]\label{cor:det-traces-bell}
For an $n\times n$ matrix $A$, set
$s_k=-(k-1)!\tr(A^k)$.  Then
\begin{equation}\label{eq:det-traces-bell}
 \det A=\frac{(-1)^n}{n!}\BellC n(s_1,s_2,\ldots,s_n).
\end{equation}
\end{corollary}
\begin{proof}
Apply \eqref{eq:elementary-via-bell} to the eigenvalues of $A$, using
$e_n=\det A$ and $p_k=\tr(A^k)$.  Both sides are polynomial in the matrix entries,
so the argument remains valid for nondiagonalizable matrices.
\end{proof}

\section{Cycle indices}
Let $a_j$ mark cycles of length $j$ in a permutation.
\begin{theorem}\label{thm:cycle-index-bell}
The cycle index of the symmetric group is
\begin{equation}\label{eq:cycle-index-bell}
 Z(S_n)=\frac1{n!}\BellC n(0!a_1,1!a_2,\ldots,(n-1)!a_n).
\end{equation}
\end{theorem}
\begin{proof}
A permutation is a set of cycles.  A labelled cycle of size $j$ has $(j-1)!$
linear realizations, so its component EGF is $(j-1)!a_jt^j/j!=a_jt^j/j$.
The exponential formula gives
$\exp(\sum_{j\ge1}a_jt^j/j)$; comparison with
\eqref{eq:complete-bell-egf} proves the identity after division by $n!$.
\end{proof}

\section{Moments and cumulants}
Let $M(t)=\E e^{tX}=\sum_{n\ge0}\mu'_nt^n/n!$ and
$K(t)=\log M(t)=\sum_{n\ge1}\kappa_nt^n/n!$.
\begin{theorem}[Moment--cumulant relations]\label{thm:moment-cumulant}
\begin{align}
 \mu'_n&=\BellC n(\kappa_1,\ldots,\kappa_n)
 =\sum_{k=1}^{n}\BellP nk(\kappa_1,\ldots),
 \label{eq:moments-from-cumulants}\\
 \kappa_n&=\sum_{k=1}^{n}(-1)^{k-1}(k-1)!
 \BellP nk(\mu'_1,\ldots,\mu'_{n-k+1}).
 \label{eq:cumulants-from-moments}
\end{align}
\end{theorem}
\begin{proof}
The identity $M=e^K$ is \eqref{eq:complete-bell-egf}; its logarithmic inverse is
\cref{thm:bell-transform-inverse}.
\end{proof}

\section{Hermite polynomials}
The probabilists' Hermite polynomials are defined by
\[
 e^{xt-t^2/2}=\sum_{n\ge0}\operatorname{He}_n(x)\frac{t^n}{n!}.
\]
Comparison with \eqref{eq:complete-bell-egf} gives and proves
\begin{equation}\label{eq:hermite-bell}
 \operatorname{He}_n(x)=\BellC n(x,-1,0,\ldots,0).
\end{equation}
In particular, the usual recurrence and differential equation follow by
differentiating this generating function.

\section{Polynomial sequences of binomial type}
Let
$h(t)=\sum_{j\ge1}a_jt^j/j!$ with $a_1\ne0$, and define
\begin{equation}\label{eq:binomial-type-bell}
 p_n(x)=\BellC n(a_1x,a_2x,\ldots,a_nx)
 =\sum_{k=0}^{n}\BellP nk(a_1,\ldots,a_{n-k+1})x^k.
\end{equation}

\begin{theorem}\label{thm:binomial-type-bell}
The polynomials $p_n$ satisfy
\begin{align}
 \sum_{n\ge0}p_n(x)\frac{t^n}{n!}&=e^{xh(t)},
 \label{eq:binomial-type-egf}\\
 p_n(x+y)&=\sum_{k=0}^{n}\binom nkp_k(x)p_{n-k}(y).
 \label{eq:binomial-type-identity}
\end{align}
If $Q=h^{-1}(D_x)$ is the associated delta operator, then
\begin{equation}\label{eq:delta-operator}
 Qp_n(x)=np_{n-1}(x).
\end{equation}
\end{theorem}
\begin{proof}
Substitute $x a_j$ in \eqref{eq:complete-bell-egf}.  The product
$e^{(x+y)h}=e^{xh}e^{yh}$ gives the binomial identity.  Finally
$h^{-1}(D_x)e^{xh(t)}=h^{-1}(h(t))e^{xh(t)}=t e^{xh(t)}$; coefficient comparison
gives \eqref{eq:delta-operator}.
\end{proof}

\begin{auditbox}
In some sources the symbol $B_n(x)$ denotes the Touchard polynomial, in others a
complete Bell polynomial, a Bernoulli polynomial, or a binary-tree sequence.  The
notations $T_n$, $\BellC n$, $\beta_n(x)$, and $C_n$ are kept separate here; this
is essential when the objects occur in the same inversion formula.
\end{auditbox}

\part{Eulerian Numbers and Descent Calculus}

\chapter{Ordinary Eulerian numbers}
\index{Eulerian numbers}

\section{Descents, symmetry, and the basic recurrence}
For a permutation $\pi=\pi_1\cdots\pi_n$, a descent is an index
$i\in\{1,\ldots,n-1\}$ with $\pi_i>\pi_{i+1}$.
\begin{definition}
The Eulerian number $\eulerian nk=A(n,k)$ counts permutations of $[n]$ with
exactly $k$ descents.  Put
\[
 A_0(t)=1,
 \qquad
 A_n(t)=\sum_{k=0}^{n-1}\eulerian nk t^k\quad(n\ge1).
\]
Outside $0\le k<n$ the number is zero.
\end{definition}

\begin{theorem}[Eulerian recurrence]\label{thm:eulerian-recurrence}
\begin{equation}\label{eq:eulerian-recurrence}
 \eulerian nk=(k+1)\eulerian{n-1}{k}
 +(n-k)\eulerian{n-1}{k-1}.
\end{equation}
Equivalently,
\begin{equation}\label{eq:eulerian-poly-recurrence}
 A_n(t)=\bigl(1+(n-1)t\bigr)A_{n-1}(t)
       +t(1-t)A_{n-1}'(t).
\end{equation}
\end{theorem}
\begin{proof}
Insert $n$ into a permutation of $[n-1]$.  If the old permutation has $k$
descents, insertion after one of those $k$ descents or at the end preserves the
number of descents, giving $k+1$ choices.  If it has $k-1$ descents, insertion in
any of the remaining $n-k$ slots creates one new descent.  This proves
\eqref{eq:eulerian-recurrence}.  Multiply by $t^k$, sum over $k$, and collect the
terms containing $k$ to obtain \eqref{eq:eulerian-poly-recurrence}.
\end{proof}

Complementation $\pi_i\mapsto n+1-\pi_i$ interchanges descents and ascents.
Therefore
\begin{equation}\label{eq:eulerian-symmetry}
 \eulerian nk=\eulerian n{n-1-k},
 \qquad
 A_n(t)=t^{n-1}A_n(t^{-1}).
\end{equation}
In particular $\eulerian n0=\eulerian n{n-1}=1$ and
$\sum_k\eulerian nk=n!$.  The next diagonal is
\begin{equation}\label{eq:eulerian-k1}
 \eulerian n1=2^n-(n+1),
\end{equation}
which follows either from the recurrence or from the explicit formula below.

\section{Exponential generating function and explicit formula}
\begin{theorem}[Eulerian EGF]\label{thm:eulerian-egf}
\begin{equation}\label{eq:eulerian-egf}
 \sum_{n\ge0}A_n(t)\frac{x^n}{n!}
 =\frac{t-1}{t-e^{(t-1)x}}
 =\frac{(1-t)e^{(1-t)x}}{1-te^{(1-t)x}}.
\end{equation}
\end{theorem}
\begin{proof}
Let $F(x,t)$ denote the left side.  Multiplying
\eqref{eq:eulerian-poly-recurrence} by $x^{n-1}/(n-1)!$ and summing gives the
first-order equation
\[
 (1-tx)F_x=t(1-t)F_t+F,
 \qquad F(0,t)=1.
\]
Direct substitution verifies that the displayed rational-exponential function
solves this equation and initial condition.  Formal uniqueness follows by
recursively comparing coefficients of $x$.
\end{proof}

Expanding the denominator geometrically gives the central power-series identity.
\begin{theorem}\label{thm:eulerian-power-series}
For $n\ge0$,
\begin{equation}\label{eq:eulerian-power-series}
 \sum_{m=0}^{\infty}(m+1)^n t^m
 =\frac{A_n(t)}{(1-t)^{n+1}}.
\end{equation}
Consequently,
\begin{equation}\label{eq:eulerian-explicit}
 \eulerian nk=\sum_{j=0}^{k}(-1)^j\binom{n+1}{j}(k+1-j)^n.
\end{equation}
\end{theorem}
\begin{proof}
From \eqref{eq:eulerian-egf},
\[
 F(x,t)=(1-t)\sum_{m\ge0}t^m e^{(m+1)(1-t)x}.
\]
Extracting $x^n/n!$ proves \eqref{eq:eulerian-power-series}.  Multiply by
$(1-t)^{n+1}$ and extract $t^k$ to obtain
\eqref{eq:eulerian-explicit}.  Formula \eqref{eq:eulerian-k1} is its case $k=1$.
\end{proof}

\begin{theorem}[A second polynomial recurrence]\label{thm:eulerian-binomial-recurrence}
For $n>1$,
\begin{equation}\label{eq:eulerian-binomial-recurrence}
 A_n(t)=\sum_{k=0}^{n-1}\binom nk A_k(t)(t-1)^{n-1-k}.
\end{equation}
\end{theorem}
\begin{proof}
The EGF of the right sides, summed over $n\ge1$, is
\[
 F(x,t)\frac{e^{(t-1)x}-1}{t-1}.
\]
Using \eqref{eq:eulerian-egf}, this equals $F(x,t)-1$, which is the EGF of the
left sides.
\end{proof}

\section{Worpitzky's identity and sums of powers}
\begin{theorem}[Worpitzky]\label{thm:worpitzky}
For every indeterminate $x$ and $n\ge1$,
\begin{equation}\label{eq:worpitzky}
 x^n=\sum_{k=0}^{n-1}\eulerian nk\binom{x+k}{n}.
\end{equation}
\end{theorem}
\begin{proof}
The coefficient of $t^{x-1}$ in
$t^{-k}(1-t)^{-n-1}$ is $\binom{x+k}{n}$ for positive integer $x$.
Equivalently, taking the coefficient of $t^{x-1}$ in
\eqref{eq:eulerian-power-series} after multiplying by the numerator gives
\eqref{eq:worpitzky}.  Since both sides are polynomials of degree $n$ in $x$ and
agree for all positive integers, they agree identically.

A direct combinatorial proof standardizes a word in $[x]^n$: its stable sorting
permutation records the relative order, and its $k$ descents specify $k$ mandatory
separations among the $n$ sorted positions.  Stars and bars gives
$\binom{x+k}{n}$ words for each permutation with $k$ descents.
\end{proof}

\begin{corollary}[Power sums]\label{cor:eulerian-power-sum}
For $m,n\ge1$,
\begin{equation}\label{eq:eulerian-power-sum}
 \sum_{r=1}^{m}r^n
 =\sum_{k=0}^{n-1}\eulerian nk\binom{m+k+1}{n+1}.
\end{equation}
\end{corollary}
\begin{proof}
Sum \eqref{eq:worpitzky} over $x=1,\ldots,m$ and use the hockey-stick identity
$\sum_{x=1}^{m}\binom{x+k}{n}=\binom{m+k+1}{n+1}$; the omitted lower boundary
terms vanish for $0\le k<n$.
\end{proof}

\section{Negative polylogarithms}
For $|z|<1$, $\Li_{-n}(z)=\sum_{m\ge1}m^nz^m$.  Shifting $m$ in
\eqref{eq:eulerian-power-series} gives
\begin{equation}\label{eq:negative-polylog-eulerian}
 \Li_{-n}(z)=\frac{zA_n(z)}{(1-z)^{n+1}}
 =\frac{1}{(1-z)^{n+1}}
   \sum_{k=0}^{n-1}\eulerian nk z^{n-k},
\end{equation}
where the second equality uses symmetry \eqref{eq:eulerian-symmetry}.  Since the
right side is rational, it also supplies the meromorphic continuation of the
negative polylogarithm.

\section{Eulerian and Stirling bases}
\begin{theorem}\label{thm:eulerian-stirling}
\begin{equation}\label{eq:eulerian-stirling}
 A_n(t)=\sum_{k=0}^{n}\stirtwo nk k!(t-1)^{n-k}.
\end{equation}
\end{theorem}
\begin{proof}
Use $m^n=\sum_k\stirtwo nk\falling mk$ in
$\sum_{m\ge0}m^nt^m$.  Since
$\sum_{m\ge0}\falling mk t^m=k!t^k/(1-t)^{k+1}$, multiplication by
$(1-t)^{n+1}/t$ yields
$A_n(t)=\sum_k\stirtwo nk k!t^{k-1}(1-t)^{n-k}$.  Replacing $t$ by $1/t$ and
using \eqref{eq:eulerian-symmetry}, or simply expanding around $t=1$, gives the
equivalent displayed form.
\end{proof}

\section{An Irwin--Hall integral and hypersimplex volumes}
\begin{theorem}\label{thm:eulerian-irwin-hall}
For every function $f:\R\to\C$ for which the following finite sum is defined,
\begin{equation}\label{eq:eulerian-irwin-hall}
 \sum_{k=0}^{n-1}\eulerian nk f(k)
 =n!\int_{[0,1]^n}f\!\left(\left\lfloor x_1+\cdots+x_n\right\rfloor\right)
 \dd x_1\cdots\dd x_n.
\end{equation}
In particular, the slab
\[
 \Delta_{n,k}=\{x\in[0,1]^n:k\le x_1+\cdots+x_n<k+1\}
\]
has volume $\eulerian nk/n!$.
\end{theorem}
\begin{proof}
The volume of $\{x\in[0,1]^n:\sum x_i<y\}$ is obtained from the simplex
$\{x_i\ge0:\sum x_i<y\}$ by inclusion--exclusion over coordinates exceeding $1$:
\[
 V_n(y)=\frac1{n!}\sum_{j=0}^{\lfloor y\rfloor}(-1)^j\binom nj(y-j)^n.
\]
Thus the volume of the $k$th slab is $V_n(k+1)-V_n(k)$.  Combining adjacent terms
by Pascal's identity gives
\[
 \frac1{n!}\sum_{j=0}^{k}(-1)^j\binom{n+1}{j}(k+1-j)^n
 =\frac1{n!}\eulerian nk
\]
by \eqref{eq:eulerian-explicit}.  Summing $f(k)$ times these volumes proves the
integral identity.
\end{proof}

\section{Alternating sums and Bernoulli numbers}
\begin{theorem}\label{thm:eulerian-alternating}
For $n>0$,
\begin{equation}\label{eq:eulerian-alternating}
 \sum_{k=0}^{n-1}(-1)^k\eulerian nk
 =2^{n+1}(2^{n+1}-1)\frac{\beta_{n+1}}{n+1}.
\end{equation}
For $n>1$,
\begin{align}
 \sum_{k=0}^{n-1}\frac{(-1)^k\eulerian nk}{\binom{n-1}{k}}&=0,
 \label{eq:eulerian-reciprocal-zero}\\
 \sum_{k=0}^{n-1}\frac{(-1)^k\eulerian nk}{\binom nk}&=(n+1)\beta_n.
 \label{eq:eulerian-reciprocal-bernoulli}
\end{align}
\end{theorem}
\begin{proof}
At $t=-1$, \eqref{eq:eulerian-egf} becomes
$2/(1+e^{-2x})=1+\tanh x$.  The Bernoulli generating function gives
\[
 \tanh x=\sum_{n\ge1}
 2^{n+1}(2^{n+1}-1)\frac{\beta_{n+1}}{n+1}\frac{x^n}{n!},
\]
which proves \eqref{eq:eulerian-alternating}.

For the remaining identities use
\[
 \binom{n-1}{k}^{-1}=n\int_0^1u^k(1-u)^{n-k-1}\dd u,
 \quad
 \binom nk^{-1}=(n+1)\int_0^1u^k(1-u)^{n-k}\dd u.
\]
Substitute $t=-u/(1-u)$ in \eqref{eq:eulerian-stirling}; then
\begin{equation}\label{eq:eulerian-beta-transform}
 (1-u)^nA_n\!\left(-\frac{u}{1-u}\right)
 =\sum_{j=0}^{n}(-1)^{n-j}\stirtwo njj!(1-u)^j.
\end{equation}
For the first integral divide by $1-u$ and integrate.  The result is
$n\sum_{j=1}^n(-1)^{n-j}\stirtwo nj(j-1)!=0$ for $n>1$, because it is the
coefficient of $x^n/n!$ in
$\log(1+(e^x-1))=x$.  For the second integral one obtains
$(n+1)\sum_j(-1)^{n-j}\stirtwo njj!/(j+1)$.  The standard Stirling formula
$\beta_n=\sum_j(-1)^jj!\stirtwo nj/(j+1)$ proves the claim; for odd $n>1$ both
sides vanish, and for even $n$ the signs coincide.
\end{proof}

\chapter{Geometric interpretations}

\section{The cube and its Ehrhart numerator}
For a lattice polytope $P$ of dimension $d$, its Ehrhart series has the form
$\sum_{m\ge0}|mP\cap\Z^d|t^m=h^*(t)/(1-t)^{d+1}$.
For the unit cube $C=[0,1]^n$ one has $|mC\cap\Z^n|=(m+1)^n$; hence
\eqref{eq:eulerian-power-series} proves:
\begin{theorem}
The $h^*$-polynomial of the standard $n$-cube is $A_n(t)$, so its $h^*$-vector is
the $n$th Eulerian row.
\end{theorem}

The regions $\Delta_{n,k}$ in \cref{thm:eulerian-irwin-hall} are half-open
hypersimplices.  Their normalized volumes are $n!\operatorname{vol}\Delta_{n,k}
=\eulerian nk$.

\section{The permutohedron}
The type-$A$ permutohedron is the convex hull of the permutations of
$(1,2,\ldots,n)$; its dimension is $n-1$.  Faces of its braid fan are indexed by
ordered set partitions.  There are $k!\stirtwo nk$ ordered partitions into $k$
blocks.

\begin{theorem}
The $h$-polynomial of the simple polytope dual to the type-$A$ permutohedron is
$A_n(t)$.
\end{theorem}
\begin{proof}
For a simple $(n-1)$-polytope the $h$-polynomial is the transform
$h(t)=\sum_F(t-1)^{\dim F}$ over faces of the dual simplicial complex.  Grouping
braid-fan faces by their $k$ ordered blocks gives
$\sum_k k!\stirtwo nk(t-1)^{n-k}$, which is $A_n(t)$ by
\eqref{eq:eulerian-stirling}.
\end{proof}

\chapter{Type-B and second-order Eulerian numbers}

\section{Signed permutations and the corrected type-B indexing}
A signed permutation is a word $\pi_1\cdots\pi_n$ whose absolute values form a
permutation of $[n]$.  Put $\pi_0=0$ and call
$i\in\{0,\ldots,n-1\}$ a type-$B$ descent when $\pi_i>\pi_{i+1}$.
Let $E_B(n,k)$ count signed permutations with $k$ such descents, and set
$M_n(t)=\sum_{k=0}^{n}E_B(n,k)t^k$.

\begin{auditbox}
The archived display reverses the two entries in the descent inequality and pairs
that convention with an incompatible explicit formula.  The definition above is
the standard Coxeter descent convention.  It gives the archived table
$1,6,1$ for $n=2$ and the midpoint-Eulerian generating function below.
\end{auditbox}

\begin{theorem}[Type-$B$ recurrence and Worpitzky series]\label{thm:typeB-eulerian}
\begin{align}
 E_B(n,k)&=(2k+1)E_B(n-1,k)+(2n-2k+1)E_B(n-1,k-1),
 \label{eq:typeB-recurrence}\\
 \sum_{m\ge0}(2m+1)^nt^m&=\frac{M_n(t)}{(1-t)^{n+1}},
 \label{eq:typeB-power-series}\\
 E_B(n,k)&=\sum_{j=0}^{k}(-1)^j\binom{n+1}{j}
             \bigl(2(k-j)+1\bigr)^n.
 \label{eq:typeB-explicit}
\end{align}
The first rows are
\[
1;\quad 1,1;\quad1,6,1;\quad1,23,23,1;\quad1,76,230,76,1.
\]
\end{theorem}
\begin{proof}
Insert either $n$ or $-n$ into a signed permutation of size $n-1$.  Marking the
slots according to whether insertion preserves or creates a descent gives
$2k+1$ preserving slots for an object with $k$ descents and
$2n-2k+1$ creating slots for one with $k-1$ descents; this proves the recurrence.

A type-$B$ standardization partitions the $(2m+1)^n$ words in the ordered alphabet
$-m<\cdots<-1<0<1<\cdots<m$ by their signed sorting permutation.  If that
permutation has $k$ type-$B$ descents, stars and bars gives
$\binom{m+n-k}{n}$ compatible weakly increasing absolute levels.  Hence
\[
 (2m+1)^n=\sum_{k=0}^{n}E_B(n,k)\binom{m+n-k}{n}.
\]
Multiply by $t^m$, sum in $m$, and use
$\sum_{m\ge0}\binom{m+n-k}{n}t^m=t^k/(1-t)^{n+1}$ to obtain
\eqref{eq:typeB-power-series}.  Multiplication by $(1-t)^{n+1}$ and coefficient
extraction gives \eqref{eq:typeB-explicit}.
\end{proof}

As in type $A$, these coefficients form the $h$-vector of the simple polytope dual
to the type-$B$ permutohedron; the proof repeats the braid-fan argument with signed
ordered partitions.

\section{Stirling permutations and second-order Eulerian numbers}
A \emph{Stirling permutation of order $n$} is a permutation of the multiset
$\{1,1,2,2,\ldots,n,n\}$ such that all entries between the two copies of $i$ are
larger than $i$.  Let $\eulertwo nk$ count those with $k$ ordinary descents and put
\[
 P_n(t)=\sum_{k=0}^{n-1}\eulertwo nk t^k,
 \qquad P_0(t)=1.
\]

\begin{theorem}[Second-order recurrence]\label{thm:second-eulerian-recurrence}
\begin{align}
 \eulertwo nk&=(2n-k-1)\eulertwo{n-1}{k-1}
 +(k+1)\eulertwo{n-1}{k},
 \label{eq:second-eulerian-recurrence}\\
 P_{n+1}(t)&=(2nt+1)P_n(t)-t(t-1)P_n'(t),
 \label{eq:second-eulerian-poly-recurrence}\\
 P_n(1)&=(2n-1)!!.
 \label{eq:second-eulerian-row-sum}
\end{align}
\end{theorem}
\begin{proof}
The two copies of the largest label $n$ must be adjacent.  Insert the pair $nn$
into one of the $2n-1$ gaps of a Stirling permutation of order $n-1$.  A descent
gap or the terminal gap preserves the number of descents, giving $k+1$ choices
from an object with $k$ descents.  Each of the remaining $2n-k-1$ gaps creates one
descent from an object with $k-1$ descents.  This is the first recurrence; summing
against $t^k$ gives the second.  At each stage there are $2n-1$ insertion gaps, so
the total count is $(2n-1)!!$.
\end{proof}

The differential recurrence has two useful equivalent forms:
\begin{align}
 (t-1)^{-2n-2}P_{n+1}(t)
 &=\left(t(1-t)^{-2n-1}P_n(t)\right)',
 \label{eq:second-eulerian-differential-form}\\
 u_{n+1}(t)&=\left(\frac{t}{1-t}u_n(t)\right)',
 \quad u_n(t)=(t-1)^{-2n}P_n(t),\quad u_0=1.
 \label{eq:second-eulerian-u}
\end{align}
They follow by the product rule from
\eqref{eq:second-eulerian-poly-recurrence}.

\begin{theorem}[Diagonal Stirling generating function]\label{thm:second-eulerian-stirling-gf}
For $n\ge1$,
\begin{equation}\label{eq:second-eulerian-stirling-gf}
 \sum_{m=0}^{\infty}\stirtwo{n+m}{m}t^m
 =\frac{tP_n(t)}{(1-t)^{2n+1}}.
\end{equation}
\end{theorem}
\begin{proof}
Let $F_n(t)$ denote the left side.  The second-kind recurrence gives
$(1-t)F_{n+1}=tF_n'$, with $F_0=(1-t)^{-1}$.  Hence
$F_1=t/(1-t)^3$.  If the displayed formula holds for $n$, differentiating
$tP_n/(1-t)^{2n+1}$ and using
\eqref{eq:second-eulerian-poly-recurrence} gives the formula for $n+1$.
\end{proof}


\section{Second-order Eulerian interpolation of first-kind diagonals}
The same second-order Eulerian row that controls
\eqref{eq:second-eulerian-stirling-gf} also gives a binomial-basis expansion for
the near-diagonals of the unsigned first-kind triangle.

\begin{theorem}[Second-order Eulerian diagonal formula]
\label{thm:second-eulerian-first-diagonal}
For $p,m\ge0$,
\begin{equation}
\label{eq:second-eulerian-first-diagonal}
 \stirone m{m-p}
 =\sum_{j=0}^{p}\eulertwo pj\binom{m+j}{2p},
\end{equation}
where an out-of-range first-kind number is zero.  Equivalently, the polynomial
continuation in a variable $x$ satisfies
\[
 \stirone{x}{x-p}
 =\sum_{j=0}^{p}\eulertwo pj\binom{x+j}{2p}.
\]
The first instances are
\begin{align*}
 \stirone{x}{x-1}&=\binom{x}{2},\\
 \stirone{x}{x-2}&=\binom{x}{4}+2\binom{x+1}{4},\\
 \stirone{x}{x-3}&=\binom{x}{6}+8\binom{x+1}{6}
                    +6\binom{x+2}{6}.
\end{align*}
\end{theorem}
\begin{proof}
For fixed $p$ define
\[
 D_p(t)=\sum_{m\ge0}\stirone m{m-p}t^m.
\]
Writing $a_{m,p}=\stirone m{m-p}$, the first-kind recurrence becomes
$a_{m,p}=a_{m-1,p}+(m-1)a_{m-1,p-1}$.  Therefore
\begin{equation}
\label{eq:first-diagonal-ogf-recurrence}
 (1-t)D_p(t)=t^2D_{p-1}'(t),
 \qquad D_0(t)=\frac1{1-t}.
\end{equation}
Let
\[
 Q_p(t)=t^{2p}P_p(t^{-1}).
\]
The polynomial recurrence \eqref{eq:second-eulerian-poly-recurrence} is equivalent,
after the substitution $t\mapsto t^{-1}$, to
\[
 Q_p(t)=t^2\bigl((1-t)Q_{p-1}'(t)+(2p-1)Q_{p-1}(t)\bigr).
\]
It follows inductively from \eqref{eq:first-diagonal-ogf-recurrence} that
\begin{equation}
\label{eq:first-diagonal-second-eulerian-ogf}
 D_p(t)=\frac{Q_p(t)}{(1-t)^{2p+1}}
 =\frac{t^{2p}P_p(t^{-1})}{(1-t)^{2p+1}}.
\end{equation}
Since $P_p(u)=\sum_{j=0}^{p}\eulertwo pj u^j$, coefficient extraction gives
\[
 [t^m]D_p(t)
 =\sum_{j=0}^{p}\eulertwo pj
   [t^{m-2p+j}](1-t)^{-2p-1}
 =\sum_{j=0}^{p}\eulertwo pj\binom{m+j}{2p},
\]
which is \eqref{eq:second-eulerian-first-diagonal}.  Both sides are polynomials in
$m$ of degree at most $2p$, so the same identity gives the stated polynomial
continuation.
\end{proof}

\part{Bernoulli, Euler, and Summation Calculus}

\chapter{Bernoulli numbers, Bernoulli polynomials, and N\"orlund extensions}
\index{Bernoulli number}\index{Bernoulli polynomial}

\section{Generating functions and first consequences}
The Bernoulli numbers have already appeared in the Eulerian alternating sums.
We now develop their own coefficient calculus.  Our convention is
$\beta_1=-1/2$, consistent with the generating-function normalization used in
the DLMF \cite{NISTDLMF}.

\begin{definition}
The Bernoulli polynomials $\beta_n(x)$ are defined by
\begin{equation}\label{eq:bernoulli-polynomial-egf}
 \frac{te^{xt}}{e^t-1}
 =\sum_{n\ge0}\beta_n(x)\frac{t^n}{n!},
 \qquad
 \beta_n:=\beta_n(0).
\end{equation}
\end{definition}
The first values are
\begin{align*}
 \beta_0(x)&=1,\\
 \beta_1(x)&=x-\frac12,\\
 \beta_2(x)&=x^2-x+\frac16,\\
 \beta_3(x)&=x^3-\frac32x^2+\frac12x,\\
 \beta_4(x)&=x^4-2x^3+x^2-\frac1{30},\\
 \beta_5(x)&=x^5-\frac52x^4+\frac53x^3-\frac16x.
\end{align*}

\begin{theorem}[Appell, addition, difference, and reflection laws]
\label{thm:bernoulli-basic-laws}
For $n\ge0$,
\begin{align}
 \beta_n'(x)&=n\beta_{n-1}(x),
 \label{eq:bernoulli-appell}\\
 \beta_n(x+y)&=\sum_{k=0}^{n}\binom nk\beta_k(x)y^{n-k},
 \label{eq:bernoulli-addition}\\
 \beta_n(x+1)-\beta_n(x)&=nx^{n-1}\quad(n\ge1),
 \label{eq:bernoulli-difference}\\
 \beta_n(1-x)&=(-1)^n\beta_n(x).
 \label{eq:bernoulli-reflection}
\end{align}
In particular,
\begin{equation}\label{eq:bernoulli-polynomial-explicit}
 \beta_n(x)=\sum_{k=0}^{n}\binom nk\beta_kx^{n-k}.
\end{equation}
Moreover, $\beta_{2m+1}=0$ for every $m\ge1$.
\end{theorem}
\begin{proof}
Differentiate \eqref{eq:bernoulli-polynomial-egf} with respect to $x$;
coefficient comparison gives \eqref{eq:bernoulli-appell}.  Multiplication of
the generating function by $e^{yt}$ gives
\eqref{eq:bernoulli-addition}; setting $x=0$ gives
\eqref{eq:bernoulli-polynomial-explicit}.  Next,
\[
 \frac{te^{(x+1)t}}{e^t-1}-\frac{te^{xt}}{e^t-1}=te^{xt},
\]
whose coefficient of $t^n/n!$ is $nx^{n-1}$ for $n\ge1$.  Finally,
\[
 \frac{te^{(1-x)t}}{e^t-1}
 =\frac{(-t)e^{x(-t)}}{e^{-t}-1},
\]
so replacing $t$ by $-t$ proves \eqref{eq:bernoulli-reflection}.  At $x=0$
this says $\beta_n(1)=(-1)^n\beta_n$.  For $n\ne1$ the difference law gives
$\beta_n(1)=\beta_n$, hence odd $\beta_n$ vanish for $n>1$.
\end{proof}

\section{Faulhaber summation and the multiplication theorem}
\begin{theorem}[Faulhaber's formula in translated form]
\label{thm:faulhaber-bernoulli}
For $p\ge0$, an integer $N\ge0$, and arbitrary $x$,
\begin{equation}\label{eq:faulhaber-bernoulli}
 \sum_{r=0}^{N-1}(x+r)^p
 =\frac{\beta_{p+1}(x+N)-\beta_{p+1}(x)}{p+1}.
\end{equation}
\end{theorem}
\begin{proof}
By \eqref{eq:bernoulli-difference},
\[
 \frac{\beta_{p+1}(u+1)-\beta_{p+1}(u)}{p+1}=u^p.
\]
Set $u=x+r$ and sum from $r=0$ to $N-1$; the left side telescopes.
\end{proof}
Thus the Bernoulli basis is adapted to discrete antidifferentiation just as the
falling-factorial basis is adapted to the forward-difference operator.

\begin{theorem}[Raabe multiplication formula]
\label{thm:bernoulli-multiplication}
For every positive integer $q$,
\begin{equation}\label{eq:bernoulli-multiplication}
 \sum_{r=0}^{q-1}\beta_n\!\left(x+\frac rq\right)
 =q^{1-n}\beta_n(qx).
\end{equation}
\end{theorem}
\begin{proof}
Sum the generating functions and use the finite geometric series:
\begin{align*}
 \sum_{r=0}^{q-1}\frac{te^{(x+r/q)t}}{e^t-1}
 &=\frac{te^{xt}}{e^t-1}\frac{e^t-1}{e^{t/q}-1}\\
 &=q\frac{(t/q)e^{(qx)(t/q)}}{e^{t/q}-1}.
\end{align*}
The coefficient of $t^n/n!$ on the right is $q^{1-n}\beta_n(qx)$.
\end{proof}

\section{A Stirling-number formula}
\begin{theorem}\label{thm:bernoulli-stirling}
For $n\ge0$,
\begin{equation}\label{eq:bernoulli-stirling}
 \boxed{\displaystyle
 \beta_n=\sum_{k=0}^{n}(-1)^k\frac{k!}{k+1}\stirtwo nk.}
\end{equation}
\end{theorem}
\begin{proof}
The formal identity
\[
 \sum_{k\ge0}\frac{(-1)^ku^k}{k+1}=\frac{\log(1+u)}u
\]
with $u=e^t-1$ becomes
\[
 \frac{t}{e^t-1}
 =\sum_{k\ge0}\frac{(-1)^k}{k+1}(e^t-1)^k.
\]
Now use
$(e^t-1)^k=k!\sum_{n\ge k}\stirtwo nk t^n/n!$ and compare
coefficients.  Only $k\le n$ contributes, so every coefficient calculation is
finite.
\end{proof}

This formula is also the integral identity
\[
 \beta_n=\int_0^1 T_n(-u)\,\dd u,
\]
because integration of $u^k$ supplies the factor $1/(k+1)$.  It places
Bernoulli numbers, Touchard polynomials, and Stirling numbers in one triangular
transform.

\section{Generalized Bernoulli polynomials}
\index{Norlund polynomial@N\"orlund polynomial}
\begin{definition}
For a scalar parameter $\alpha$, define the generalized Bernoulli, or
N\"orlund, polynomials by
\begin{equation}\label{eq:norlund-egf}
 \left(\frac{t}{e^t-1}\right)^{\!\alpha}e^{xt}
 =\sum_{n\ge0}\beta_n^{(\alpha)}(x)\frac{t^n}{n!}.
\end{equation}
Thus $\beta_n^{(1)}(x)=\beta_n(x)$ and
$\beta_n^{(0)}(x)=x^n$.
\end{definition}

\begin{theorem}[Convolution calculus]
\label{thm:norlund-convolution}
For arbitrary parameters $\alpha,\gamma$,
\begin{align}
 \frac{\partial}{\partial x}\beta_n^{(\alpha)}(x)
 &=n\beta_{n-1}^{(\alpha)}(x),
 \label{eq:norlund-appell}\\
 \beta_n^{(\alpha+\gamma)}(x+y)
 &=\sum_{k=0}^{n}\binom nk
   \beta_k^{(\alpha)}(x)\beta_{n-k}^{(\gamma)}(y).
 \label{eq:norlund-convolution}
\end{align}
Furthermore,
\begin{equation}\label{eq:norlund-diagonal}
 \beta_{n-1}^{(n)}=(-1)^{n-1}(n-1)!\qquad(n\ge1).
\end{equation}
\end{theorem}
\begin{proof}
The derivative and convolution laws follow respectively by differentiating and
multiplying the generating functions.  For the diagonal value, Lagrange
inversion applied to the inverse pair $e^t-1$ and $\log(1+z)$ gives
\[
 [t^{n-1}]\left(\frac{t}{e^t-1}\right)^n=(-1)^{n-1}.
\]
By \eqref{eq:norlund-egf}, the left side is
$\beta_{n-1}^{(n)}/(n-1)!$.
\end{proof}

The complete Bell polynomials give an explicit construction.  From
\begin{equation}\label{eq:log-bernoulli-generator}
 \log\frac{t}{e^t-1}
 =-\frac t2-\sum_{r\ge1}
   \frac{\beta_{2r}}{2r(2r)!}t^{2r},
\end{equation}
we obtain
\begin{equation}\label{eq:norlund-complete-bell}
 \beta_n^{(\alpha)}(x)
 =\BellC n(c_1,c_2,\ldots,c_n),
\end{equation}
where
\[
 c_1=x-\frac\alpha2,\qquad
 c_{2r}=-\frac{\alpha\beta_{2r}}{2r},\qquad
 c_{2r+1}=0\quad(r\ge1).
\]
Indeed, the logarithm of the generating function in
\eqref{eq:norlund-egf} is $\sum_{j\ge1}c_jt^j/j!$, and the complete Bell
polynomial is the exponential of exactly such a series.  Formula
\eqref{eq:log-bernoulli-generator} follows by differentiating both sides,
using the generating function for $\beta_n$, and matching the value at $t=0$.

\section{Fourier expansion and the even zeta values}
Let $\widetilde\beta_n(x)=\beta_n(\{x\})$ be the periodic Bernoulli function,
where $\{x\}$ is the fractional part.

\begin{theorem}[Fourier series of periodic Bernoulli polynomials]
\label{thm:bernoulli-fourier}
For $n\ge2$ and real $x$,
\begin{equation}\label{eq:bernoulli-fourier}
 \widetilde\beta_n(x)
 =-\frac{n!}{(2\pi i)^n}
   \sum_{k\in\Z\setminus\{0\}}\frac{e^{2\pi ikx}}{k^n}.
\end{equation}
The series is absolutely and uniformly convergent.  Consequently, for
$m\ge1$,
\begin{equation}\label{eq:zeta-even-bernoulli}
 \zeta(2m)=(-1)^{m+1}
 \frac{(2\pi)^{2m}}{2(2m)!}\beta_{2m}.
\end{equation}
\end{theorem}
\begin{proof}
For $k\ne0$, let
$c_k=\int_0^1\beta_n(x)e^{-2\pi ikx}\,\dd x$.  Since
$\beta_r(1)=\beta_r(0)$ for $r\ne1$, integration by parts and
\eqref{eq:bernoulli-appell} give
\[
 c_k=\frac{n}{2\pi ik}c_k^{(n-1)}
     =\cdots
     =\frac{n!}{(2\pi ik)^{n-1}}
       \int_0^1\beta_1(x)e^{-2\pi ikx}\,\dd x.
\]
One final integration by parts uses
$\beta_1(1)-\beta_1(0)=1$ and yields
$c_k=-n!/(2\pi ik)^n$.  The mean coefficient is zero because
$\int_0^1\beta_n(x)\dd x=0$ for $n\ge1$, by
\eqref{eq:bernoulli-appell}.  Since $\sum_{k\ne0}|k|^{-n}$ converges for
$n\ge2$, the Fourier series converges absolutely and uniformly to the
continuous periodic function $\widetilde\beta_n$.  Setting $x=0$ and
$n=2m$ gives
\[
 \beta_{2m}=-\frac{(2m)!}{(2\pi i)^{2m}}2\zeta(2m),
\]
which rearranges to \eqref{eq:zeta-even-bernoulli}.
\end{proof}

\section{A modified even Bernoulli sequence}
One of the supplied MathWorld notes records the coefficients
$\mathfrak b_n$ defined by
\begin{equation}\label{eq:modified-bernoulli-egf}
 \sum_{n\ge0}\mathfrak b_nx^n
 =\frac12\log\!\left(
   \frac{e^{x/2}-e^{-x/2}}{x/2}\right).
\end{equation}
The constant term is $\mathfrak b_0=\frac12\log2$, all positive odd
coefficients vanish, and for even $n\ge2$,
\begin{equation}\label{eq:modified-bernoulli-closed}
 \mathfrak b_n=\frac{\beta_n}{2n\,n!}.
\end{equation}
To prove this, differentiate \eqref{eq:modified-bernoulli-egf} and use
\[
 \coth u=\frac1u+\sum_{r\ge1}
 \frac{2^{2r}\beta_{2r}}{(2r)!}u^{2r-1}.
\]
The singular terms cancel and termwise integration, with the value at $x=0$,
gives \eqref{eq:modified-bernoulli-closed}.  Thus this ``modified'' sequence
is a simple ordinary-generating normalization of the even Bernoulli numbers.

\chapter{Euler polynomials and Genocchi numbers}
\index{Euler polynomial}\index{Genocchi number}

\section{Euler polynomials}
\begin{definition}
The Euler polynomials $E_n(x)$ are defined by
\begin{equation}\label{eq:euler-polynomial-egf}
 \frac{2e^{xt}}{e^t+1}
 =\sum_{n\ge0}E_n(x)\frac{t^n}{n!}.
\end{equation}
\end{definition}
The first few are
\begin{align*}
 E_0(x)&=1,\\
 E_1(x)&=x-\frac12,\\
 E_2(x)&=x^2-x,\\
 E_3(x)&=x^3-\frac32x^2+\frac14,\\
 E_4(x)&=x^4-2x^3+x.
\end{align*}
These Euler polynomials should not be confused with the Eulerian polynomials of
the preceding part.

\begin{theorem}[Basic Euler-polynomial laws]
\label{thm:euler-polynomial-laws}
For $n\ge0$,
\begin{align}
 E_n'(x)&=nE_{n-1}(x),
 \label{eq:euler-poly-appell}\\
 E_n(x+y)&=\sum_{k=0}^{n}\binom nkE_k(x)y^{n-k},
 \label{eq:euler-poly-addition}\\
 E_n(x+1)+E_n(x)&=2x^n,
 \label{eq:euler-poly-difference}\\
 E_n(1-x)&=(-1)^nE_n(x).
 \label{eq:euler-poly-reflection}
\end{align}
Consequently,
\begin{equation}\label{eq:alternating-power-sum-euler}
 2\sum_{r=0}^{N-1}(-1)^r(x+r)^n
 =E_n(x)+(-1)^{N-1}E_n(x+N).
\end{equation}
\end{theorem}
\begin{proof}
The derivative and addition laws follow from
\eqref{eq:euler-polynomial-egf} exactly as for any Appell sequence.  Adding
the generating functions at $x$ and $x+1$ gives $2e^{xt}$, proving
\eqref{eq:euler-poly-difference}.  Replacing $t$ by $-t$ gives the reflection
law.  Finally, apply \eqref{eq:euler-poly-difference} at $x+r$, multiply by
$(-1)^r$, and sum; all interior Euler-polynomial terms cancel.
\end{proof}

\begin{theorem}[Bernoulli--Euler connection]
\label{thm:bernoulli-euler-connection}
For $n\ge0$,
\begin{equation}\label{eq:bernoulli-euler-connection}
 E_n(x)=\frac{2^{n+1}}{n+1}
 \left[
 \beta_{n+1}\!\left(\frac{x+1}{2}\right)
 -\beta_{n+1}\!\left(\frac{x}{2}\right)
 \right].
\end{equation}
\end{theorem}
\begin{proof}
Multiply the right side by $t^n/n!$ and sum over $n$.  After putting
$m=n+1$, the result is
\begin{align*}
 \frac1t\left[
 \frac{2t e^{(x+1)t}}{e^{2t}-1}
 -\frac{2t e^{xt}}{e^{2t}-1}
 \right]
 &=\frac{2e^{xt}(e^t-1)}{e^{2t}-1}\\
 &=\frac{2e^{xt}}{e^t+1},
\end{align*}
which is the defining generating function.
\end{proof}

The centered values produce the classical signed secant numbers.  Define
$\mathcal E_n=2^nE_n(1/2)$.  Then
\begin{equation}\label{eq:secant-euler-egf}
 \sum_{n\ge0}\mathcal E_n\frac{z^n}{n!}=\operatorname{sech}z.
\end{equation}
Hence $\mathcal E_{2m+1}=0$ and
$\mathcal E_0,\mathcal E_2,\mathcal E_4,\ldots=1,-1,5,-61,1385,\ldots$.
This follows by setting $x=1/2$ in \eqref{eq:euler-polynomial-egf} and then
replacing $t$ by $2z$.

\section{Genocchi polynomials and numbers}
\begin{definition}
Define Genocchi polynomials by
\begin{equation}\label{eq:genocchi-egf}
 \frac{2te^{xt}}{e^t+1}
 =\sum_{n\ge0}G_n(x)\frac{t^n}{n!},
 \qquad G_n:=G_n(0).
\end{equation}
\end{definition}

\begin{theorem}\label{thm:genocchi-relations}
For $n\ge1$,
\begin{align}
 G_n(x)&=nE_{n-1}(x),
 \label{eq:genocchi-euler}\\
 G_n&=2(1-2^n)\beta_n,
 \label{eq:genocchi-bernoulli}\\
 G_n(x+1)+G_n(x)&=2n x^{n-1}.
 \label{eq:genocchi-difference}
\end{align}
Thus $G_0=0,G_1=1,G_2=-1,G_3=0,G_4=1,G_6=-3,G_8=17$.
\end{theorem}
\begin{proof}
Multiplying the Euler-polynomial generating function by $t$ proves
\eqref{eq:genocchi-euler}.  For \eqref{eq:genocchi-bernoulli}, use
\[
 \frac{2t}{e^t+1}
 =\frac{2t}{e^t-1}-\frac{4t}{e^{2t}-1}
 =2\sum_{n\ge0}\beta_n\frac{t^n}{n!}
  -2\sum_{n\ge0}\beta_n\frac{(2t)^n}{n!}.
\]
Coefficient comparison gives the result.  The difference identity follows by
adding the generating functions at $x$ and $x+1$.
\end{proof}

\chapter{Bernoulli polynomials of the second kind}
\index{Bernoulli polynomial!second kind}
These polynomials are also called Cauchy polynomials of the first kind.  Their
natural delta series is $\log(1+t)$ rather than $t$, so they belong to the
Sheffer rather than the Appell class.

\begin{definition}
Define $b_n(x)$ by
\begin{equation}\label{eq:bernoulli-second-kind-egf}
 \frac{t(1+t)^x}{\log(1+t)}
 =\sum_{n\ge0}b_n(x)\frac{t^n}{n!},
 \qquad b_n:=b_n(0).
\end{equation}
\end{definition}
The first polynomials are
\begin{align*}
 b_0(x)&=1,\\
 b_1(x)&=x+\frac12,\\
 b_2(x)&=x^2-\frac16,\\
 b_3(x)&=x^3-\frac32x^2+\frac14,\\
 b_4(x)&=x^4-4x^3+4x^2-\frac{19}{30}.
\end{align*}

\begin{theorem}[Integral and finite-difference representations]
\label{thm:bernoulli-second-kind-laws}
For $n\ge0$,
\begin{align}
 b_n(x)&=\int_0^1\falling{x+u}{n}\,\dd u,
 \label{eq:bernoulli-second-kind-integral}\\
 b_n(x+1)-b_n(x)&=n\falling{x}{n-1},
 \label{eq:bernoulli-second-kind-difference}\\
 b_n'(x)&=n\falling{x}{n-1},
 \label{eq:bernoulli-second-kind-derivative}\\
 b_n(x)&=b_n+n\sum_{k=1}^{n}\frac{s(n-1,k-1)}{k}x^k,
 \label{eq:bernoulli-second-kind-stirling}\\
 b_n\!\left(\frac n2-1-x\right)
 &=(-1)^n b_n\!\left(\frac n2-1+x\right).
 \label{eq:bernoulli-second-kind-reflection}
\end{align}
\end{theorem}
\begin{proof}
Since
$\sum_{n\ge0}\falling{x+u}{n}t^n/n!=(1+t)^{x+u}$,
integration over $0\le u\le1$ gives
\[
 \int_0^1(1+t)^{x+u}\dd u
 =(1+t)^x\frac{t}{\log(1+t)},
\]
proving \eqref{eq:bernoulli-second-kind-integral}.  Subtracting the generating
functions at $x+1$ and $x$ gives $t(1+t)^x$, which proves the difference law.
Differentiating with respect to $x$ cancels the logarithm and gives the same
right side, proving \eqref{eq:bernoulli-second-kind-derivative}.  Since
$\falling{x}{n-1}=\sum_{j=0}^{n-1}s(n-1,j)x^j$, integration from $0$ to $x$
gives \eqref{eq:bernoulli-second-kind-stirling}.  Finally use
\[
 \falling{z}{n}=(-1)^n\falling{n-1-z}{n}
\]
in the integral representation and change $u$ to $1-u$.
\end{proof}

The inverse relation between $\log(1+t)$ and $e^t-1$ gives the diagonal
coefficient
\begin{equation}\label{eq:cauchy-diagonal-preview}
 [t^{n-1}]\left(\frac{t}{\log(1+t)}\right)^n=\frac1{(n-1)!},
\end{equation}
which will reappear as a direct Lagrange-inversion example.  The polynomials
$b_n(x)$ will also provide a concrete Sheffer sequence in the final synthesis
part.

\chapter{Euler--Maclaurin summation and Darboux--Taylor remainders}
\index{Euler--Maclaurin formula}

\section{Periodic Bernoulli functions}
For $r\ge1$, put
$\widetilde\beta_r(x)=\beta_r(x-\lfloor x\rfloor)$ away from integers.  At
integers the value of $\widetilde\beta_1$ may be chosen arbitrarily in
integration formulas; the symmetric convention sets it to $0$.  For $r\ge2$
the function is continuous and $1$-periodic.  On every open unit interval,
\begin{equation}\label{eq:periodic-bernoulli-derivative}
 \frac{\dd}{\dd x}\widetilde\beta_r(x)
 =r\widetilde\beta_{r-1}(x).
\end{equation}

\begin{theorem}[Euler--Maclaurin with an exact remainder]
\label{thm:euler-maclaurin}
Let $a<b$ be integers, let $m\ge1$, and let
$f\in C^{2m}([a,b])$.  Then
\begin{align}
 \sum_{k=a}^{b}f(k)
 ={}&\int_a^bf(x)\,\dd x+\frac{f(a)+f(b)}2
 \notag\\
 &+\sum_{r=1}^{m}\frac{\beta_{2r}}{(2r)!}
   \left(f^{(2r-1)}(b)-f^{(2r-1)}(a)\right)+R_m,
 \label{eq:euler-maclaurin}\\
 R_m={}&-\frac1{(2m)!}\int_a^b
 \widetilde\beta_{2m}(x)f^{(2m)}(x)\,\dd x.
 \label{eq:euler-maclaurin-remainder}
\end{align}
Moreover,
\begin{equation}\label{eq:euler-maclaurin-bound}
 |R_m|\le
 \frac{2\zeta(2m)}{(2\pi)^{2m}}
 \int_a^b|f^{(2m)}(x)|\,\dd x.
\end{equation}
\end{theorem}
\begin{proof}
On one interval $[k,k+1]$, integration by parts gives
\[
 \int_k^{k+1}\widetilde\beta_1(x)f'(x)\,\dd x
 =\frac{f(k)+f(k+1)}2-\int_k^{k+1}f(x)\,\dd x.
\]
Summing from $k=a$ to $b-1$ and then adding
$(f(a)+f(b))/2$ gives
\begin{equation}\label{eq:euler-maclaurin-start}
 \sum_{k=a}^{b}f(k)
 =\int_a^bf(x)\,\dd x+\frac{f(a)+f(b)}2
  +\int_a^b\widetilde\beta_1(x)f'(x)\,\dd x.
\end{equation}
Use \eqref{eq:periodic-bernoulli-derivative} and integrate by parts
successively.  Odd Bernoulli numbers beyond $\beta_1$ vanish, so the boundary
terms occur only at even indices and are exactly the displayed sum; stopping
after $2m$ derivatives leaves \eqref{eq:euler-maclaurin-remainder}.  Finally,
\cref{thm:bernoulli-fourier} gives
\[
 \sup_x|\widetilde\beta_{2m}(x)|
 \le\frac{2(2m)!}{(2\pi)^{2m}}\zeta(2m),
\]
and insertion in the remainder integral proves
\eqref{eq:euler-maclaurin-bound}.
\end{proof}

\section{Power sums, harmonic numbers, and factorials}
For a polynomial $f$, the Euler--Maclaurin process terminates and is equivalent
to Faulhaber's formula.  For nonpolynomial $f$ it produces asymptotic
expansions with computable remainders.

\begin{corollary}\label{cor:harmonic-euler-maclaurin}
For every fixed $m\ge1$,
\begin{equation}\label{eq:harmonic-asymptotic}
 H_N=\log N+\gamma+\frac1{2N}
 -\sum_{r=1}^{m-1}\frac{\beta_{2r}}{2r\,N^{2r}}
 +O(N^{-2m})
 \qquad(N\to\infty).
\end{equation}
\end{corollary}
\begin{proof}
Apply \cref{thm:euler-maclaurin} to $f(x)=1/x$ on $[N,M]$, then let
$M\to\infty$.  Every differentiated term at $M$ vanishes, and the remainder
bound is $O(N^{-2m})$.  The constant left after subtracting $\log N$ is, by
definition, Euler's constant $\gamma$.
\end{proof}

\begin{corollary}[Stirling expansion]\label{cor:stirling-euler-maclaurin}
For every fixed $m\ge1$,
\begin{multline}\label{eq:stirling-euler-maclaurin}
 \log N!=\left(N+\frac12\right)\log N-N+\frac12\log(2\pi)\\
 {}+\sum_{r=1}^{m-1}
 \frac{\beta_{2r}}{2r(2r-1)N^{2r-1}}
 +O(N^{-(2m-1)}).
\end{multline}
\end{corollary}
\begin{proof}
Apply Euler--Maclaurin to $f(x)=\log x$ and absorb the lower-end terms into a
constant.  The derivatives
$f^{(j)}(x)=(-1)^{j-1}(j-1)!x^{-j}$ give the displayed inverse powers and the
remainder estimate.  The constant is $\frac12\log(2\pi)$, as follows from
Wallis' product (equivalently, from the central-binomial-coefficient estimate
obtained from that product).  This is the usual proof of Stirling's formula
with all higher corrections retained.
\end{proof}

\section{Darboux's repeated-integration formula}
The Darboux item in the supplementary material is the integration-by-parts
mechanism behind Taylor expansion.  In a normalization convenient here it is
the following exact identity.

\begin{theorem}[Taylor formula with integral remainder]
\label{thm:darboux-taylor}
If $f\in C^m$ on the segment joining $a$ and $x$, then
\begin{equation}\label{eq:darboux-taylor}
 f(x)=\sum_{r=0}^{m-1}\frac{f^{(r)}(a)}{r!}(x-a)^r
 +\frac1{(m-1)!}\int_a^x(x-t)^{m-1}f^{(m)}(t)\,\dd t.
\end{equation}
\end{theorem}
\begin{proof}
For $m=1$ this is the fundamental theorem of calculus.  Integrate the
remainder by parts, taking
$u=f^{(m)}(t)$ and
$\dd v=(x-t)^{m-1}\dd t/(m-1)!$.  The boundary at $t=x$ vanishes and the
boundary at $t=a$ contributes
$f^{(m-1)}(a)(x-a)^{m-1}/(m-1)!$; the remaining integral is the remainder of
order $m-1$.  Induction gives the formula.
\end{proof}

Both \cref{thm:darboux-taylor,thm:euler-maclaurin} are repeated-integration
identities.  Taylor's kernel is an ordinary power, whereas Euler--Maclaurin's
kernel is a periodic Bernoulli polynomial adapted to summation over a lattice.


\chapter{Pochhammer symbols, polygamma functions, and Barnes' \texorpdfstring{$G$}{G}-function}
\index{Pochhammer symbol}\index{polygamma function}\index{Barnes G-function@Barnes $G$-function}
The supplementary MathWorld material connects the factorial bases of the first
parts with gamma quotients and their parameter derivatives.  Bell polynomials
then reappear as the universal coefficients for differentiating those
quotients.  Gamma, polygamma, and Barnes normalizations are cross-checked against
the DLMF \cite{NISTDLMF}.

\section{Pochhammer symbols and the binomial series}
For $n\in\N$, the rising Pochhammer symbol is
\[
 (z)_n=\rising zn=z(z+1)\cdots(z+n-1),\qquad (z)_0=1.
\]
Where the quotient is defined,
\begin{equation}\label{eq:pochhammer-gamma-ratio}
 (z)_n=\frac{\Gamma(z+n)}{\Gamma(z)}.
\end{equation}
This follows immediately by iterating $\Gamma(z+1)=z\Gamma(z)$ and then holds
meromorphically in $z$.  More generally one may define
\begin{equation}\label{eq:generalized-pochhammer}
 (z)_\alpha=\frac{\Gamma(z+\alpha)}{\Gamma(z)}
\end{equation}
for arbitrary complex $\alpha$ away from poles.

\begin{theorem}[Generalized binomial theorem]
\label{thm:generalized-binomial-pochhammer}
As a formal power series, and analytically for $|t|<1$,
\begin{equation}\label{eq:generalized-binomial-pochhammer}
 (1-t)^{-z}=\sum_{n\ge0}\frac{(z)_n}{n!}t^n.
\end{equation}
Consequently,
\begin{equation}\label{eq:pochhammer-vandermonde}
 (x+y)_n=\sum_{k=0}^{n}\binom nk(x)_k(y)_{n-k}.
\end{equation}
\end{theorem}
\begin{proof}
Let $A(t)=\sum a_nt^n$ be the formal solution of
$(1-t)A'(t)=zA(t)$ with $A(0)=1$.  Coefficient comparison gives
$(n+1)a_{n+1}=(z+n)a_n$, hence $a_n=(z)_n/n!$.  The analytic solution of the
same initial-value problem is $(1-t)^{-z}$, proving
\eqref{eq:generalized-binomial-pochhammer}.  Multiply the series at parameters
$x$ and $y$ and compare the coefficient of $t^n$ to obtain
\eqref{eq:pochhammer-vandermonde}.
\end{proof}

Since $(z)_n=\sum_k\stirone nk z^k$, the unsigned first-kind Stirling triangle
is exactly the monomial expansion of the Pochhammer basis.  The binomial
identity above says that this basis is of binomial type.

\section{Digamma and polygamma series}
Define
\[
 \psi(z)=\frac{\Gamma'(z)}{\Gamma(z)},\qquad
 \psi^{(m)}(z)=\frac{\dd^m}{\dd z^m}\psi(z).
\]
The Weierstrass product
\begin{equation}\label{eq:weierstrass-gamma}
 \frac1{\Gamma(z)}
 =ze^{\gamma z}\prod_{k=1}^{\infty}
 \left(1+\frac zk\right)e^{-z/k}
\end{equation}
converges locally uniformly.  Logarithmic differentiation therefore gives,
for $z$ away from the nonpositive integers,
\begin{equation}\label{eq:digamma-series}
 \psi(z)=-\gamma+\sum_{k=0}^{\infty}
 \left(\frac1{k+1}-\frac1{k+z}\right).
\end{equation}
For $m\ge1$ and $\Re z>0$, termwise differentiation yields
\begin{equation}\label{eq:polygamma-series}
 \psi^{(m)}(z)=(-1)^{m+1}m!
 \sum_{k=0}^{\infty}\frac1{(z+k)^{m+1}}.
\end{equation}
The differentiated series converges locally uniformly in the right half-plane,
which justifies every derivative.  The recurrence
\begin{equation}\label{eq:polygamma-recurrence}
 \psi^{(m)}(z+1)-\psi^{(m)}(z)
 =(-1)^m\frac{m!}{z^{m+1}}
\end{equation}
follows either from the series or by differentiating
$\psi(z+1)=\psi(z)+1/z$.

At positive integers these formulas become
\begin{align}
 \psi(n+1)&=H_n-\gamma,
 \label{eq:digamma-harmonic}\\
 \psi^{(m)}(n+1)&=(-1)^{m+1}m!
 \left(\zeta(m+1)-H_n^{(m+1)}\right),
 \label{eq:polygamma-harmonic}
\end{align}
where $H_n^{(r)}=\sum_{j=1}^{n}j^{-r}$.  Thus harmonic-number formulas for
first-kind Stirling diagonals are finite-difference shadows of polygamma
identities.

\section{Parameter derivatives of factorials}
For $r\ge1$, put
\begin{equation}\label{eq:shifted-harmonic-sum}
 H_n^{(r)}(z)=\sum_{j=0}^{n-1}\frac1{(z+j)^r}.
\end{equation}
Then \eqref{eq:polygamma-recurrence} gives
\begin{equation}\label{eq:polygamma-difference-harmonic}
 \psi^{(r-1)}(z+n)-\psi^{(r-1)}(z)
 =(-1)^{r-1}(r-1)!H_n^{(r)}(z).
\end{equation}

\begin{theorem}[Bell-polynomial parameter derivatives]
\label{thm:pochhammer-parameter-derivatives}
For $m\ge0$,
\begin{equation}\label{eq:pochhammer-parameter-derivatives}
 \frac{\dd^m}{\dd z^m}(z)_n
 =(z)_n\BellC m\bigl(L_1,L_2,\ldots,L_m\bigr),
\end{equation}
where
\begin{equation}\label{eq:pochhammer-log-derivatives}
 L_r=(-1)^{r-1}(r-1)!H_n^{(r)}(z).
\end{equation}
More generally, for the generalized symbol $(z)_\alpha$ one replaces $L_r$ by
$\psi^{(r-1)}(z+\alpha)-\psi^{(r-1)}(z)$.
\end{theorem}
\begin{proof}
Write
\[
 (z)_n=\exp\!\left(\log\Gamma(z+n)-\log\Gamma(z)\right).
\]
The $r$th derivative of the exponent is the polygamma difference in
\eqref{eq:polygamma-difference-harmonic}.  The exponential specialization
\eqref{eq:exp-faa} now gives
\eqref{eq:pochhammer-parameter-derivatives}.  The same argument with $n$
replaced by $\alpha$ proves the generalized statement.
\end{proof}

For example,
\begin{align}
 \frac{\dd}{\dd z}(z)_n
 &=(z)_nH_n^{(1)}(z),
 \label{eq:pochhammer-first-derivative}\\
 \frac{\dd^2}{\dd z^2}(z)_n
 &=(z)_n\left(H_n^{(1)}(z)^2-H_n^{(2)}(z)\right),
 \label{eq:pochhammer-second-derivative}\\
 \frac{\dd^3}{\dd z^3}(z)_n
 &=(z)_n\left(H_n^{(1)}(z)^3
 -3H_n^{(1)}(z)H_n^{(2)}(z)+2H_n^{(3)}(z)\right).
 \label{eq:pochhammer-third-derivative}
\end{align}
The same calculation gives the finite parameter-shift expansion
\begin{equation}\label{eq:pochhammer-shift-exponential}
 \frac{(z+h)_n}{(z)_n}
 =\exp\!\left(
 \sum_{r\ge1}\frac{(-1)^{r-1}}rH_n^{(r)}(z)h^r
 \right),
\end{equation}
valid formally in $h$ and analytically until a shifted pole is encountered.
Extracting powers of $h$ again produces complete Bell polynomials.

\section{Barnes' \texorpdfstring{$G$}{G}-function}
The Barnes function is the canonical interpolation of a product of factorials.
Define it initially by the locally uniformly convergent Weierstrass product
\begin{multline}\label{eq:barnes-g-product}
 G(z+1)=(2\pi)^{z/2}
 \exp\!\left(-\frac12z(z+1)-\frac12\gamma z^2\right)\\
 {}\times\prod_{k=1}^{\infty}
 \left(1+\frac zk\right)^k
 \exp\!\left(-z+\frac{z^2}{2k}\right).
\end{multline}
Indeed,
\[
 k\log\left(1+\frac zk\right)-z+\frac{z^2}{2k}=O(k^{-2})
\]
locally uniformly in $z$, so the logarithm converges normally away from the
zeros and the product defines an entire function.

\begin{theorem}[Functional equation and factorial interpolation]
\label{thm:barnes-g-functional}
The product \eqref{eq:barnes-g-product} satisfies
\begin{equation}\label{eq:barnes-g-functional}
 G(z+1)=\Gamma(z)G(z),\qquad G(1)=1.
\end{equation}
Consequently,
\begin{equation}\label{eq:barnes-g-integers}
 G(N+1)=\prod_{k=1}^{N-1}k!
 \qquad(N\ge1).
\end{equation}
\end{theorem}
\begin{proof}
The value $G(1)=1$ follows by setting $z=0$ in the product.  Divide the product
at $z$ by the one at $z-1$ and truncate at $k=N$.  The algebraic part
telescopes:
\begin{equation}\label{eq:barnes-product-telescope}
 \prod_{k=1}^{N}
 \left(\frac{k+z}{k+z-1}\right)^k
 =\frac{(N+z)^N\Gamma(z)}{\Gamma(N+z)}.
\end{equation}
Including the exponential factors gives
\begin{multline*}
 \frac{G(z+1)}{G(z)}
 =\Gamma(z)\sqrt{2\pi}
 e^{-z-\gamma z+\gamma/2}\\
 {}\times\lim_{N\to\infty}
 \frac{(N+z)^N}{\Gamma(N+z)}
 e^{-N+(z-1/2)H_N}.
\end{multline*}
Stirling's formula and $H_N=\log N+\gamma+o(1)$ show that the limit on the
second line is
\[
 \frac{e^{z+\gamma(z-1/2)}}{\sqrt{2\pi}},
\]
which cancels the preceding prefactor.
Thus the ratio is $\Gamma(z)$.  Iteration at positive integers proves
\eqref{eq:barnes-g-integers}.
\end{proof}

\begin{theorem}[Alexeiewsky--Barnes integral]
\label{thm:barnes-integral}
On any simply connected domain where consistent logarithm branches are chosen,
\begin{multline}\label{eq:barnes-integral}
 \log G(z+1)=\frac z2\log(2\pi)-\frac12z(z+1)
 +z\log\Gamma(z+1)\\
 {}-\int_0^z\log\Gamma(t+1)\,\dd t.
\end{multline}
\end{theorem}
\begin{proof}
Logarithmic differentiation of \eqref{eq:barnes-g-product}, justified by
normal convergence, gives
\[
 \frac{\dd}{\dd z}\log G(z+1)
 =\frac12\log(2\pi)-z-\frac12-\gamma z
 +\sum_{k\ge1}\frac{z^2}{k(k+z)}.
\]
From \eqref{eq:digamma-series},
$\psi(z+1)=-\gamma+\sum_{k\ge1}z/[k(k+z)]$; hence the last two terms combine
to $z\psi(z+1)$.  The derivative of the right side of
\eqref{eq:barnes-integral} is also
$\frac12\log(2\pi)-z-\frac12+z\psi(z+1)$.  Both sides vanish at $z=0$, so
they are equal.
\end{proof}

\section{Euler--Maclaurin asymptotics for Barnes \texorpdfstring{$G$}{G}}
Let $A$ be the Glaisher--Kinkelin constant, characterized by
\begin{equation}\label{eq:glaisher-definition}
 \log A=\lim_{N\to\infty}\left
 \{\sum_{k=1}^{N}k\log k
 -\left(\frac{N^2}{2}+\frac N2+\frac1{12}\right)\log N
 +\frac{N^2}{4}\right\}.
\end{equation}

\begin{theorem}[Barnes asymptotic expansion]
\label{thm:barnes-asymptotic}
For every fixed $m\ge1$ and positive integers $N\to\infty$,
\begin{multline}\label{eq:barnes-asymptotic-gamma-form}
 \log G(N+1)
 =\frac{N^2}{4}+N\log\Gamma(N+1)
 -\left(\frac{N(N+1)}2+\frac1{12}\right)\log N-\log A\\
 {}+\sum_{r=1}^{m-1}
 \frac{\beta_{2r+2}}
 {2r(2r+1)(2r+2)N^{2r}}
 +O(N^{-2m}).
\end{multline}
Equivalently,
\begin{multline}\label{eq:barnes-asymptotic-elementary}
 \log G(N+1)
 =\frac{N^2}{2}\log N-\frac{3N^2}{4}
 +\frac N2\log(2\pi)-\frac1{12}\log N
 +\left(\frac1{12}-\log A\right)\\
 {}+\sum_{r=1}^{m-1}
 \frac{\beta_{2r+2}}{4r(r+1)N^{2r}}
 +O(N^{-2m}).
\end{multline}
The constant $\frac1{12}-\log A$ is also $\zeta'(-1)$.
\end{theorem}
\begin{proof}
Apply \cref{thm:euler-maclaurin} to $f(x)=x\log x$.  Since
$f^{(q)}(x)=(-1)^q(q-2)!x^{1-q}$ for $q\ge2$, collecting the upper-end terms
gives
\begin{multline}\label{eq:hyperfactorial-asymptotic}
 \sum_{k=1}^{N}k\log k
 =\left(\frac{N^2}{2}+\frac N2+\frac1{12}\right)\log N
 -\frac{N^2}{4}+\log A\\
 {}-\sum_{r=1}^{m-1}
 \frac{\beta_{2r+2}}
 {2r(2r+1)(2r+2)N^{2r}}+O(N^{-2m}).
\end{multline}
The existence of the constant and the stated error follow from the exact
Euler--Maclaurin remainder bound.  From
\eqref{eq:barnes-g-integers}, rearranging the triangular sum of logarithms
gives the exact identity
\[
 \log G(N+1)=N\log\Gamma(N+1)-\sum_{k=1}^{N}k\log k.
\]
Substitute \eqref{eq:hyperfactorial-asymptotic} to obtain
\eqref{eq:barnes-asymptotic-gamma-form}.  Substitution of Stirling's expansion
\eqref{eq:stirling-euler-maclaurin} and collection of equal inverse powers gives
\eqref{eq:barnes-asymptotic-elementary}; the coefficient identity is
\[
 \frac{\beta_{2r+2}}{(2r+2)(2r+1)}
 +\frac{\beta_{2r+2}}{2r(2r+1)(2r+2)}
 =\frac{\beta_{2r+2}}{4r(r+1)}.
\]
The equality $\zeta'(-1)=1/12-\log A$ follows by differentiating the
Euler--Maclaurin continuation of the zeta function at $s=-1$.
\end{proof}

The sectorial complex expansion has the same coefficients and holds away from
the negative real axis.  The integer proof above already exhibits the essential
mechanism: Barnes $G$ packages a second summation of $\log z$, so one more layer
of Bernoulli corrections appears beyond Stirling's expansion for $\Gamma$.

\part{Differentiation, Composition, and Reversion}

\chapter{Fa\`a di Bruno's formula}
\index{Faà di Bruno formula@Fa\`a di Bruno formula}

\section{The set-partition form}
For a block $B\subseteq[n]$, write $|B|$ for its size.
\begin{theorem}[Fa\`a di Bruno, partition form]\label{thm:faa-partition}
For $n\ge0$ and sufficiently differentiable $f,g$,
\begin{equation}\label{eq:faa-partition}
 D^n(f\circ g)(x)=
 \sum_{\pi\in\setpart([n])}
 f^{(|\pi|)}(g(x))\prod_{B\in\pi}g^{(|B|)}(x).
\end{equation}
For $n=0$, the empty partition contributes $f(g(x))$.
\end{theorem}
\begin{proof}
Induct on $n$.  Differentiate a term indexed by a partition $\pi$ of $[n]$.
If the derivative hits $f^{(|\pi|)}(g)$, the chain rule contributes $g'$ and
creates the new singleton block $\{n+1\}$.  If it hits the factor associated with
a block $B$, it changes $g^{(|B|)}$ to $g^{(|B|+1)}$, which inserts $n+1$ into
$B$.  Every partition of $[n+1]$ arises uniquely in exactly one of these ways:
remove $n+1$, either deleting its singleton block or removing it from its
non-singleton block.  Hence all terms occur once.
\end{proof}

\section{Multiplicity and Bell-polynomial forms}
Grouping partitions by the number $m_j$ of blocks of size $j$ gives the classical
formula.
\begin{theorem}[Fa\`a di Bruno, multiplicity form]\label{thm:faa-multiplicity}
\begin{multline}\label{eq:faa-multiplicity}
 D^n f(g(x))=
 \sum_{\substack{m_1,\ldots,m_n\ge0\\
                  m_1+2m_2+\cdots+nm_n=n}}
 \frac{n!}{m_1!m_2!\cdots m_n!}\\
 {}\times f^{(m_1+\cdots+m_n)}(g(x))
 \prod_{j=1}^{n}\left(\frac{g^{(j)}(x)}{j!}\right)^{m_j}.
\end{multline}
Equivalently,
\begin{equation}\label{eq:faa-bell}
 D^n f(g(x))=
 \sum_{k=0}^{n}f^{(k)}(g(x))
 \BellP nk\bigl(g'(x),g''(x),\ldots,g^{(n-k+1)}(x)\bigr).
\end{equation}
\end{theorem}
\begin{proof}
For fixed multiplicities $(m_j)$, the number of set partitions of $[n]$ with
$m_j$ blocks of size $j$ is
$n!/\prod_j((j!)^{m_j}m_j!)$.  Substitution in
\eqref{eq:faa-partition} gives \eqref{eq:faa-multiplicity}.  Grouping further by
$k=\sum m_j$ and invoking \eqref{eq:partial-bell-definition} gives
\eqref{eq:faa-bell}.
\end{proof}

For $n=4$ this reads
\begin{align}
 (f\circ g)^{(4)}={}&f^{(4)}(g)(g')^4
 +6f^{(3)}(g)(g')^2g''+3f''(g)(g'')^2\notag\\
 &+4f''(g)g'g^{(3)}+f'(g)g^{(4)}.
 \label{eq:faa-n4}
\end{align}
The coefficients $1,6,3,4,1$ count partitions of types
$1+1+1+1$, $2+1+1$, $2+2$, $3+1$, and $4$ respectively.  This also proves the
expanded example in the archive without four separate differentiations.

\section{The exponential special case}
Because every derivative of $e^u$ equals $e^u$, summing the partial Bell
polynomials in \eqref{eq:faa-bell} yields
\begin{equation}\label{eq:exp-faa}
 D^n e^{g(x)}=e^{g(x)}
 \BellC n\bigl(g'(x),g''(x),\ldots,g^{(n)}(x)\bigr).
\end{equation}
When $g$ is a cumulant generating function, this is precisely the
moment--cumulant relation of \cref{thm:moment-cumulant}.

\section{Mixed partial derivatives}
For $B\subseteq[n]$, write
$\partial_B=\prod_{j\in B}\partial/\partial x_j$.
\begin{theorem}[Multivariate Fa\`a di Bruno for distinct variables]
\label{thm:faa-multivariate}
If $y=g(x_1,\ldots,x_n)$, then
\begin{equation}\label{eq:faa-multivariate}
 \frac{\partial^n}{\partial x_1\cdots\partial x_n}f(y)
 =\sum_{\pi\in\setpart([n])}f^{(|\pi|)}(y)
   \prod_{B\in\pi}\partial_B y.
\end{equation}
\end{theorem}
\begin{proof}
The proof of \cref{thm:faa-partition} applies verbatim, except that the new
derivative is $\partial/\partial x_{n+1}$.  It either creates a singleton block or
joins the new index to one existing block.
\end{proof}

For three variables, \eqref{eq:faa-multivariate} becomes
\begin{align*}
 \partial_1\partial_2\partial_3f(y)
 ={}&f'(y)\partial_{123}y\\
 &+f''(y)\bigl((\partial_1y)(\partial_{23}y)
 +(\partial_2y)(\partial_{13}y)+(\partial_3y)(\partial_{12}y)\bigr)\\
 &+f^{(3)}(y)(\partial_1y)(\partial_2y)(\partial_3y).
\end{align*}
For repeated variables or vector-valued inner maps, one labels derivative slots
first and then contracts tensor indices; the same partition lattice remains the
combinatorial skeleton.


\section{Canonical outer functions}
Fa\`a di Bruno becomes a library of higher-derivative rules once the derivatives
of the outer function are inserted.

\begin{theorem}[Power, logarithm, reciprocal, and trigonometric rules]
\label{thm:faa-canonical-outer}
For $n\ge1$ and at points where the expressions are defined,
\begin{align}
 D^n g(x)^\alpha
 &=\sum_{k=1}^{n}\falling{\alpha}{k}g(x)^{\alpha-k}
   \BellP nk(g',g'',\ldots),
 \label{eq:faa-power}\\
 D^n\log g(x)
 &=\sum_{k=1}^{n}(-1)^{k-1}(k-1)!\,
   \frac{\BellP nk(g',g'',\ldots)}{g(x)^k},
 \label{eq:faa-log}\\
 D^n\frac1{g(x)}
 &=\sum_{k=1}^{n}(-1)^kk!\,
   \frac{\BellP nk(g',g'',\ldots)}{g(x)^{k+1}},
 \label{eq:faa-reciprocal}\\
 D^n\sin(g(x))
 &=\sum_{k=1}^{n}\sin\!\left(g(x)+\frac{k\pi}{2}\right)
   \BellP nk(g',g'',\ldots),
 \label{eq:faa-sine}\\
 D^n\cos(g(x))
 &=\sum_{k=1}^{n}\cos\!\left(g(x)+\frac{k\pi}{2}\right)
   \BellP nk(g',g'',\ldots).
 \label{eq:faa-cosine}
\end{align}
For $n=0$, the original function is recovered separately.
\end{theorem}
\begin{proof}
In \eqref{eq:faa-bell}, insert respectively
\[
 D^ku^\alpha=\falling{\alpha}{k}u^{\alpha-k},\qquad
 D^k\log u=(-1)^{k-1}(k-1)!u^{-k},
\]
\[
 D^ku^{-1}=(-1)^kk!u^{-k-1},\qquad
 D^k\sin u=\sin(u+k\pi/2),\qquad
 D^k\cos u=\cos(u+k\pi/2).
\]
No further combinatorial calculation is required.
\end{proof}

\section{Leibniz rules and higher quotients}
\begin{theorem}[Multinomial Leibniz rule]
\label{thm:multinomial-leibniz}
For $q$ sufficiently differentiable functions,
\begin{equation}\label{eq:multinomial-leibniz}
 D^n\prod_{r=1}^{q}f_r(x)
 =\sum_{j_1+\cdots+j_q=n}
   \binom{n}{j_1,\ldots,j_q}
   \prod_{r=1}^{q}f_r^{(j_r)}(x).
\end{equation}
\end{theorem}
\begin{proof}
Taylor-expand every factor at $x$:
\[
 \prod_{r=1}^{q}f_r(x+h)
 =\prod_{r=1}^{q}\sum_{j_r\ge0}f_r^{(j_r)}(x)\frac{h^{j_r}}{j_r!}.
\]
Extracting $n![h^n]$ gives \eqref{eq:multinomial-leibniz}.
\end{proof}

\begin{corollary}[Higher quotient rule]\label{cor:higher-quotient}
If $g(x)\ne0$, then
\begin{equation}\label{eq:higher-quotient}
 D^n\frac{f(x)}{g(x)}
 =\sum_{r=0}^{n}\binom nr f^{(n-r)}(x)
   \sum_{k=0}^{r}(-1)^kk!\,
   \frac{\BellP rk(g',g'',\ldots)}{g(x)^{k+1}},
\end{equation}
where $\BellP00=1$.
\end{corollary}
\begin{proof}
Apply the two-factor case of \cref{thm:multinomial-leibniz} to
$f\cdot g^{-1}$ and use \eqref{eq:faa-reciprocal}, including its $r=0$
term.
\end{proof}

\section{Coordinate-free Fa\`a di Bruno formula}
\index{Frechet derivative@Fr\'echet derivative}
Let $E,F,G$ be Banach spaces, let $g:E\to F$ and $f:F\to G$, and write
$D^r$ for the symmetric $r$-linear Fr\'echet derivative.

\begin{theorem}[Fr\'echet Fa\`a di Bruno]
\label{thm:faa-frechet}
For directions $h_1,\ldots,h_n\in E$,
\begin{multline}\label{eq:faa-frechet}
 D^n(f\circ g)(x)[h_1,\ldots,h_n]\\
 =\sum_{\pi\in\setpart([n])}
 D^{|\pi|}f(g(x))
 \Bigl[
   D^{|B|}g(x)[h_i:i\in B]
 \Bigr]_{B\in\pi}.
\end{multline}
One vector of $F$ is supplied to $D^{|\pi|}f$ for each block $B$ of $\pi$.
\end{theorem}
\begin{proof}
Induct on $n$.  The case $n=1$ is the Fr\'echet chain rule.  Differentiate a
term indexed by a partition $\pi$ of $[n]$ in the new direction $h_{n+1}$.
If the derivative strikes the outer multilinear map, the chain rule creates a
new argument $Dg(x)[h_{n+1}]$, corresponding to the singleton block
$\{n+1\}$.  If it strikes the inner factor belonging to a block $B$, that
factor becomes
\[
 D^{|B|+1}g(x)[h_i:i\in B,h_{n+1}],
\]
which adjoins $n+1$ to $B$.  Every partition of $[n+1]$ is obtained uniquely
by deleting $n+1$ either from its singleton block or from its unique
non-singleton block.  Hence every term occurs exactly once.
\end{proof}

When $E=F=G=\R$ and all directions equal $1$, grouping partitions by their
block-size multiplicities recovers \eqref{eq:faa-multiplicity}.  Coordinate
indices are therefore bookkeeping added after the partition structure, not a
new combinatorial theorem.

\section{Derivatives of inverse functions}
Let $y=f(x)$ with $f'(x)\ne0$, and put $x=g(y)=f^{-1}(y)$.

\begin{theorem}[Recursive inverse-derivative formula]
\label{thm:inverse-derivative-recursion}
For $n\ge2$,
\begin{equation}\label{eq:inverse-derivative-recursion}
 g^{(n)}(y)=-\frac1{f'(g(y))}
 \sum_{k=2}^{n}f^{(k)}(g(y))
 \BellP nk\bigl(g'(y),g''(y),\ldots,g^{(n-k+1)}(y)\bigr),
\end{equation}
with $g'(y)=1/f'(g(y))$.  In particular, writing derivatives of $f$ at
$x=g(y)$,
\begin{align}
 g'&=\frac1{f'},
 \label{eq:inverse-derivative-1}\\
 g''&=-\frac{f''}{(f')^3},
 \label{eq:inverse-derivative-2}\\
 g'''&=\frac{3(f'')^2-f'f'''}{(f')^5},
 \label{eq:inverse-derivative-3}\\
 g^{(4)}&=\frac{-15(f'')^3+10f'f''f'''-(f')^2f^{(4)}}{(f')^7}.
 \label{eq:inverse-derivative-4}
\end{align}
\end{theorem}
\begin{proof}
For $n\ge2$, differentiate $f(g(y))=y$ $n$ times and apply
\eqref{eq:faa-bell}.  The term with $k=1$ is
\[
 f'(g(y))g^{(n)}(y).
\]
Moving all terms with $k\ge2$ to the other side gives
\eqref{eq:inverse-derivative-recursion}.  Substitution of the already computed
lower derivatives gives the four displayed cases.  The recursion is triangular
because $\BellP nk$ with $k\ge2$ involves derivatives of $g$ of order at most
$n-1$.
\end{proof}

Lagrange inversion will replace this triangular recursion by a direct
coefficient formula.  The two descriptions are complementary: Fa\`a di Bruno is
local in derivative order, whereas Lagrange--B\"urmann is global in the defining
series.

\section{Nested composition}
For three scalar functions, repeated Fa\`a di Bruno gives the explicit finite
formula
\begin{align}
 D^n f(g(h(x)))
 &={}
 \sum_{k=0}^{n}f^{(k)}(g(h(x)))
 \BellP nk\bigl((g\circ h)',(g\circ h)'',\ldots\bigr),
 \label{eq:threefold-composition-outer}\\
 (g\circ h)^{(j)}
 &={}
 \sum_{\ell=0}^{j}g^{(\ell)}(h(x))
 \BellP j\ell(h',h'',\ldots).
 \label{eq:threefold-composition-inner}
\end{align}
Substituting the second line into the first is a closed formula involving only
derivatives of $f,g,h$.  Combinatorially, the inner partition groups derivative
labels that strike one occurrence of $h$, and the outer partition groups those
blocks according to occurrences of $g$.  The canonical flattening of nested
partitions is the combinatorial reason formal composition is associative.

\begin{auditbox}[title={Composition normalizations}]
For ordinary series, $\OBellP nk$ extracts $[z^n]G(z)^k$; for exponential
series, $\BellP nk$ extracts $n![z^n]G(z)^k/k!$; in derivative formulas the
arguments are the actual derivatives $g^{(j)}(x)$.  The conversion
\[
 \OBellP nk(g_1,g_2,\ldots)
 =\frac{k!}{n!}\BellP nk(1!g_1,2!g_2,\ldots)
\]
accounts for the factors $n!$, $k!$, and $j!$ that differ among the five
drafts.
\end{auditbox}

\chapter{Composition of ordinary and exponential series}

\section{Ordinary formal power series}
Let
\[
 F(u)=\sum_{k\ge0}a_ku^k,
 \qquad
 G(x)=\sum_{j\ge1}b_jx^j,
 \qquad
 F(G(x))=\sum_{n\ge0}c_nx^n.
\]
\begin{theorem}[Ordinary composition]\label{thm:ordinary-composition}
For $n\ge1$,
\begin{align}
 c_n&=\sum_{k=1}^{n}a_k
 \sum_{\substack{i_1+\cdots+i_k=n\\i_r\ge1}}
 b_{i_1}\cdots b_{i_k},
 \label{eq:ordinary-composition-compositions}\\
 &=\sum_{k=1}^{n}a_k
 \sum_{\substack{\pi_1+\cdots+\pi_n=k\\
                   \pi_1+2\pi_2+\cdots+n\pi_n=n}}
 \binom{k}{\pi_1,\ldots,\pi_n}
 b_1^{\pi_1}\cdots b_n^{\pi_n}
 \label{eq:ordinary-composition-multiplicities}\\
 &=\sum_{k=1}^{n}a_k\OBellP nk(b_1,\ldots,b_{n-k+1}),
 \label{eq:ordinary-composition-bell}
\end{align}
and $c_0=a_0$.
\end{theorem}
\begin{proof}
The coefficient of $x^n$ in $G(x)^k$ is obtained by choosing an ordered
composition $i_1+\cdots+i_k=n$, proving the first formula.  Collect equal part
sizes; the number of orderings with multiplicities $\pi_j$ is the multinomial
coefficient, proving the second.  The third is the definition of the ordinary
Bell polynomial.
\end{proof}

Two useful specializations are immediate.  If $F(u)=e^u$ and the series are
exponentially normalized, one obtains the exponential formula.  If
$A(x)=1+\sum_{n\ge1}a_nx^n$, then
\begin{equation}\label{eq:reciprocal-ordinary-bell}
 [x^n]\frac1{A(x)}
 =\sum_{k=1}^{n}(-1)^k\OBellP nk(a_1,\ldots,a_{n-k+1})
 \qquad(n\ge1),
\end{equation}
by taking $F(u)=1/(1-u)$ and $G(x)=-\sum_{n\ge1}a_nx^n$.

\section{Exponential formal power series}
Let
\[
 F(x)=\sum_{j\ge1}a_j\frac{x^j}{j!},\qquad
 G(u)=\sum_{k\ge0}b_k\frac{u^k}{k!},\qquad
 G(F(x))=\sum_{n\ge0}c_n\frac{x^n}{n!}.
\]
\begin{theorem}[Exponential composition]\label{thm:exponential-composition}
\begin{equation}\label{eq:exponential-composition}
 c_n=\sum_{k=0}^{n}b_k\BellP nk(a_1,\ldots,a_{n-k+1}).
\end{equation}
Equivalently,
\begin{equation}\label{eq:exponential-composition-partitions}
 c_n=\sum_{\pi\in\setpart([n])}b_{|\pi|}\prod_{B\in\pi}a_{|B|}.
\end{equation}
\end{theorem}
\begin{proof}
Substitute $F(x)$ into $G$ and apply
\eqref{eq:partial-bell-egf}.  The set-partition form is
\cref{thm:bell-poly-partitions} with an additional weight $b_k$ for the number of
blocks.
\end{proof}

Taking $b_k=1$ gives the complete Bell-polynomial exponential identity
\begin{equation}\label{eq:exp-complete-bell-again}
 \exp\!\left(\sum_{j\ge1}a_j\frac{x^j}{j!}\right)
 =\sum_{n\ge0}\BellC n(a_1,\ldots,a_n)\frac{x^n}{n!}.
\end{equation}
If the series converge, \eqref{eq:exponential-composition} is exactly
\eqref{eq:faa-bell} evaluated at the expansion point; formally, no convergence
notion is needed.

\chapter{Lagrange inversion and Lagrange--Bürmann}
\index{Lagrange inversion}

\section{Formal theorem}
Let $R$ be a commutative $\Q$-algebra and let $\phi(w)\in R[[w]]$ have invertible
constant term.  There is a unique $g(z)\in zR[[z]]$ satisfying
\begin{equation}\label{eq:lagrange-functional}
 g(z)=z\phi(g(z)).
\end{equation}
Uniqueness follows recursively because the coefficient of $z^n$ on the right
depends only on coefficients of $g$ below degree $n$.

\begin{theorem}[Lagrange--Bürmann]\label{thm:lagrange-burmann}
For $n\ge1$ and every formal series $H$,
\begin{align}
 \coeff zn H(g(z))
 &=\frac1n\coeff w{n-1}H'(w)\phi(w)^n,
 \label{eq:lagrange-burmann}\\
 &=\coeff wn H(w)\phi(w)^{n-1}\bigl(\phi(w)-w\phi'(w)\bigr).
 \label{eq:lagrange-burmann-alt}
\end{align}
In particular,
\begin{equation}\label{eq:lagrange-basic}
 \coeff zn g(z)=\frac1n\coeff w{n-1}\phi(w)^n.
\end{equation}
\end{theorem}
\begin{proof}
Put $f(w)=w/\phi(w)$, so that $f(g(z))=z$.  Formal residues and the change of
variables theorem \cref{thm:res-subst} give
\begin{align*}
 \coeff zn H(g(z))
 &=\res_z H(g(z))z^{-n-1}\dd z\\
 &=\res_w H(w)f(w)^{-n-1}f'(w)\dd w\\
 &=-\frac1n\res_w H(w)(f(w)^{-n})'\dd w\\
 &=\frac1n\res_w H'(w)f(w)^{-n}\dd w,
\end{align*}
where the residue of a derivative is zero.  Since
$f(w)^{-n}=w^{-n}\phi(w)^n$, this is
\eqref{eq:lagrange-burmann}.  Applying the same zero-residue rule to
$H(w)\phi(w)^nw^{-n}$ gives
\[
 \frac1n\res H'\phi^nw^{-n}
 =\res H\phi^nw^{-n-1}-\res H\phi^{n-1}\phi'w^{-n},
\]
which is \eqref{eq:lagrange-burmann-alt}.  Take $H(w)=w$ for
\eqref{eq:lagrange-basic}.
\end{proof}

\section{Analytic inverse functions}
\begin{theorem}[Analytic Lagrange inversion]\label{thm:analytic-lagrange}
Let $f$ be analytic near $a$, let $b=f(a)$, and suppose $f'(a)\ne0$.  Its local
inverse $g$ is analytic near $b$ and has
\begin{equation}\label{eq:analytic-lagrange-series}
 g(z)=a+\sum_{n\ge1}g_n\frac{(z-b)^n}{n!},
\end{equation}
where
\begin{equation}\label{eq:analytic-lagrange-coeff}
 g_n=\left.\frac{\dd^{n-1}}{\dd w^{n-1}}
 \left(\frac{w-a}{f(w)-f(a)}\right)^n\right|_{w=a}.
\end{equation}
The series has a positive radius of convergence.
\end{theorem}
\begin{proof}
The analytic inverse-function theorem gives an analytic local inverse and therefore
a convergent Taylor series.  Put $u=w-a$, $\zeta=z-b$, and
\[
 \phi(u)=\frac{u}{f(a+u)-f(a)};
\]
its constant term is $1/f'(a)$.  The inverse displacement satisfies
$u=\zeta\phi(u)$.  Formula \eqref{eq:lagrange-basic}, multiplied by $n!$, is
exactly \eqref{eq:analytic-lagrange-coeff}.
\end{proof}

If $f'(a)=0$ but $f(w)-f(a)=c(w-a)^q+O((w-a)^{q+1})$ with $c\ne0$, factor
$f(w)-f(a)=(w-a)^qU(w)$ with $U(a)=c$.  Choosing one of the $q$ values of
$c^{-1/q}$ reduces the equation to ordinary inversion in the variable
$(z-f(a))^{1/q}$.  Thus the inverse has $q$ local Puiseux branches.  This is the
precise sense in which Lagrange inversion extends to a critical point; a
single-valued Taylor inverse cannot exist there.

\section{Inverse coefficients in Bell-polynomial form}
Suppose
\[
 f(w)=\sum_{j\ge1}f_j\frac{w^j}{j!},\qquad f_1\ne0,
 \qquad
 g(z)=f^{\langle-1\rangle}(z)=\sum_{n\ge1}g_n\frac{z^n}{n!}.
\]
Define
\begin{equation}\label{eq:fhat}
 \widehat f_j=\frac{f_{j+1}}{(j+1)f_1}.
\end{equation}
\begin{theorem}[Explicit inverse coefficients]\label{thm:inverse-bell-coeff}
One has $g_1=f_1^{-1}$ and, for $n\ge2$,
\begin{equation}\label{eq:inverse-bell-coeff}
 g_n=\frac1{f_1^n}\sum_{k=1}^{n-1}(-1)^k\rising nk
 \BellP{n-1}{k}(\widehat f_1,\ldots,\widehat f_{n-k}).
\end{equation}
\end{theorem}
\begin{proof}
Write
\[
 \frac{f(w)}{f_1w}=1+U(w),
 \qquad U(w)=\sum_{j\ge1}\widehat f_j\frac{w^j}{j!}.
\]
By \eqref{eq:lagrange-basic},
$g_n=(n-1)!\coeff w{n-1}(w/f(w))^n$.  Now
\[
 (1+U)^{-n}=\sum_{k\ge0}(-1)^k\binom{n+k-1}{k}U^k
 =\sum_{k\ge0}(-1)^k\frac{\rising nk}{k!}U^k.
\]
Apply \eqref{eq:partial-bell-egf} and extract degree $n-1$; the $k=0$ term
vanishes when $n\ge2$.
\end{proof}

\section{Associahedra, plane trees, and signed face sums}
\index{associahedron}
The supplied drafts use both exponential and ordinary coefficient conventions
near the associahedral formula.  The geometric statement is an ordinary-series
identity, so we isolate that normalization.  Let
\begin{equation}\label{eq:associa-series}
 f(x)=x+\sum_{n\ge1}a_nx^{n+1},\qquad
 g(y)=f^{\langle-1\rangle}(y)=y+\sum_{n\ge1}b_ny^{n+1}.
\end{equation}
Let $K_n$ denote the associahedron of dimension $n-1$.  Its faces may be
represented by plane rooted trees with $n+1$ leaves and no unary internal
vertices.  If such a tree has $m_j$ internal vertices of arity $j+1$, put
\begin{equation}\label{eq:associa-profile}
 k=\sum_{j\ge1}m_j,\qquad
 \sum_{j\ge1}jm_j=n,\qquad
 a_F=\prod_{j\ge1}a_j^{m_j}.
\end{equation}
The corresponding face has dimension $n-k$.

\begin{lemma}[Plane-tree profile count]\label{lem:plane-tree-profile}
For a profile $(m_1,m_2,\ldots)$ satisfying \eqref{eq:associa-profile}, the
number of plane rooted trees with that profile and $n+1$ leaves is
\begin{equation}\label{eq:plane-tree-profile-count}
 \frac{(n+k)!}{(n+1)!\prod_{j\ge1}m_j!}.
\end{equation}
\end{lemma}
\begin{proof}
List the vertices in preorder and record $d(v)-1$, where $d(v)$ is the number
of children.  A leaf contributes $-1$, while an internal vertex of arity
$j+1$ contributes $j$.  The word has $n+k+1$ entries, total sum $-1$, and
multiplicities $n+1,m_1,m_2,\ldots$.  The cycle lemma \cite{DvoretzkyMotzkin1947} says that exactly one
cyclic shift of each such word has all proper partial sums nonnegative; that
shift is precisely the preorder degree word of a plane rooted tree.  Because
the total sum is $-1$, no word has a nontrivial period, so every cyclic orbit
has size $n+k+1$.  Hence the number of trees is
\[
 \frac1{n+k+1}
 \binom{n+k+1}{n+1,m_1,m_2,\ldots}
 =\frac{(n+k)!}{(n+1)!\prod_jm_j!}.
\]
\end{proof}

\begin{theorem}[Associahedral Lagrange inversion]
\label{thm:associahedral-inversion}
For every $n\ge1$,
\begin{align}
 b_n
 &=\sum_{F\text{ a face of }K_n}(-1)^{n-\dim F}a_F,
 \label{eq:associahedral-face-sum}\\
 &=\sum_{\substack{m_1,m_2,\ldots\ge0\\\sum j m_j=n}}
 (-1)^k\frac{(n+k)!}{(n+1)!\prod_{j\ge1}m_j!}
 \prod_{j\ge1}a_j^{m_j},
 \qquad k=\sum_jm_j.
 \label{eq:associahedral-frequency}
\end{align}
\end{theorem}
\begin{proof}
Write $f(x)=xA(x)$, where
$A(x)=1+\sum_{j\ge1}a_jx^j$.  The inverse equation is
$g=y/A(g)$, so Lagrange inversion gives
\[
 b_n=[y^{n+1}]g(y)=\frac1{n+1}[x^n]A(x)^{-(n+1)}.
\]
The generalized binomial theorem and the multinomial theorem give
\begin{align*}
 b_n
 &=\frac1{n+1}\sum_{k=0}^{n}(-1)^k\binom{n+k}{k}
 [x^n]\left(\sum_{j\ge1}a_jx^j\right)^k\\
 &=\sum_{\substack{\sum jm_j=n\\\sum m_j=k}}
 (-1)^k\frac{(n+k)!}{(n+1)!\prod_jm_j!}
 \prod_ja_j^{m_j},
\end{align*}
which is \eqref{eq:associahedral-frequency}.  By
\cref{lem:plane-tree-profile}, the unsigned coefficient counts the trees of
that profile.  The standard face-tree correspondence identifies these trees
with faces of $K_n$; contraction of an internal edge merges two internal
vertices, and a tree with $k$ internal vertices represents a face of dimension
$n-k$.  Thus $(-1)^k=(-1)^{n-\dim F}$, proving
\eqref{eq:associahedral-face-sum}.
\end{proof}

\begin{corollary}[The face numbers of the associahedron]
\label{cor:associahedron-f-vector}
The number $f_d(K_n)$ of $d$-dimensional faces, $0\le d\le n-1$, is
\begin{equation}\label{eq:associahedron-f-vector}
 f_d(K_n)=\frac1{n+1}\binom{n-1}{d}\binom{2n-d}{n}.
\end{equation}
In particular, $K_n$ has the Catalan number
$\frac1{n+1}\binom{2n}{n}$ of vertices, and
$\sum_d(-1)^df_d(K_n)=1$.
\end{corollary}
\begin{proof}
A $d$-face has $k=n-d$ internal vertices.  Summing
\eqref{eq:plane-tree-profile-count} over profiles with
$\sum m_j=k$ uses
\[
 \sum_{\substack{\sum jm_j=n\\\sum m_j=k}}
 \frac1{\prod_jm_j!}
 =\frac1{k!}[x^n]\left(\frac{x}{1-x}\right)^k
 =\frac1{k!}\binom{n-1}{k-1}.
\]
Substitution and $k=n-d$ give \eqref{eq:associahedron-f-vector}.  The vertex
case is $d=0$.  Finally set every $a_j=1$ in
\eqref{eq:associahedral-face-sum}.  Then $f(x)=x/(1-x)$ and
$g(y)=y/(1+y)$, so $b_n=(-1)^n$; multiplying the resulting identity by
$(-1)^n$ gives the Euler relation.
\end{proof}

The first inverse polynomials are
\begin{align*}
 b_1&=-a_1,\\
 b_2&=-a_2+2a_1^2,\\
 b_3&=-a_3+5a_1a_2-5a_1^3,\\
 b_4&=-a_4+6a_1a_3+3a_2^2-21a_1^2a_2+14a_1^4.
\end{align*}
For $K_4$, the coefficients $1,6,3,21,14$ enumerate the whole polytope, the
two product types of facets, the edges, and the vertices.  The alternating
face sum is therefore not a metaphor for inversion: it is the inverse
coefficient polynomial itself.

\chapter{Algebraic and transcendental examples}

\section{The Fuss--Catalan algebraic equation}
Let $p\ge2$ be an integer and consider
\begin{equation}\label{eq:fuss-equation}
 x^p-x+z=0,
 \qquad x(0)=0.
\end{equation}
It is equivalent to
$x=z\phi(x)$ with $\phi(x)=(1-x^{p-1})^{-1}$.

\begin{theorem}\label{thm:fuss-series}
The distinguished solution is
\begin{equation}\label{eq:fuss-series}
 x(z)=\sum_{k=0}^{\infty}\binom{pk}{k}
       \frac{z^{(p-1)k+1}}{(p-1)k+1}.
\end{equation}
Its radius of convergence is
\begin{equation}\label{eq:fuss-radius}
 R_p=(p-1)p^{-p/(p-1)},
\end{equation}
and the series converges absolutely also on $|z|=R_p$.
\end{theorem}
\begin{proof}
By \eqref{eq:lagrange-basic},
\[
 \coeff znx(z)=\frac1n\coeff w{n-1}(1-w^{p-1})^{-n}.
\]
This vanishes unless $n-1=(p-1)k$.  In that case the coefficient is
$\frac1n\binom{n+k-1}{k}=\binom{pk}{k}/((p-1)k+1)$, proving the series.

The inverse map is $z=f(x)=x-x^p$.  Its nearest critical points satisfy
$f'(x)=1-px^{p-1}=0$ and all have modulus $p^{-1/(p-1)}$; their critical values
have modulus $(p-1)p^{-p/(p-1)}$.  No local inverse can pass analytically through
a critical value, while the inverse-function theorem continues the branch up to
one, so this is the radius.  Stirling's formula gives coefficients of order
$Ck^{-3/2}R_p^{-((p-1)k+1)}$, which proves absolute convergence on the boundary.
\end{proof}

\section{Lambert's \texorpdfstring{$W$}{W} function at the origin}
The Lambert function satisfies $W(z)e^{W(z)}=z$.
\begin{theorem}\label{thm:lambert-W-zero}
On its principal branch near $0$,
\begin{equation}\label{eq:lambert-W-zero}
 W(z)=\sum_{n=1}^{\infty}(-n)^{n-1}\frac{z^n}{n!}
 =z-z^2+\frac32z^3-\frac83z^4+O(z^5),
\end{equation}
with radius of convergence $e^{-1}$.
\end{theorem}
\begin{proof}
Here $f(w)=we^w$, so $w/f(w)=e^{-w}$.  Formula
\eqref{eq:analytic-lagrange-coeff} gives
$g_n=D^{n-1}e^{-nw}|_{w=0}=(-n)^{n-1}$.  The inverse map has its nearest critical
point at $w=-1$, where $f(-1)=-e^{-1}$; this square-root branch point determines
the radius.
\end{proof}

\section{A rapidly useful expansion around \texorpdfstring{$W(e)=1$}{W(e)=1}}
Put $\xi=\log z-1$ and $u=W(z)-1$.  Then
\begin{equation}\label{eq:lambert-near-e-implicit}
 u+\log(1+u)=\xi.
\end{equation}
Reversion gives
\begin{equation}\label{eq:lambert-near-e}
 W(e^{1+\xi})=1+\frac{\xi}{2}+\frac{\xi^2}{16}
 -\frac{\xi^3}{192}-\frac{\xi^4}{3072}
 +\frac{13\xi^5}{61440}+O(\xi^6).
\end{equation}
\begin{proof}
The derivative at $u=0$ of the left side of
\eqref{eq:lambert-near-e-implicit} is $2$, so it has an ordinary inverse.
Substitute $u=c_1\xi+\cdots+c_5\xi^5$ into
$u+\log(1+u)=\xi$ and compare successive powers.  The triangular equations give
$c_1=1/2,c_2=1/16,c_3=-1/192,c_4=-1/3072,c_5=13/61440$.
\end{proof}
The closest critical points of $u+\log(1+u)$ satisfy
$1+(1+u)^{-1}=0$, hence $u=-2$; using the two adjacent logarithm values gives
images $-2\pm i\pi$.  Therefore the radius in the $\xi$-plane is
$\sqrt{4+\pi^2}$.  For positive real $z$, the expansion converges whenever
$|\log z-1|<\sqrt{4+\pi^2}$, approximately
$0.0655<z<112.63$.  Taking $\xi=-1$ yields a convergent evaluation of
$W(1)=0.567143\ldots$.

\section{Binary trees and Catalan numbers}
Let $C_n$ count rooted ordered full binary trees with $n$ internal vertices, and
let $B(z)=\sum_{n\ge0}C_nz^n$.  Removing the root gives
\begin{equation}\label{eq:binary-tree-functional}
 B(z)=1+zB(z)^2.
\end{equation}
Put $C(z)=B(z)-1$; then $C=z(1+C)^2$.  Lagrange inversion gives
\begin{equation}\label{eq:catalan-lagrange}
 C_n=\coeff znC(z)=\frac1n\coeff w{n-1}(1+w)^{2n}
 =\frac1n\binom{2n}{n-1}=\frac1{n+1}\binom{2n}{n}.
\end{equation}
This simultaneously proves the functional equation, its unique formal solution,
and the Catalan closed form.  The symbol $C_n$ is used here to avoid confusing
these binary-tree numbers with Bell numbers.


\section{Inverse trigonometric coefficient identities}
Inverse functions can evaluate coefficient extractions that are otherwise
surprisingly opaque.  Since
\[
 \frac{\dd}{\dd y}\arcsin y=(1-y^2)^{-1/2}
 =\sum_{m\ge0}\binom{2m}{m}\frac{y^{2m}}{4^m},
\]
integration gives
\begin{equation}\label{eq:arcsin-series}
 \arcsin y=\sum_{m\ge0}
 \frac1{2m+1}\binom{2m}{m}\frac{y^{2m+1}}{4^m}.
\end{equation}
On the other hand, Lagrange inversion applied to $f(x)=\sin x$ gives
\[
 [y^{2m+1}]\arcsin y
 =\frac1{2m+1}[x^{2m}]
 \left(\frac{x}{\sin x}\right)^{2m+1}.
\]
Comparison proves
\begin{equation}\label{eq:sine-power-diagonal}
 \boxed{\displaystyle
 [x^{2m}]\left(\frac{x}{\sin x}\right)^{2m+1}
 =\frac1{4^m}\binom{2m}{m}.}
\end{equation}
Likewise,
$\arctan y=\sum_{m\ge0}(-1)^my^{2m+1}/(2m+1)$ and inversion of
$f(x)=\tan x$ give
\begin{equation}\label{eq:tangent-power-diagonal}
 \boxed{\displaystyle
 [x^{2m}]\left(\frac{x}{\tan x}\right)^{2m+1}=(-1)^m.}
\end{equation}
Both proofs are complete because the elementary derivative expansions identify
the inverse coefficients, while \cref{thm:lagrange-burmann} identifies the same
coefficients with the displayed high powers.

\section{Exponential and logarithmic diagonal identities}
The inverse pair $f(x)=e^x-1$ and $g(y)=\log(1+y)$ gives
\begin{equation}\label{eq:norlund-lagrange-diagonal}
 [x^{n-1}]\left(\frac{x}{e^x-1}\right)^n=(-1)^{n-1},
 \qquad n\ge1.
\end{equation}
Indeed, \eqref{eq:lagrange-basic} says that the left side divided by $n$ is
$[y^n]\log(1+y)=(-1)^{n-1}/n$.  In the notation of
\cref{eq:norlund-egf}, this is precisely
$\beta_{n-1}^{(n)}=(-1)^{n-1}(n-1)!$.

Conversely, $f(x)=\log(1+x)$ and $g(y)=e^y-1$ give
\begin{equation}\label{eq:cauchy-lagrange-diagonal}
 \boxed{\displaystyle
 [x^{n-1}]\left(\frac{x}{\log(1+x)}\right)^n
 =\frac1{(n-1)!}.}
\end{equation}
Here the left side divided by $n$ equals
$[y^n](e^y-1)=1/n!$.  This evaluates a diagonal of the higher powers of the
Cauchy-number generating function from
\cref{eq:bernoulli-second-kind-egf}.

\section{Kepler's equation and the Fourier--Bessel resummation}
\index{Kepler equation}
A classical analytic application of reversion is Kepler's equation
\begin{equation}\label{eq:kepler-equation}
 E=M+e\sin E,\qquad 0\le e<1.
\end{equation}
Taking $x=M$, $y=e$, and $f(u)=\sin u$ in
\eqref{eq:shifted-lagrange-v} gives the convergent reversion series
\begin{equation}\label{eq:kepler-reversion}
 E=M+\sum_{k\ge1}\frac{e^k}{k!}
 \left(\frac{\dd}{\dd M}\right)^{k-1}\sin^kM.
\end{equation}
It admits a more useful harmonic resummation.

\begin{theorem}[Fourier--Bessel solution of Kepler's equation]
\label{thm:kepler-bessel}
For $0\le e<1$, the unique real solution of
\eqref{eq:kepler-equation} is
\begin{equation}\label{eq:kepler-bessel}
 E=M+2\sum_{n\ge1}\frac{J_n(ne)}{n}\sin(nM),
\end{equation}
where
\begin{equation}\label{eq:bessel-integral}
 J_n(t)=\frac1{2\pi}\int_0^{2\pi}
 \cos(nu-t\sin u)\,\dd u.
\end{equation}
The series converges absolutely and uniformly in $M$ for every fixed $e<1$,
and uniformly for $e$ in compact subsets of $[0,1)$.
\end{theorem}
\begin{proof}
The map $E\mapsto M=E-e\sin E$ has derivative
$1-e\cos E\ge1-e>0$, so it is strictly increasing and has a unique analytic
inverse satisfying $E(M+2\pi)=E(M)+2\pi$.  Thus
$q(M)=E(M)-M$ is a real-analytic, odd, $2\pi$-periodic function.  Its Fourier
series contains only sine terms and converges absolutely and uniformly.

Let $c_n$ be its $n$th sine coefficient.  Since $q=e\sin E$ and
$\dd M=(1-e\cos E)\dd E$,
\begin{align*}
 c_n
 &=\frac1\pi\int_0^{2\pi}q(M)\sin(nM)\,\dd M\\
 &=\frac e\pi\int_0^{2\pi}
 \sin E\,\sin\bigl(n(E-e\sin E)\bigr)(1-e\cos E)\,\dd E.
\end{align*}
Now
\[
 \frac{\dd}{\dd E}\cos\bigl(n(E-e\sin E)\bigr)
 =-n(1-e\cos E)\sin\bigl(n(E-e\sin E)\bigr).
\]
Integration by parts, with vanishing boundary term, gives
\begin{equation}\label{eq:kepler-cn-integral}
 c_n=\frac e{\pi n}\int_0^{2\pi}
 \cos E\cos(nE-ne\sin E)\,\dd E.
\end{equation}
From \eqref{eq:bessel-integral} and
$2\cos E\cos\theta=\cos(\theta+E)+\cos(\theta-E)$,
this integral is
$\pi(J_{n-1}(ne)+J_{n+1}(ne))$.  The recurrence
\begin{equation}\label{eq:bessel-recurrence}
 J_{n-1}(t)+J_{n+1}(t)=\frac{2n}{t}J_n(t)
\end{equation}
follows directly from the integral: integrate
$\frac{\dd}{\dd u}\sin(nu-t\sin u)$ over a period to obtain
\[
 n\int_0^{2\pi}\cos(nu-t\sin u)\,\dd u
 =t\int_0^{2\pi}\cos u\cos(nu-t\sin u)\,\dd u.
\]
Substitution in \eqref{eq:kepler-cn-integral} yields
$c_n=2J_n(ne)/n$, proving \eqref{eq:kepler-bessel}.  Uniform convergence on
compact $e$-intervals also follows directly from analytic dependence of the
periodic inverse, or from standard contour shifts in its Fourier coefficients.
\end{proof}

Expanding the exponential in the integral gives the familiar entire series
\begin{equation}\label{eq:bessel-power-series}
 J_n(t)=\sum_{r\ge0}\frac{(-1)^r}{r!(n+r)!}
 \left(\frac t2\right)^{n+2r}\qquad(n\ge0).
\end{equation}
One proof is to extract the constant Fourier coefficient of
$e^{inu}\exp\{\frac t2(e^{-iu}-e^{iu})\}$; matching equal powers of
$e^{iu}$ and $e^{-iu}$ gives exactly the displayed sum.  Hence every term in
\eqref{eq:kepler-bessel} can itself be expanded into the perturbation series
\eqref{eq:kepler-reversion}; the two formulas are different reorganizations of
the same analytic inverse.

\chapter{Lagrange's reversion theorem in shifted form}

\section{The two-parameter implicit equation}
Let $v=v(x,y)$ be the unique formal solution of
\begin{equation}\label{eq:shifted-reversion-equation}
 v=x+yf(v),
 \qquad v(x,0)=x.
\end{equation}

\begin{theorem}[Lagrange reversion]\label{thm:shifted-lagrange-reversion}
For every formal or analytic $g$,
\begin{equation}\label{eq:shifted-lagrange-reversion}
 g(v)=g(x)+\sum_{k=1}^{\infty}\frac{y^k}{k!}
 \left(\frac{\partial}{\partial x}\right)^{k-1}
 \bigl(f(x)^kg'(x)\bigr).
\end{equation}
In particular,
\begin{equation}\label{eq:shifted-lagrange-v}
 v=x+\sum_{k=1}^{\infty}\frac{y^k}{k!}
 \left(\frac{\partial}{\partial x}\right)^{k-1}f(x)^k.
\end{equation}
\end{theorem}
\begin{proof}
Put $u=v-x$.  Then $u=y\phi(u)$ with $\phi(u)=f(x+u)$, treating $x$ as a
coefficient parameter.  Apply \eqref{eq:lagrange-burmann} to
$H(u)=g(x+u)$:
\[
 [y^k]g(v)=\frac1k[u^{k-1}]g'(x+u)f(x+u)^k
 =\frac1{k!}D_x^{k-1}(g'(x)f(x)^k).
\]
Summation proves the first formula; take $g(z)=z$ for the second.
\end{proof}

\section{The delta-function proof decoded}
The archive also gives a distributional derivation.  Assume temporarily that
$x,y\in\R$, the equation $z-x-yf(z)=0$ has one simple root $v$, and
$1-yf'(v)>0$.  The elementary change-of-variables rule for the Dirac delta gives
\begin{equation}\label{eq:delta-root}
 g(v)=\int_{\R}\delta(yf(z)-z+x)g(z)(1-yf'(z))\,\dd z.
\end{equation}
Using
$\delta(s)=\int_{\R}e^{iks}\dd k/(2\pi)$ and expanding only in the formal
parameter $y$,
\begin{align}
 g(v)
 &=\sum_{n\ge0}D_x^n\left[
 \frac{y^nf(x)^ng(x)}{n!}(1-yf'(x))\right]\notag\\
 &=\sum_{n\ge0}D_x^n\left[
 \frac{y^nf(x)^ng(x)}{n!}
 -\frac{y^{n+1}}{(n+1)!}\bigl((gf^{n+1})'-g'f^{n+1}\bigr)
 \right].
 \label{eq:delta-reversion-middle}
\end{align}
The first term of level $n+1$ cancels the total derivative in the second term of
level $n$; what remains is exactly
\eqref{eq:shifted-lagrange-reversion}.  Thus the Fourier--delta calculation is a
compact analytic avatar of the formal residue proof.  Formula
\eqref{eq:delta-root} requires an absolute Jacobian in general; the displayed sign
is valid in the stated orientation-preserving neighbourhood.


\chapter{Multivariate Lagrange--Good inversion}
\index{Lagrange inversion!multivariate}\index{Good inversion formula}
One variable hides the Jacobian in an integration-by-parts identity.  In
several variables the Jacobian survives as a determinant and is the essential
new feature.  The determinant form developed here is Good's multivariate
extension \cite{Good1960}; its combinatorial meaning is illuminated by Gessel's
tree and cancellation methods \cite{Gessel1987}.

\section{Multivariate formal residues}
For a Laurent series in $d$ variables, define
\[
 \res_{\boldsymbol x}H(\boldsymbol x)\,
 \dd x_1\cdots\dd x_d
 =[x_1^{-1}\cdots x_d^{-1}]H(\boldsymbol x).
\]

\begin{lemma}[Formal multivariate change of variables]
\label{lem:multivariate-residue-substitution}
Let
$f_i(\boldsymbol x)=x_i u_i(\boldsymbol x)$, where every
$u_i(\boldsymbol0)$ is invertible.  Then
\begin{multline}\label{eq:multivariate-residue-substitution}
 \res_{\boldsymbol t}H(\boldsymbol t)\,
 \dd t_1\cdots\dd t_d\\
 =\res_{\boldsymbol x}H(\boldsymbol f(\boldsymbol x))
 \det\!\left(\frac{\partial f_i}{\partial x_j}\right)_{i,j=1}^{d}
 \dd x_1\cdots\dd x_d.
\end{multline}
\end{lemma}
\begin{proof}
By linearity it is enough to consider the differential form
$\boldsymbol t^{\boldsymbol m}\dd t_1\wedge\cdots\wedge\dd t_d$.
If some $m_j\ne-1$, that form is exact: up to the sign needed to move
$\dd t_j$ into position, it is
\[
 \dd\!\left(
 \frac{t_j^{m_j+1}}{m_j+1}
 \prod_{i\ne j}t_i^{m_i}
 \dd t_1\wedge\cdots\widehat{\dd t_j}\cdots\wedge\dd t_d
 \right).
\]
The residue of an exact formal top form is zero, because it is a coefficient
of a formal partial derivative.  Pullback commutes with exterior
differentiation, so the residue after substitution is also zero.

It remains to consider $m_1=\cdots=m_d=-1$.  Since
$f_i=x_i u_i$,
\[
 \frac{\dd f_i}{f_i}=\frac{\dd x_i}{x_i}+\frac{\dd u_i}{u_i}.
\]
In the wedge product of these $d$ forms, the product of all
$\dd x_i/x_i$ has residue $1$.  Every other term contains at least one
regular form $\dd u_i/u_i$, whose coefficients have no negative exponents;
such a term cannot contain $x_1^{-1}\cdots x_d^{-1}$.  Its residue is zero.
Thus the substitution preserves the unique nonexact Laurent top form and
annihilates all the others, proving the formula.
\end{proof}

\section{Good's determinant formula}
Let $\boldsymbol w=(w_1,\ldots,w_d)$ be the unique vector of formal series
satisfying
\begin{equation}\label{eq:good-implicit-system}
 w_i=t_i\phi_i(\boldsymbol w),\qquad
 \phi_i(\boldsymbol0)\ne0,
 \qquad 1\le i\le d.
\end{equation}
Existence and uniqueness follow coefficient by coefficient: the constant term
is zero, the coefficient of $t_i$ is $\phi_i(0)$, and every coefficient of
total degree $N$ depends only on lower total degrees.

\begin{theorem}[Lagrange--Good inversion]
\label{thm:lagrange-good}
For a multi-index $\boldsymbol n=(n_1,\ldots,n_d)\in\N^d$ and any formal
series $F$,
\begin{multline}\label{eq:lagrange-good}
 [\boldsymbol t^{\boldsymbol n}]F(\boldsymbol w(\boldsymbol t))\\
 =[\boldsymbol x^{\boldsymbol n}]
 F(\boldsymbol x)\prod_{i=1}^{d}\phi_i(\boldsymbol x)^{n_i}
 \det\!\left(
 \delta_{ij}-\frac{x_j}{\phi_i(\boldsymbol x)}
 \frac{\partial\phi_i}{\partial x_j}(\boldsymbol x)
 \right)_{i,j=1}^{d}.
\end{multline}
\end{theorem}
\begin{proof}
Represent the coefficient on the left as an iterated residue and make the
substitution
\[
 t_i=\frac{x_i}{\phi_i(\boldsymbol x)}.
\]
Under this substitution, the solution $\boldsymbol w(\boldsymbol t)$ becomes
$\boldsymbol x$.  The Jacobian entries are
\[
 \frac{\partial t_i}{\partial x_j}
 =\frac1{\phi_i}\left(
 \delta_{ij}-\frac{x_i}{\phi_i}
 \frac{\partial\phi_i}{\partial x_j}
 \right).
\]
The factors $t_i^{-n_i-1}$ contribute
$x_i^{-n_i-1}\phi_i^{n_i+1}$, while factoring $1/\phi_i$ from each Jacobian
row leaves $\phi_i^{n_i}$.  We therefore obtain the coefficient on the right
with determinant
\[
 \det\!\left(
 \delta_{ij}-\frac{x_i}{\phi_i}\frac{\partial\phi_i}{\partial x_j}
 \right).
\]
Let $D_x=\diag(x_1,\ldots,x_d)$,
$D_\phi=\diag(\phi_1,\ldots,\phi_d)$, and
$J_\phi=(\partial\phi_i/\partial x_j)$.  The last matrix is
$I-D_xD_\phi^{-1}J_\phi$.  The identity
$\det(I-AB)=\det(I-BA)$ changes its determinant to
$\det(I-D_\phi^{-1}J_\phi D_x)$, whose $(i,j)$ entry is the one in
\eqref{eq:lagrange-good}.  The residue substitution is justified by
\cref{lem:multivariate-residue-substitution}.
\end{proof}

\begin{corollary}[Recovery of one-variable Lagrange--B\"urmann]
For $d=1$, \eqref{eq:lagrange-good} reduces to
\[
 [t^n]F(w(t))=[x^n]F(x)\phi(x)^n
 -[x^{n-1}]F(x)\phi(x)^{n-1}\phi'(x).
\]
This equals
\begin{equation}\label{eq:good-one-variable-reduction}
 \frac1n[x^{n-1}]F'(x)\phi(x)^n.
\end{equation}
\end{corollary}
\begin{proof}
Differentiate $F(x)\phi(x)^n$ and take the coefficient of $x^{n-1}$:
\[
 n[x^n]F\phi^n
 =[x^{n-1}]F'\phi^n
 +n[x^{n-1}]F\phi^{n-1}\phi'.
\]
Rearrangement gives \eqref{eq:good-one-variable-reduction}, which is
\eqref{eq:lagrange-burmann}.
\end{proof}

\section{A multitype tree specialization}
Let $A=(a_{ij})$ be a $d\times d$ matrix of weights and consider
\begin{equation}\label{eq:multitype-tree-system}
 w_i=t_i\exp\!\left(\sum_{j=1}^{d}a_{ij}w_j\right).
\end{equation}
This is the exponential generating system for rooted labelled trees with
vertex types: a type-$i$ root carries an unordered set of type-$j$ subtrees,
with edge weight $a_{ij}$.  Good's formula gives, for every $F$,
\begin{multline}\label{eq:multitype-tree-good}
 [\boldsymbol t^{\boldsymbol n}]F(\boldsymbol w)
 =[\boldsymbol x^{\boldsymbol n}]F(\boldsymbol x)
 \exp\!\left(\sum_{i,j}n_i a_{ij}x_j\right)\\
 {}\times\det\!\left(\delta_{ij}-a_{ij}x_j\right)_{i,j=1}^{d}.
\end{multline}
Indeed, $\partial_j\phi_i=a_{ij}\phi_i$, so the logarithmic Jacobian in
\eqref{eq:lagrange-good} collapses to the displayed linear determinant.  This
example shows why a determinant, rather than an unwieldy alternating sum, is
the natural multivariate correction.

\section{Analytic convergence and critical sets}
Formal inversion determines all coefficients without convergence assumptions.
For analytic maps, the inverse centered at a regular point continues until it
meets either a singularity of the defining map or the image of the critical
set where the Jacobian determinant vanishes.  In one variable this condition is
$f'(w)=0$; in the system \eqref{eq:good-implicit-system} it is precisely the
vanishing of the determinant in \eqref{eq:lagrange-good}.  Simple critical
points generically create square-root Puiseux singularities and the familiar
$n^{-3/2}$ coefficient factor.  The branch point $-e^{-1}$ of Lambert $W$, the
point $1/4$ of the Catalan generating function, and the points $\pm1$ of
$\arcsin$ are the one-variable manifestations of this same geometric
criterion.

\chapter{Laplace-type asymptotic integrals}
\index{Laplace asymptotics}

\section{Endpoint expansion}
Consider
\begin{equation}\label{eq:laplace-integral}
 I(\lambda)=\int_a^b e^{-\lambda f(x)}g(x)\,\dd x,
 \qquad \lambda\to+\infty,
\end{equation}
where the endpoint $a$ is the unique relevant minimum and, with $t=x-a$,
\begin{align}
 f(a+t)&\sim f(a)+\sum_{k\ge0}a_kt^{k+\alpha},
 &a_0>0,\quad\alpha>0,
 \label{eq:laplace-f-expansion}\\
 g(a+t)&\sim\sum_{k\ge0}b_kt^{k+\beta-1},
 &\beta>0.
 \label{eq:laplace-g-expansion}
\end{align}
Let $\OBellP{k}{j}$ denote the ordinary Bell polynomial.

\begin{theorem}[Laplace--Erd\'elyi coefficient formula]
\label{thm:laplace-bell}
Under the standard differentiability and remainder hypotheses that permit Watson's
lemma,
\begin{equation}\label{eq:laplace-asymptotic}
 I(\lambda)\sim e^{-\lambda f(a)}
 \sum_{n=0}^{\infty}
 \Gamma\!\left(\frac{n+\beta}{\alpha}\right)
 \frac{c_n}{\lambda^{(n+\beta)/\alpha}},
\end{equation}
where
\begin{equation}\label{eq:laplace-cn}
 c_n=\frac1{\alpha a_0^{(n+\beta)/\alpha}}
 \sum_{k=0}^{n}b_{n-k}
 \sum_{j=0}^{k}
 \binom{-\frac{n+\beta}{\alpha}}{j}
 \frac{\OBellP{k}{j}(a_1,\ldots,a_{k-j+1})}{a_0^j}.
\end{equation}
\end{theorem}
\begin{proof}
Write
$f(a+t)-f(a)=t^\alpha A(t)$ with
$A(t)=a_0+a_1t+\cdots$, and set
$s=(f(a+t)-f(a))^{1/\alpha}=tA(t)^{1/\alpha}$.  This has an inverse $t=t(s)$.
Expand
\[
 g(a+t(s))t'(s)=\sum_{n\ge0}d_ns^{n+\beta-1}.
\]
A generalized Lagrange--Bürmann calculation, first for integral $\beta$ and then
by formal continuation in the exponent, gives
\[
 d_n=\sum_{k=0}^{n}b_{n-k}[t^k]A(t)^{-(n+\beta)/\alpha}.
\]
Indeed, apply \eqref{eq:lagrange-burmann} to an antiderivative of
$t^{m+\beta-1}$ and differentiate back after substitution.

Put $\rho=(n+\beta)/\alpha$.  The ordinary Bell expansion yields
\begin{align*}
 [t^k]A(t)^{-\rho}
 &=a_0^{-\rho}[t^k]
 \sum_{j\ge0}\binom{-\rho}{j}
 \left(\sum_{r\ge1}\frac{a_r}{a_0}t^r\right)^j\\
 &=a_0^{-\rho}\sum_{j=0}^{k}\binom{-\rho}{j}a_0^{-j}
 \OBellP{k}{j}(a_1,\ldots,a_{k-j+1}).
\end{align*}
Thus $c_n=d_n/\alpha$ is \eqref{eq:laplace-cn}.  Finally
\[
 \int_0^\infty e^{-\lambda s^\alpha}s^{n+\beta-1}\dd s
 =\frac1\alpha\Gamma\!\left(\frac{n+\beta}{\alpha}\right)
 \lambda^{-(n+\beta)/\alpha},
\]
and Watson's lemma controls the replacement of the finite endpoint interval by
the asymptotic expansion.
\end{proof}

\begin{remark}
The appearance of ordinary rather than exponential Bell polynomials is forced by
the ordinary Taylor expansion of $A(t)^{-\rho}$.  Lagrange inversion enters in the
change from $t$ to the natural Laplace coordinate $s$; the gamma factor enters only
in the final monomial integral.
\end{remark}


\part{M\"obius, Sheffer, and Riordan Synthesis}

\chapter{M\"obius inversion on the partition lattice}
\index{Mobius inversion@M\"obius inversion!partition lattice}
Let $\Pi_n$ be the lattice of partitions of $[n]$, ordered by refinement.
The moment--cumulant formulas already exhibit its zeta transform; here we make
the inverse explicit.

\begin{theorem}[M\"obius function of the partition lattice]
\label{thm:partition-lattice-mobius}
For $\pi\le\sigma$ in $\Pi_n$, let $m_C$ be the number of blocks of $\pi$
contained in a block $C$ of $\sigma$.  Then
\begin{equation}\label{eq:partition-lattice-mobius-general}
 \mu(\pi,\sigma)=\prod_{C\in\sigma}
 (-1)^{m_C-1}(m_C-1)!.
\end{equation}
In particular,
\begin{equation}\label{eq:partition-lattice-mobius-top}
 \mu(\pi,\widehat1)=(-1)^{|\pi|-1}(|\pi|-1)!.
\end{equation}
\end{theorem}
\begin{proof}
The interval $[\pi,\sigma]$ is the direct product, over $C\in\sigma$, of the
partition lattices on the $m_C$ blocks of $\pi$ lying in $C$.  M\"obius
functions multiply on direct products, so it is enough to prove
$\mu_{\Pi_m}(\widehat0,\widehat1)=(-1)^{m-1}(m-1)!$ for every $m$.

We prove this by induction on $m$.  If $\tau<\widehat1$ is a partition of
$[m]$, then $[\widehat0,\tau]$ is the product of the smaller lattices
$\Pi_{|B|}$ over $B\in\tau$.  The induction hypothesis and multiplicativity
therefore give
\[
 \mu_{\Pi_m}(\widehat0,\tau)
 =\prod_{B\in\tau}(-1)^{|B|-1}(|B|-1)!.
\]
For a fixed $\tau$, the absolute value on the right counts choices of a cyclic
ordering on every block, hence permutations of $[m]$ whose cycle-support
partition is $\tau$; its sign $(-1)^{m-|\tau|}$ is the sign of each such
permutation.  Thus the sum of the displayed candidates over all
$\tau\in\Pi_m$ is the sum of the signs of all permutations in $S_m$.  For
$m>1$ this sum is zero, since multiplication by $(1\ 2)$ is a sign-reversing
involution.  The defining recurrence
$\sum_{\tau\le\widehat1}\mu(\widehat0,\tau)=0$ therefore forces the remaining
top value to be the candidate
$(-1)^{m-1}(m-1)!$.  The case $m=1$ is immediate.  Multiplicativity on
$[\pi,\sigma]$ now proves \eqref{eq:partition-lattice-mobius-general}.
\end{proof}

\begin{corollary}[Moments and cumulants as M\"obius inverses]
\label{cor:moment-cumulant-mobius}
If
\begin{equation}\label{eq:moment-partition-lattice}
 m_n=\sum_{\pi\in\Pi_n}\prod_{B\in\pi}\kappa_{|B|},
\end{equation}
then
\begin{equation}\label{eq:cumulant-partition-lattice}
 \kappa_n=\sum_{\pi\in\Pi_n}
 (-1)^{|\pi|-1}(|\pi|-1)!
 \prod_{B\in\pi}m_{|B|}.
\end{equation}
\end{corollary}
\begin{proof}
Equation \eqref{eq:moment-partition-lattice} is the zeta transform on
$\Pi_n$ of the multiplicative block weight $\kappa$.  Apply M\"obius inversion
and use \eqref{eq:partition-lattice-mobius-top}.  Grouping partitions by their
block-size profiles gives exactly
\eqref{eq:moments-from-cumulants} and
\eqref{eq:cumulants-from-moments}.
\end{proof}

The Boolean lattice controls inclusion--exclusion, while the partition lattice
controls passage between arbitrary structures and connected components.  Their
M\"obius functions explain why signs and factorials occur in the inverse
binomial and inverse Bell transforms.

\chapter{Sheffer sequences and delta operators}
\index{Sheffer sequence}\index{delta operator}
The binomial-type families of \cref{thm:binomial-type-bell} are the associated
sequences of the wider Sheffer class.  We use the exponential convention of
Roman \cite{Roman1984}; the group-theoretic identification with Riordan arrays
is developed in \cite{HeHsuShiue2007}.

\begin{definition}
Let $A(z)$ be a unit, $A(0)\ne0$, and let
$\beta(z)=\beta_1z+\beta_2z^2+\cdots$ with $\beta_1\ne0$.  The Sheffer sequence
for $(A,\beta)$ is defined by
\begin{equation}\label{eq:sheffer-egf}
 \sum_{n\ge0}s_n(x)\frac{z^n}{n!}
 =A(z)e^{x\beta(z)}.
\end{equation}
The associated binomial-type sequence $p_n(x)$ has generating function
$e^{x\beta(z)}$.
\end{definition}

\begin{theorem}[Sheffer identity and characterization]
\label{thm:sheffer-characterization}
The polynomials in \eqref{eq:sheffer-egf} satisfy
\begin{equation}\label{eq:sheffer-identity}
 s_n(x+y)=\sum_{k=0}^{n}\binom nk s_k(x)p_{n-k}(y).
\end{equation}
Conversely, if $p_n$ is of binomial type and a polynomial sequence $s_n$
satisfies \eqref{eq:sheffer-identity}, then its EGF has the form
\eqref{eq:sheffer-egf} for a unique unit $A(z)$.
\end{theorem}
\begin{proof}
Multiply the EGF $A(z)e^{x\beta(z)}$ by
$e^{y\beta(z)}$ and compare coefficients to prove the identity.  Conversely,
let $S(x,z)=\sum s_n(x)z^n/n!$ and
$P(y,z)=\sum p_n(y)z^n/n!=e^{y\beta(z)}$.  The assumed identity is
$S(x+y,z)=S(x,z)P(y,z)$.  Setting $x=0$ gives
$S(y,z)=S(0,z)e^{y\beta(z)}$; hence $A(z)=S(0,z)$, uniquely.
\end{proof}

Writing $A(z)=\sum a_jz^j/j!$ immediately gives
\begin{equation}\label{eq:sheffer-convolution}
 s_n(x)=\sum_{j=0}^{n}\binom nj a_jp_{n-j}(x).
\end{equation}
Thus the unit $A$ performs an invertible Appell convolution on an associated
binomial-type family.

Let $\bar\beta$ be the compositional inverse of $\beta$.  Functional calculus
in the differentiation operator gives the delta operator
\begin{equation}\label{eq:sheffer-delta-operator}
 Q=\bar\beta(D_x).
\end{equation}

\begin{theorem}[Lowering law]
\label{thm:sheffer-lowering}
Both the associated and the Sheffer sequences satisfy
\begin{equation}\label{eq:sheffer-lowering}
 Qp_n(x)=np_{n-1}(x),\qquad
 Qs_n(x)=ns_{n-1}(x).
\end{equation}
\end{theorem}
\begin{proof}
Since $D_xe^{x\beta(z)}=\beta(z)e^{x\beta(z)}$,
\[
 \bar\beta(D_x)e^{x\beta(z)}
 =\bar\beta(\beta(z))e^{x\beta(z)}=ze^{x\beta(z)}.
\]
The factor $A(z)$ is independent of $x$, so the same equation holds for the
Sheffer EGF.  Coefficient comparison gives both lowering laws.
\end{proof}

Lagrange inversion is therefore internal to Sheffer calculus: it computes the
coefficients of $\bar\beta$, and hence the differential operator that lowers
the sequence.

\section{Principal examples}
The following table uses exactly one convention, so the comparisons are
literal rather than up to rescaling.
\begin{center}
\begin{tabularx}{\textwidth}{>{\raggedright\arraybackslash}p{0.24\textwidth}
 >{\raggedright\arraybackslash}p{0.25\textwidth}X}
\toprule
Family & $(A(z),\beta(z))$ & EGF and lowering operator\\
\midrule
Monomials & $(1,z)$ & $e^{xz}$; $Q=D$.\\
Falling factorials & $(1,\log(1+z))$ & $(1+z)^x$;
$Q=e^D-1=\Delta$.\\
Rising factorials & $(1,-\log(1-z))$ & $(1-z)^{-x}$;
$Q=1-e^{-D}$.\\
Touchard polynomials & $(1,e^z-1)$ & $e^{x(e^z-1)}$;
$Q=\log(1+D)$.\\
Bernoulli polynomials & $(z/(e^z-1),z)$ &
$ze^{xz}/(e^z-1)$; $Q=D$.\\
Euler polynomials & $(2/(e^z+1),z)$ &
$2e^{xz}/(e^z+1)$; $Q=D$.\\
Bernoulli polynomials of the second kind &
$(z/\log(1+z),\log(1+z))$ &
$z(1+z)^x/\log(1+z)$; $Q=e^D-1$.\\
\bottomrule
\end{tabularx}
\end{center}
The first four rows are of binomial type.  The last three have nontrivial
prefactors and are genuine Sheffer sequences; the two Bernoulli families are
adapted respectively to differentiation and finite differences.

\chapter{Exponential Riordan arrays}
\index{Riordan array!exponential}

\section{Definition, action, and group law}
For a unit $g(z)$ and a delta series $f(z)$, define the lower-triangular array
\begin{equation}\label{eq:riordan-entry}
 R_{n,k}=\frac{n!}{k!}[z^n]g(z)f(z)^k.
\end{equation}
If $a=(a_k)$ has EGF $A(z)$, then the transformed sequence
$b_n=\sum_kR_{n,k}a_k$ has EGF
\begin{equation}\label{eq:riordan-action}
 B(z)=g(z)A(f(z)).
\end{equation}
Indeed, summing \eqref{eq:riordan-entry} against $a_k$ gives
$g(z)\sum_ka_kf(z)^k/k!$.

\begin{theorem}[Riordan group law]
\label{thm:riordan-group}
With arrays denoted by pairs,
\begin{align}
 (g,f)(h,\ell)&=(g\,h\!\circ f,\ \ell\!\circ f),
 \label{eq:riordan-product}\\
 (g,f)^{-1}&=\left(\frac1{g\circ\bar f},\bar f\right),
 \label{eq:riordan-inverse}
\end{align}
where $\bar f=f^{\langle-1\rangle}$.
\end{theorem}
\begin{proof}
Apply first $(h,\ell)$ and then $(g,f)$ to an arbitrary EGF $A$:
\[
 A(z)\longmapsto h(z)A(\ell(z))
 \longmapsto g(z)h(f(z))A(\ell(f(z))),
\]
which proves the product.  Substitution of
$(1/(g\circ\bar f),\bar f)$ on either side gives $(1,z)$, the identity array.
\end{proof}

The row polynomials of $(A,\beta)$ are precisely the Sheffer sequence:
\begin{equation}\label{eq:riordan-sheffer-row}
 s_n(x)=\sum_{k=0}^{n}R_{n,k}x^k,
 \qquad
 \sum_ns_n(x)\frac{z^n}{n!}=A(z)e^{x\beta(z)}.
\end{equation}
Thus Sheffer composition is ordinary matrix multiplication in the exponential
Riordan group.

\section{Bell and Stirling arrays}
The partial Bell-polynomial array is the Riordan array
\begin{equation}\label{eq:bell-riordan-array}
 \BellP nk(x_1,x_2,\ldots)
 =\frac{n!}{k!}[z^n]
 \left(\sum_{j\ge1}x_j\frac{z^j}{j!}\right)^k.
\end{equation}
The Stirling arrays are
\begin{align}
 \stirtwo nk
 &=\frac{n!}{k!}[z^n](e^z-1)^k,
 \label{eq:stirling-second-riordan}\\
 s(n,k)
 &=\frac{n!}{k!}[z^n]\log(1+z)^k.
 \label{eq:stirling-first-riordan}
\end{align}
Since $e^z-1$ and $\log(1+z)$ are compositional inverses, their Riordan arrays
are inverse matrices.  This recovers Stirling inversion as a special case of
\eqref{eq:riordan-inverse}; Fa\`a di Bruno composition is the corresponding
matrix product, and Lagrange inversion computes its inverse entries.

\begin{theorem}[Connection constants for two Sheffer families]
\label{thm:sheffer-connection-constants}
Let $s_n$ and $r_n$ be Sheffer for $(A,\beta)$ and $(C,\gamma)$, respectively,
and let $\bar\gamma$ be the inverse of $\gamma$.  In
\begin{equation}\label{eq:sheffer-connection-expansion}
 s_n(x)=\sum_{k=0}^{n}c_{n,k}r_k(x),
\end{equation}
the connection coefficients are
\begin{equation}\label{eq:sheffer-connection-coefficients}
 c_{n,k}=\frac{n!}{k!}[z^n]
 \frac{A(z)}{C(\bar\gamma(\beta(z)))}
 \bigl(\bar\gamma(\beta(z))\bigr)^k.
\end{equation}
\end{theorem}
\begin{proof}
In vector notation, $s=(A,\beta)x^{\bullet}$ and
$r=(C,\gamma)x^{\bullet}$.  Hence
$s=(A,\beta)(C,\gamma)^{-1}r$.  Use
\eqref{eq:riordan-product} and \eqref{eq:riordan-inverse}; the resulting
Riordan pair is
\[
 \left(
 \frac{A(z)}{C(\bar\gamma(\beta(z)))},\,
 \bar\gamma(\beta(z))
 \right).
\]
Its entry formula is exactly
\eqref{eq:sheffer-connection-coefficients}.
\end{proof}

This theorem systematically produces all change-of-basis triangles among
monomials, factorials, Touchard, Bernoulli, Euler, and Cauchy polynomials.  It
also explains why apparently unrelated coefficient tables in the five drafts
obey the same matrix laws.

\section{Abel polynomials from tree reversion}
For a parameter $a$, let $T_a(z)$ be the unique solution of
\begin{equation}\label{eq:tree-function-a}
 T_a(z)=ze^{aT_a(z)}.
\end{equation}
Lagrange inversion gives
\begin{equation}\label{eq:abel-exponential-coefficient}
 [z^n]e^{xT_a(z)}
 =\frac{x}{n}[u^{n-1}]e^{(x+an)u}
 =\frac{x(x+an)^{n-1}}{n!}.
\end{equation}
Replacing $a$ by $-a$ shows that the Abel polynomials
\begin{equation}\label{eq:abel-polynomial}
 A_n(x)=x(x-an)^{n-1},\qquad A_0(x)=1,
\end{equation}
are of binomial type, with EGF
$\exp(xT_{-a}(z))$.  Therefore
\begin{multline}\label{eq:abel-binomial-identity}
 (x+y)(x+y-an)^{n-1}\\
 =\sum_{k=0}^{n}\binom nk
 x(x-ak)^{k-1}y\bigl(y-a(n-k)\bigr)^{n-k-1},
\end{multline}
with the endpoint factors interpreted by $A_0=1$.  This is Abel's binomial
identity, obtained without guessing: it is the binomial law forced by the tree
reversion series.

\appendix
\part{Synthesis, Verification, and Source Record}

\chapter{Editorial synthesis of the five drafts}
\label{app:editorial-synthesis}

\section{Method of consolidation}
The five draft manuscripts are not independent treatments of five disjoint
subjects.  They are overlapping realizations of one coefficient calculus, with
substantially different chapter orders, proof granularity, notation, and choices
of examples.  A literal concatenation would therefore repeat the elementary
Stirling and Bell theory several times while separating formulas that are best
understood as instances of one theorem.  The consolidation used the following
rules.
\begin{enumerate}
\item An identity appearing in more than one draft is stated once, in the
      strongest useful form, and is followed by one complete proof.  Alternative
      arguments are retained when they expose genuinely different mechanisms,
      such as a bijection versus a formal-residue calculation.
\item Definitions are moved before every theorem that depends on them.  Formal
      power-series statements are separated from analytic statements, so that
      convergence assumptions are neither silently imposed nor silently omitted.
\item Notational collisions are resolved globally.  In particular, Bell numbers,
      Bernoulli numbers, Bell polynomials, Euler polynomials, Eulerian numbers,
      type-$B$ Eulerian numbers, and binary-tree numbers are never denoted by the
      same undecorated symbol in the same context.
\item Schematic formulas are replaced by checkable formulas.  Missing index
      ranges, boundary values, Jacobians, branch choices, and normalization factors
      are made explicit.  Where an archived proof stopped after suggesting a
      pattern, the argument is completed by coefficient extraction, M\"obius
      inversion, a cycle lemma, or formal residues.
\item Material in the supplementary Wikipedia and MathWorld folders is not copied
      as a second encyclopedia.  It is absorbed into the existing dependency
      structure and proved using the same coefficient machinery.
\end{enumerate}

\section{Draft-by-draft concordance}
\renewcommand{\arraystretch}{1.13}
\begin{longtable}{@{}>{\raggedright\arraybackslash}p{0.24\textwidth}
                      >{\raggedright\arraybackslash}p{0.31\textwidth}
                      >{\raggedright\arraybackslash}p{0.39\textwidth}@{}}
\toprule
\textbf{Draft} & \textbf{Principal emphasis} & \textbf{Material retained in the unified text}\\
\midrule
\endfirsthead
\toprule
\textbf{Draft} & \textbf{Principal emphasis} & \textbf{Material retained in the unified text}\\
\midrule
\endhead
\midrule
\multicolumn{3}{r}{\small Continued on the next page}\\
\endfoot
\bottomrule
\endlastfoot

\nolinkurl{combinatorial_coefficient_calculus.tex}
& The most developed dependency-driven treatment; strong coverage of Stirling,
  Bell, Eulerian, Fa\`a di Bruno, Lagrange, congruence, and asymptotic methods.
& Used as the structural spine.  Its formal coefficient calculus, full triangular
  array theory, exact Bell-number analysis, Eulerian chapters, residue proof of
  Lagrange inversion, and Laplace coefficient formula are retained and expanded.\\[1ex]

\nolinkurl{Formulae_Compendium.tex}
& A compact but wide survey, with particularly useful higher-derivative
  specializations and systematic reversion examples.
& Supplies the expanded library of power, logarithm, reciprocal, trigonometric,
  product, quotient, Fr\'echet, and inverse-function derivative formulas, together
  with normalization cross-checks and the Kepler reversion example.\\[1ex]

\nolinkurl{Formulae_Monograph.tex} in
\nolinkurl{Formulae_Monograph_Package}
& A proof-oriented monograph emphasizing transitions among factorial bases,
  labelled structures, special-function series, and certificate procedures.
& Its alternative combinatorial proofs, generalized factorial triangles,
  finite-difference identities, special-function viewpoints, and exact
  low-order consistency checks are merged into the corresponding chapters.\\[1ex]

\nolinkurl{Formulae_Monograph.tex} in
\nolinkurl{Formulae_Monograph_Source_and_PDF}
& A parallel monograph with a different order and additional attention to
  normalization, multivariate differentiation, and auditability.
& Contributes the global notation policy, several repaired boundary conventions,
  the coordinate-free composition viewpoint, the formula-dependency map, and the
  source-audit methodology used in this appendix.\\[1ex]

\nolinkurl{Combinatorial_Formulae_and_Inversion_Theorems.tex}
& A broad synthesis linking Stirling transforms, Bell polynomials, probability,
  inversion, implicit functions, and algebraic examples.
& Contributes the Boolean- and partition-lattice perspectives, the expanded
  moment--cumulant discussion, inverse-series examples, and the bridge from
  one-variable inversion to Lagrange--Good inversion.\\
\end{longtable}
\renewcommand{\arraystretch}{1}

The entries above describe provenance, not mutually exclusive ownership.  Many
results occur in three or more drafts.  Repetition of a theorem is removed, while
distinct proofs, examples, or normalizations are retained in the theorem's natural
chapter.

\section{Supplementary dossiers absorbed}
\begin{longtable}{@{}>{\raggedright\arraybackslash}p{0.25\textwidth}
                      >{\raggedright\arraybackslash}p{0.69\textwidth}@{}}
\toprule
\textbf{Folder or item} & \textbf{Integrated contribution}\\
\midrule
\endfirsthead
\toprule
\textbf{Folder or item} & \textbf{Integrated contribution}\\
\midrule
\endhead
\bottomrule
\endlastfoot

\nolinkurl{Wikipedia/Stirling_*}
& Definitions, triangular recurrences, factorial-basis transitions, generating
  functions, finite differences, asymptotic regimes, and variants.  These are
  unified in the Stirling parts and supplied with consistent signed/unsigned
  conventions.\\

\nolinkurl{Wikipedia/Bell_*}
& Set partitions, Touchard polynomials, Dobi\'nski's formula, Bell-polynomial
  composition, congruences, and asymptotics.  Complementary Bell numbers are added
  as the evaluation $T_n(-1)$.\\

\nolinkurl{Wikipedia/Eulerian_number.txt}
& Descent statistics, Worpitzky's identity, power sums, rational generating
  functions, and geometric interpretations, consolidated with the type-$B$ and
  second-order variants.\\

\nolinkurl{Wikipedia/Fa*_di_Bruno's_formula.txt}
& Partition, multiplicity, Bell-polynomial, and multivariate forms of the higher
  chain rule, expanded here to Banach-space derivatives, higher quotients, and
  nested composition.\\

\nolinkurl{Wikipedia/Lagrange_*}
& Formal, analytic, shifted, and combinatorial inversion formulas, including the
  plane-tree and associahedral interpretation.  The ordinary/exponential
  normalization ambiguity is explicitly resolved.\\

\nolinkurl{Wikipedia/Associahedron.*}
& Face--tree correspondence and face numbers.  These are connected directly to
  the coefficient polynomial of a compositional inverse and proved using the
  cycle lemma and Lagrange inversion.\\

MathWorld Bernoulli, Euler, Genocchi, and Euler--Maclaurin items
& Bernoulli and Euler Appell sequences, N\"orlund polynomials, Cauchy--Bernoulli
  polynomials of the second kind, Genocchi numbers, modified Bernoulli
  coefficients, Fourier expansions, Faulhaber summation, and exact
  Euler--Maclaurin remainders.\\

MathWorld Pochhammer, digamma, polygamma, and Barnes items
& Gamma-quotient factorials, parameter derivatives through complete Bell
  polynomials, harmonic and zeta evaluations, Barnes' functional equation,
  Alexeiewsky's integral, and the Euler--Maclaurin expansion of $\log G$.\\

MathWorld Darboux item
& The repeated-integration form of Taylor's theorem and its relationship to the
  periodic-Bernoulli kernel in Euler--Maclaurin summation.\\
\end{longtable}

\section{Externally researched extensions}
The multivariate determinant formula is developed in the form introduced by Good
and is connected with Gessel's combinatorial treatment
\cite{Good1960,Gessel1987}.  The final Sheffer--Riordan synthesis follows the
formal group viewpoint of Roman and the explicit comparison of the Sheffer and
Riordan groups by He, Hsu, and Shiue
\cite{Roman1984,HeHsuShiue2007}.  Analytic normalizations for Bernoulli, gamma,
polygamma, and Barnes functions were checked against the NIST Digital Library of
Mathematical Functions \cite{NISTDLMF}.  These sources are used to extend and
cross-check the supplied material; all identities needed in the body are proved
there rather than delegated to a citation.

\chapter{Normalization ledger}
\label{app:normalization-ledger}

\section{Ordinary versus exponential coefficients}
Let
\[
 X(z)=\sum_{j\ge1}x_j\frac{z^j}{j!},\qquad
 \widehat X(z)=\sum_{j\ge1}\widehat x_jz^j.
\]
The exponential and ordinary partial Bell polynomials are characterized by
\begin{align*}
 \frac{X(z)^k}{k!}
 &=\sum_{n\ge k}\BellP nk(x_1,x_2,\ldots)\frac{z^n}{n!},\\
 \widehat X(z)^k
 &=\sum_{n\ge k}\OBellP nk(\widehat x_1,\widehat x_2,\ldots)z^n.
\end{align*}
Putting $x_j=j!\widehat x_j$ gives the conversion
\begin{equation}\label{eq:appendix-bell-normalization}
 \OBellP nk(\widehat x_1,\widehat x_2,\ldots)
 =\frac{k!}{n!}\BellP nk(1!\widehat x_1,2!\widehat x_2,\ldots).
\end{equation}
This single identity accounts for nearly every factorial discrepancy among the
five drafts.  It also explains why derivative formulas use exponential Bell
polynomials, while ordinary Taylor substitution and the Laplace coefficients use
ordinary Bell polynomials.

\section{First-kind Stirling signs}
The two first-kind conventions are
\[
 \falling{x}{n}=\sum_{k=0}^ns(n,k)x^k,
 \qquad
 \rising{x}{n}=\sum_{k=0}^n\stirone nkx^k,
 \qquad
 s(n,k)=(-1)^{n-k}\stirone nk.
\]
Consequently the inverse matrices are
\[
 (s(n,k))_{n,k\ge0}^{-1}=(\stirtwo nk)_{n,k\ge0},
\]
whereas the unsigned array requires the sign diagonal:
\[
 (\stirone nk)^{-1}_{n,k}=(-1)^{n-k}\stirtwo nk.
\]
Any formula moved between rising and falling factorials must carry this sign.

\section{Bernoulli, Euler, and Eulerian objects}
The Bernoulli convention is fixed by
\[
 \frac{t}{e^t-1}=\sum_{n\ge0}\beta_n\frac{t^n}{n!},
 \qquad \beta_1=-\frac12.
\]
Euler polynomials are the Appell sequence generated by $2e^{xt}/(e^t+1)$.
Eulerian polynomials instead enumerate descents; they are not an Appell sequence.
Type-$B$ Eulerian numbers refer to the Coxeter group of signed permutations and
are unrelated to Bernoulli notation.  The disambiguation is semantic as well as
typographical: these families diagonalize different operators.

\section{Compositional inverses and branches}
For a formal series $f(z)=a_1z+O(z^2)$ with invertible $a_1$, the notation
$f^{\langle-1\rangle}$ always means the unique formal compositional inverse with
zero constant term.  In analytic applications it means the germ selected by the
stated base point.  Global multivalued functions such as Lambert $W$, logarithms,
and algebraic roots require a branch designation; a coefficient formula about one
germ is not silently asserted on every branch.

\section{Residues and Jacobians}
In one variable,
\[
 \res_t H(t)\,\dd t=\res_x H(f(x))f'(x)\,\dd x.
\]
In several variables, $f'$ becomes the full Jacobian determinant.  The determinant
in Good's theorem is therefore not an optional correction term: it is exactly the
factor needed to make multivariate residue invariant under substitution.

\chapter{Reproducible coefficient algorithms}
\label{app:algorithms}

\section{Triangular arrays}
For any target order $N$, the Stirling, Lah, and Eulerian triangles are generated
by their proved two-term recurrences.  A row-at-a-time implementation requires
$O(N)$ memory and $O(N^2)$ ring operations.  Exact integer arithmetic should be
used whenever the inputs are integral; floating-point evaluation is unnecessary
and obscures divisibility checks.

A generic lower-triangular recurrence has the form
\[
 a_{n,k}=u_{n,k}a_{n-1,k-1}+v_{n,k}a_{n-1,k}.
\]
Initialize $a_{0,0}=1$, declare out-of-range entries zero, and fill rows in
increasing $n$.  The boundary convention eliminates special branches in the inner
loop and is part of the mathematical specification, not merely a programming
convenience.

\section{Partial Bell polynomials}
The recurrence
\[
 \BellP nk(x_1,x_2,\ldots)
 =\sum_{j=1}^{n-k+1}\binom{n-1}{j-1}x_j
   \BellP{n-j}{k-1}(x_1,x_2,\ldots)
\]
computes the full table through order $N$.  It follows by marking the block
containing label $1$, so every generated polynomial has a direct combinatorial
certificate.  Homogeneity provides two independent checks:
\begin{align*}
 \BellP nk(cx_1,cx_2,\ldots)&=c^k\BellP nk(x_1,x_2,\ldots),\\
 \BellP nk(c x_1,c^2x_2,\ldots)&=c^n\BellP nk(x_1,x_2,\ldots).
\end{align*}

\section{Formal composition and reversion}
To compose truncated ordinary series modulo $z^{N+1}$, evaluate powers of the
inner series only up to the degree needed by the current outer coefficient.  Bell
polynomials expose all coefficient dependencies explicitly.  For reversion,
Lagrange--B\"urmann gives
\[
 [z^n]f^{\langle-1\rangle}(z)
 =\frac1n[u^{n-1}]\left(\frac{u}{f(u)}\right)^n.
\]
This is usually the clearest symbolic algorithm because every target coefficient
is independent.  Newton iteration
\[
 g_{\mathrm{new}}=g-\frac{f(g)-z}{f'(g)}
\]
doubles the number of correct coefficients at each step and is preferable at
very high orders when fast multiplication is available.  A reliable
implementation computes by one method and verifies
$f(g(z))\equiv z\pmod{z^{N+1}}$ by the other.

\section{Multivariate Good coefficients}
For fixed multi-index $\boldsymbol n$, only monomials bounded componentwise by
$\boldsymbol n$ are needed.  Truncate every $\phi_i$ and $F$ to that finite box,
form the logarithmic Jacobian matrix
\[
 J_{ij}=\delta_{ij}-\frac{x_j}{\phi_i}\partial_j\phi_i,
\]
compute its determinant in the truncated multivariate ring, multiply by
$F\prod_i\phi_i^{n_i}$, and extract $\boldsymbol x^{\boldsymbol n}$.  Substitution
of the computed coefficients back into
$w_i=t_i\phi_i(\boldsymbol w)$ is an exact finite certificate.

\section{Asymptotic coefficients}
Euler--Maclaurin, saddle-point, and Laplace expansions should be stored as formal
coefficient generators rather than as fixed decimal tables.  For an expansion
through order $r$, retain one additional derivative or coefficient and verify that
it contributes only to the stated remainder order.  The exact periodic-Bernoulli
remainder in Euler--Maclaurin and the Cauchy bounds for analytic coefficient
extraction provide quantitative, independently checkable error estimates.

\chapter{Independent consistency checks}
\label{app:consistency-checks}

\section{Low-order composition check}
Writing
\[
 f(u)=f_0+f_1u+\frac{f_2}{2!}u^2+\frac{f_3}{3!}u^3+\frac{f_4}{4!}u^4+\cdots
\]
and $g(x+h)=g_0+g_1h+g_2h^2/2!+\cdots$, direct multiplication through $h^4$
gives
\begin{align*}
 D^4(f\circ g)
 ={}&f_1g_4+f_2(4g_1g_3+3g_2^2)
      +6f_3g_1^2g_2+f_4g_1^4.
\end{align*}
The coefficients $1,4,3,6,1$ count set partitions of types
$4$, $3+1$, $2+2$, $2+1+1$, and $1+1+1+1$.  Their sum is $\Bell4=15$, checking
simultaneously the partition, multiplicity, and Bell-polynomial forms of Fa\`a di
Bruno.

\section{Low-order inverse check}
For
\[
 f(x)=x+a_1x^2+a_2x^3+a_3x^4+a_4x^5+O(x^6),
\]
the associahedral formula gives
\begin{align*}
 f^{\langle-1\rangle}(y)
 ={}&y-a_1y^2+(2a_1^2-a_2)y^3
 +(5a_1a_2-5a_1^3-a_3)y^4\\
 &+(14a_1^4-21a_1^2a_2+3a_2^2+6a_1a_3-a_4)y^5+O(y^6).
\end{align*}
Direct substitution into $f(g(y))$ cancels each coefficient through $y^5$.
The absolute coefficients $1,6,3,21,14$ in the fifth-order polynomial are the
corresponding face-profile counts of $K_4$.

\section{Matrix check for Stirling and Riordan inversion}
Truncating at any order $N$, multiply the signed first-kind matrix
$(s(n,k))_{0\le n,k\le N}$ by the second-kind matrix
$(\stirtwo nk)_{0\le n,k\le N}$.  Factorial-basis substitution proves that the
product is the identity over $\Z$; direct multiplication is an independent finite
certificate.  The same test applies to a general exponential Riordan pair:
truncate $(g,f)$ and $(1/(g\circ\bar f),\bar f)$ and verify their product is
$(1,z)$ modulo $z^{N+1}$.

\section{Analytic versus formal checks}
A formal identity can be checked in a finite quotient ring without evaluating a
single transcendental function.  An analytic identity needs a second layer:
locate singularities and critical points, select branches, and bound the tail.
Kepler's equation illustrates the distinction.  The Lagrange series is formally
forced by $E=M+e\sin E$; the inequality $1-e\cos E>0$ for $0\le e<1$ proves the
real analytic inverse, and the Fourier--Bessel series supplies a globally periodic
resummation.

\chapter{Formula-dependency map}
\label{app:dependency-map}
The principal logical routes can be summarized without suppressing the
parallel proof mechanisms:
\begin{center}
\renewcommand{\arraystretch}{1.28}
\begin{tabularx}{\textwidth}{@{}>{\raggedright\arraybackslash\bfseries\color{deepolive}}p{0.23\textwidth}X@{}}
\toprule
Route & Dependency chain \\
\midrule
Foundations
  & formal power series and residues $\longrightarrow$ factorial bases and finite
    differences $\longrightarrow$ Stirling and Lah transforms \\
Labelled structures
  & Boolean and partition M\"obius inversion $\longrightarrow$ labelled sets and
    partitions $\longrightarrow$ Bell numbers and Bell polynomials \\
Descents and summation
  & Stirling transforms and finite differences $\longrightarrow$ Eulerian arrays
    and power sums $\longrightarrow$ Bernoulli--Euler summation calculus \\
Composition and inversion
  & Bell polynomials together with formal residues $\longrightarrow$ Fa\`a di
    Bruno composition calculus $\longrightarrow$ Lagrange and Good inversion \\
Operator groups
  & factorial bases together with labelled-component calculus $\longrightarrow$
    Sheffer sequences and Riordan groups $\longrightarrow$ substitution and
    inverse-array formulas \\
Special functions
  & Bernoulli--Euler calculus together with repeated composition $\longrightarrow$
    gamma and polygamma identities $\longrightarrow$ Barnes' $G$-function and
    asymptotic coefficient expansions \\
\bottomrule
\end{tabularx}
\end{center}
The map is not merely pedagogical.  It predicts proof techniques.  A problem
about polynomial bases suggests triangular matrices or a delta operator; a problem
about labelled components suggests exponential generating functions and Bell
polynomials; an implicit substitution suggests residues and Lagrange inversion;
a lattice sum suggests Bernoulli kernels and Euler--Maclaurin.  Most formulas in
the source corpus become short once moved to the coordinate system in which their
operator is diagonal.

\chapter{Completeness and correction record}
\label{app:correction-record}

\section{Material intentionally deduplicated}
Repeated introductions to ordinary and exponential generating functions, repeated
first tables of Stirling and Bell numbers, and repeated statements of the same
recurrence are represented once.  Omitting duplicate prose does not omit a
mathematical assertion: variants are recorded when their hypotheses, indexing, or
proof mechanism differ.

\section{Systematic corrections}
The following classes of defect were repaired during consolidation.
\begin{enumerate}
\item Signed and unsigned first-kind Stirling numbers were separated whenever a
      rising-factorial formula had been combined with a falling-factorial sign.
\item The type-$B$ descent index was fixed at positions $0,1,\ldots,n-1$ with
      $\pi_0=0$, eliminating a one-step shift.
\item Ordinary inverse-series coefficients in the associahedral formula were
      separated from exponential Bell-polynomial normalizations.
\item Formal inversion hypotheses ($f(0)=0$, invertible linear coefficient) were
      separated from analytic inverse-function hypotheses and branch choices.
\item The distributional reversion proof was supplied with its Jacobian and
      orientation condition.
\item Multivariate inversion was supplied with the full determinant that one
      variable hides.
\item Asymptotic displays with unnamed correction polynomials were replaced by
      coefficient recurrences or exact remainder formulas.
\item Special-function identities were assigned their actual pole exclusions,
      half-planes, sectors, or local convergence domains.
\end{enumerate}

\section{Coverage principle}
Every distinct theorem, identity family, example class, normalization issue, and
computational certificate found in the five drafts has been assigned a location in
the unified dependency structure.  The supplementary Wikipedia and MathWorld
materials have been used to fill missing links and extend the theory, especially in
Bernoulli--Euler summation, complementary Bell numbers, associahedra, factorial
parameter derivatives, and Barnes' $G$-function.  The externally researched Good,
Gessel, and Sheffer--Riordan material closes the progression from triangular
one-variable transforms to multivariate and group-theoretic inversion.

\backmatter

\chapter*{Selected bibliography and provenance}
\addcontentsline{toc}{chapter}{Selected bibliography and provenance}

\begin{thebibliography}{99}

\bibitem{Bell1934}
E.~T. Bell,
``Exponential polynomials,''
\emph{Annals of Mathematics}, second series, vol.~35 (1934), pp.~258--277.

\bibitem{Broder1984}
A.~Z. Broder,
``The $r$-Stirling numbers,''
\emph{Discrete Mathematics}, vol.~49 (1984), pp.~241--259.

\bibitem{Charalambides2002}
C.~A. Charalambides,
\emph{Enumerative Combinatorics},
Chapman \& Hall/CRC, 2002.

\bibitem{Claesson2001}
A. Claesson,
``Generalized pattern avoidance,''
\emph{European Journal of Combinatorics}, vol.~22 (2001), pp.~961--971.

\bibitem{Comtet1974}
L. Comtet,
\emph{Advanced Combinatorics: The Art of Finite and Infinite Expansions},
translated by J.~W. Nienhuys, D. Reidel, 1974.

\bibitem{FlajoletSedgewick2009}
P. Flajolet and R. Sedgewick,
\emph{Analytic Combinatorics},
Cambridge University Press, 2009.

\bibitem{GKP1994}
R.~L. Graham, D.~E. Knuth, and O. Patashnik,
\emph{Concrete Mathematics}, second edition,
Addison--Wesley, 1994.

\bibitem{Harper1967}
L.~H. Harper,
``Stirling behavior is asymptotically normal,''
\emph{Annals of Mathematical Statistics}, vol.~38 (1967), pp.~410--414.

\bibitem{Henrici1974}
P. Henrici,
\emph{Applied and Computational Complex Analysis}, vol.~1,
Wiley, 1974.

\bibitem{HsuShiue1998}
L.~C. Hsu and P.~J.-S. Shiue,
``A unified approach to generalized Stirling numbers,''
\emph{Advances in Applied Mathematics}, vol.~20 (1998), pp.~366--384.

\bibitem{Johnson2002}
W.~P. Johnson,
``The curious history of Faà di Bruno's formula,''
\emph{American Mathematical Monthly}, vol.~109 (2002), pp.~217--234.

\bibitem{MoserWyman1955}
L. Moser and M. Wyman,
``An asymptotic formula for the Bell numbers,''
\emph{Transactions of the Royal Society of Canada}, section III, vol.~49
(1955), pp.~49--54.

\bibitem{NISTDLMF}
F.~W.~J. Olver, A.~B. Olde Daalhuis, D.~W. Lozier, B.~I. Schneider,
R.~F. Boisvert, C.~W. Clark, B.~R. Miller, B.~V. Saunders,
H.~S. Cohl, and M.~A. McClain, editors,
\emph{NIST Digital Library of Mathematical Functions},
National Institute of Standards and Technology.

\bibitem{Riordan1958}
J. Riordan,
\emph{An Introduction to Combinatorial Analysis},
Wiley, 1958.

\bibitem{Roman1984}
S. Roman,
\emph{The Umbral Calculus},
Academic Press, 1984.

\bibitem{Spivey2008}
M.~Z. Spivey,
``A generalized recurrence for Bell numbers,''
\emph{Journal of Integer Sequences}, vol.~11 (2008), Article 08.2.5.

\bibitem{Stanley2012}
R.~P. Stanley,
\emph{Enumerative Combinatorics}, vol.~1, second edition,
Cambridge University Press, 2012.

\bibitem{Temme1993}
N.~M. Temme,
``Asymptotic estimates of Stirling numbers,''
\emph{Studies in Applied Mathematics}, vol.~89 (1993), pp.~233--243.

\bibitem{Wilf2005}
H.~S. Wilf,
\emph{generatingfunctionology}, third edition,
A K Peters, 2005.

\bibitem{Apostol1976}
T.~M. Apostol,
\emph{Introduction to Analytic Number Theory},
Springer, 1976.

\bibitem{DvoretzkyMotzkin1947}
A. Dvoretzky and T. Motzkin,
``A problem of arrangements,''
\emph{Duke Mathematical Journal}, vol.~14 (1947), pp.~305--313.

\bibitem{Gessel1987}
I.~M. Gessel,
``A combinatorial proof of the multivariable Lagrange inversion formula,''
\emph{Journal of Combinatorial Theory, Series A}, vol.~45 (1987),
pp.~178--195.

\bibitem{Gessel2016}
I.~M. Gessel,
``Lagrange inversion,''
\emph{Journal of Combinatorial Theory, Series A}, vol.~144 (2016),
pp.~212--249.

\bibitem{Good1960}
I.~J. Good,
``Generalizations to several variables of Lagrange's expansion, with applications
to stochastic processes,''
\emph{Proceedings of the Cambridge Philosophical Society}, vol.~56 (1960),
pp.~367--380.

\bibitem{HeHsuShiue2007}
T.-X. He, L.~C. Hsu, and P.~J.-S. Shiue,
``The Sheffer group and the Riordan group,''
\emph{Discrete Applied Mathematics}, vol.~155 (2007), pp.~1895--1909.

\bibitem{Norlund1924}
N.~E. N\"orlund,
\emph{Vorlesungen \"uber Differenzenrechnung},
Springer, 1924.

\bibitem{ShapiroGetuWoanWoodson1991}
L.~W. Shapiro, S. Getu, W.-J. Woan, and L.~C. Woodson,
``The Riordan group,''
\emph{Discrete Applied Mathematics}, vol.~34 (1991), pp.~229--239.

\bibitem{WhittakerWatson1927}
E.~T. Whittaker and G.~N. Watson,
\emph{A Course of Modern Analysis}, fourth edition,
Cambridge University Press, 1927.

\bibitem{SupplementaryDossiers}
The supplied \emph{Wikipedia} and \emph{MathWorld} snapshot folders, containing
articles and notebooks on Stirling, Bell, Catalan, Eulerian, Fa\`a di Bruno,
Lagrange, associahedral, Bernoulli, Euler, Genocchi, Pochhammer, polygamma,
Euler--Maclaurin, Darboux, and Barnes-$G$ topics.

\bibitem{SourceArchive}
\emph{Formulae-2.zip}, the five supplied draft manuscripts:
\nolinkurl{combinatorial_coefficient_calculus.tex},
\nolinkurl{Formulae_Compendium.tex}, two independently organized drafts named
\nolinkurl{Formulae_Monograph.tex}, and
\nolinkurl{Combinatorial_Formulae_and_Inversion_Theorems.tex}, together with the
supplementary Wikipedia and MathWorld directories.

\end{thebibliography}

\printindex
\end{document}
